\documentclass[oneside, a4paper]{book}
\usepackage[left=3.5cm,right=2cm,top=3.5cm,bottom=3cm]{geometry}
\usepackage[nottoc]{tocbibind}
\usepackage{indentfirst}
\usepackage[unicode]{hyperref}
\usepackage[protrusion=false,expansion=false]{microtype}
\usepackage{cite}
\usepackage{pdfpages}
\usepackage{enumitem} 
\usepackage{float}
\usepackage{graphicx}
\usepackage{wrapfig}
\usepackage{tikz, pgf}
\usetikzlibrary{shapes}
\usetikzlibrary{patterns}
\usetikzlibrary{arrows.meta, calc}
\usepackage{listings}
\usepackage{multirow}
\usepackage{xcolor}
\usepackage{enumitem}
\definecolor{codegreen}{rgb}{0,0.6,0}
\definecolor{codegray}{rgb}{0.5,0.5,0.5}
\definecolor{codepurple}{rgb}{0.58,0,0.82}
\definecolor{backcolour}{rgb}{0.95,0.95,0.92}

\lstdefinestyle{mystyle}{
    backgroundcolor=\color{backcolour},   
    commentstyle=\color{codegreen},
    keywordstyle=\color{magenta},
    numberstyle=\tiny\color{codegray},
    stringstyle=\color{codepurple},
    basicstyle=\ttfamily\normalsize,
    breakatwhitespace=false,         
    breaklines=true,                 
    captionpos=b,                    
    keepspaces=true,                 
    numbers=left,                    
    numbersep=5pt,                  
    showspaces=false,                
    showstringspaces=false,
    showtabs=false,                  
    tabsize=2
}
\lstset{style=mystyle}

\hbadness=10000
\vbadness=10000
\hfuzz=45pt
\vfuzz=15pt

\usepackage{amsmath,amssymb,latexsym,amscd,amsxtra,graphicx,graphpap,amsthm}
\usepackage{mathtools}
\usepackage{pb-diagram}
\usepackage{caption}
\usepackage{aligned-overset}
\usepackage{picinpar,floatflt}
\usepackage{longtable}
\usepackage{fancybox}
\usepackage[utf8]{vietnam}
\usepackage{anyfontsize}
\usepackage{tabularx}
\usepackage{comment}
\usepackage{algpseudocode}
\usepackage{adjustbox}
\usepackage{algorithm}
\usepackage{graphicx}




\usepackage{tikz}
\usetikzlibrary{positioning, arrows.meta}
\usepackage{float} % Bắt buộc để dùng tham số [H] ép vị trí hình ảnh
\usepackage{tikz}
\usetikzlibrary{positioning, arrows.meta, calc}
\usepackage{amsmath}
\usetikzlibrary{positioning, arrows.meta, calc, decorations.pathreplacing}
\usepackage{float}       % Để sử dụng tham số [H] cố định vị trí hình ảnh
\usepackage{caption}     % Hỗ trợ tùy chỉnh chú thích
\usepackage{subcaption}  % Bắt buộc phải có để dùng môi trường \begin{subfigure}
\usepackage{tikz}        % Thư viện vẽ hình lõi
\usetikzlibrary{arrows.meta, positioning, decorations.pathreplacing} % Các gói mở rộng của TikZ
\usepackage{titletoc}









\makeatletter
\renewcommand{\ALG@name}{\large Thuật toán}
\makeatother
\usepackage{etoolbox}
\makeatletter
\patchcmd{\ALG@beginalgorithmic}{\tiny}{\normalsize}{}{}
\algrenewcommand\alglinenumber[1]{\large #1:}
\makeatother

%==================================
% ĐỊNH DẠNG MỤC LỤC
%==================================
\titlecontents{chapter}
  [0pt]
  {\bfseries}
  {\chaptername\ \thecontentslabel\enskip}
  {}
  {\hfill\contentspage}

%==================================
% ĐÁNH SỐ CHƯƠNG / SECTION
%==================================

\renewcommand{\thechapter}{\arabic{chapter}} 
\renewcommand{\thesection}{\thechapter.\arabic{section}} % Ghép số chương vào mục lớn (1.1, 1.2, 2.1...)
\renewcommand{\thesubsection}{\thesection.\arabic{subsection}} % Mục nhỏ: 1.1.1, 1.1.2...
\renewcommand{\thesubsubsection}{\thesubsection.\arabic{subsubsection}} % Mục nhỏ hơn: 1.1.1.1...

\newcommand{\eproof}{\hfill $\square$}
\newcommand{\chm}{{\bfseries Chứng minh.}}
\newenvironment{cm}{\chm}{\eproof}

%==================================
% CÁC MÔI TRƯỜNG ĐỊNH LÝ
%==================================
\theoremstyle{plain}
\newtheorem{bd}{Bổ đề}[chapter]
\providecommand*{\bdautorefname}{Bổ đề}
\newtheorem{md}{Mệnh đề}[chapter]
\providecommand*{\mdautorefname}{Mệnh đề}
\newtheorem{hq}{Hệ quả}[chapter]
\providecommand*{\hqautorefname}{Hệ quả}
\newtheorem{dl}{\bfseries Định lý}[chapter]
\providecommand*{\dlautorefname}{Định lý}
\newtheorem{tc}{\bfseries Tính chất}[chapter]
\providecommand*{\tcautorefname}{Tính chất}
\newtheorem{st}{\bf Bước}
\newtheorem{stc}{\hspace{0.5em} Bước}[st]
\newtheorem*{namedthm}{\namedthmname}
\newcounter{namedthm}
\makeatletter
\newenvironment{named}[1]
  {\def\namedthmname{#1}%
   \refstepcounter{namedthm}%
   \namedthm\def\@currentlabel{#1}}
  {\endnamedthm}
\makeatother

\theoremstyle{definition}
\newtheorem{dn}{Định nghĩa}[chapter]
\providecommand*{\dnautorefname}{Định nghĩa}
\newtheorem{bt}{Bài toán}[chapter]
\providecommand*{\btautorefname}{Bài toán}
\newtheorem{vd}{\bfseries Ví dụ}[chapter]
\providecommand*{\vdautorefname}{Ví dụ}
\newtheorem{assumption}{\bfseries Giả thiết}[chapter]
\providecommand*{\assumptionautorefname}{Giả thiết}

\theoremstyle{remark}
\newtheorem{kh}{ký hiệu}[chapter]
\newtheorem{nx}{Nhận xét}[chapter]
\newtheorem{ly}{Lưu ý}[chapter]
\renewcommand{\thekh}{\textit{\thechapter.\arabic{kh}}}
\renewcommand{\thenx}{\textit{\thechapter.\arabic{nx}}}
\renewcommand{\thely}{\textit{\thechapter.\arabic{ly}}}

%==================================
% BỐ CỤC TRANG
%==================================
\renewcommand{\baselinestretch}{1.5}
\setlength{\oddsidemargin}{0.3cm}
\setlength{\topmargin}{-1cm}
\setlength{\headsep}{0.5cm}
\textwidth=15.5cm
\textheight=24.5cm

%==================================
% ĐỔI TÊN CÁC THÀNH PHẦN SANG TIẾNG VIỆT
%==================================
\renewcommand{\bibname}{Tài liệu tham khảo}
\renewcommand{\listfigurename}{Danh sách hình vẽ}
\renewcommand{\listtablename}{Danh sách bảng}

%==================================
% CỠ CHỮ TÙY CHỈNH
%==================================
\renewcommand{\large}{\fontsize{14pt}{14pt}\selectfont}
\renewcommand{\Large}{\fontsize{16pt}{16pt}\selectfont}
\renewcommand{\LARGE}{\fontsize{15pt}{15pt}\selectfont}
\renewcommand{\th}{\fontsize{18pt}{18pt}\selectfont}

%==================================
% HEADER - FOOTER (fancyhdr)
%==================================
\usepackage{fancyhdr}
\usepackage{emptypage}

% Cấu hình header & footer cho trang thường (áp dụng từ chương 1)
\pagestyle{fancy}
\fancyhf{}
\fancyhead[L]{\small\textit{Đồ án tốt nghiệp}} % Trái: Đồ án tốt nghiệp
\fancyhead[R]{\small\nouppercase{\leftmark}} % Phải: Chương hiện tại và tên chương
\fancyfoot[C]{\thepage} % Giữa dưới cùng: Số trang
\renewcommand{\headrulewidth}{0.4pt}
\renewcommand{\footrulewidth}{0pt}

% Trang mở đầu chương: chỉ có số trang, không có header
\fancypagestyle{plain}{
  \fancyhf{}
  \fancyfoot[C]{\thepage}
  \renewcommand{\headrulewidth}{0pt}
}

% Định dạng lại leftmark để hiển thị "Chương X. Tên chương"
\renewcommand{\chaptermark}[1]{%
  \markboth{Chương \thechapter.\ #1}{}%
}

%==================================
% CÁC LỆNH TOÁN / DELIMITER
%==================================
\newenvironment{mproof}{\paragraph{Chứng minh:}}{\hfill$\square$}
\newenvironment{myproof}[2] {\paragraph{Proof of {#1} {#2} :}}{\hfill$\square$}
\DeclarePairedDelimiter{\pro}{\langle}{\rangle}
\DeclarePairedDelimiter{\norm}{\lVert}{\rVert}
\DeclarePairedDelimiter{\abs}{\lvert}{\rvert}
\usepackage{xpatch}
\makeatletter
\AtBeginDocument{\xpatchcmd{\@thm}{\thm@headpunct{.}}{\thm@headpunct{}}{}{}}
\makeatother

%=====================================================
%              BẮT ĐẦU TÀI LIỆU
%=====================================================
\begin{document}

%-----------------------------------------------------
% TRANG BÌA (không đánh số, không header)
%-----------------------------------------------------
\pagestyle{empty}
\pagenumbering{gobble}
\newgeometry{left=2.8cm,right=1.8cm,top=2.5cm,bottom=2cm}
\setcounter{secnumdepth}{3}

\begin{center}
    \Large
    \vspace{-1cm}
    {\bfseries ĐẠI HỌC BÁCH KHOA HÀ NỘI}\\
    {\bfseries KHOA TOÁN - TIN}

    \vspace{1.5cm}
    \includegraphics[width=4cm]{hust.jpg}
\end{center}

\vspace{1cm}


\begin{center}
    \fontsize{25pt}{16pt}\selectfont    
    {\bfseries ĐỒ ÁN TỐT NGHIỆP}

    \vspace{1cm}
    
    \fontsize{20pt}{16pt}\selectfont
    {\bfseries XÂY DỰNG HỆ THỐNG XỬ LÝ VÀ PHÂN TÍCH DỮ LIỆU GIAO DỊCH NGÂN HÀNG} %Điền tên chủ đề

    \vspace{1cm}
    
    \fontsize{16pt}{16pt}\selectfont
    {\bfseries NGUYỄN HỒNG NHUNG}
    
    \fontsize{14pt}{16pt}\selectfont {Nhung.nh227253@sis.hust.edu.vn}
    
    \vspace{0.5cm}
    
    \fontsize{13pt}{16pt}\selectfont
    {\textbf{Chuyên ngành:} \fontsize{15pt}{16pt}\selectfont Hệ thống thông tin quản lý}
\end{center}

\vspace{1.5cm}

\fontsize{14pt}{16pt}\selectfont
{
\centering
    \begin{tabular}{ l l }
        \textbf{Giảng viên hướng dẫn:}& \fontsize{16pt}{16pt}\selectfont ThS. Nguyễn Danh Tú \newcommand\tab[1][0.3cm]{\hspace*{#1}}\tab $\underset{\text{Chữ ký của GVHD}}{\rule{3.5cm}{0.5pt}}$ \\
    \end{tabular}
}

\vspace{2cm}

\begin{center}
    \LARGE HÀ NỘI -- 2026
\end{center}
    
\restoregeometry

%-----------------------------------------------------
% FRONT MATTER (Đánh số La Mã, Không Header)
%-----------------------------------------------------
\newpage
\pagenumbering{roman} % Đánh số La Mã cho các phần trước Chương 1
\setcounter{page}{1}
\pagestyle{plain} % Dùng style plain: Không header, có số trang ở footer

\thisfancypage{
    \setlength{\fboxsep}{0.5cm}
    \fbox}{}

\fontsize{13pt}{16pt}\selectfont
\bigbreak
\begin{center}
    {\bfseries NHẬN XÉT CỦA GIẢNG VIÊN HƯỚNG DẪN}
\end{center}
\bigbreak

\fontsize{12pt}{14pt}\selectfont
\begin{enumerate}
    \item [{\bfseries 1.}]{\bfseries Mục tiêu và nội dung của đồ án}
    \begin{enumerate}
        \item Mục tiêu: .
        \item Nội dung: . 
    \end{enumerate}
    \item [{\bfseries 2.}] {\bfseries Kết quả đạt được} 
    \begin{enumerate}
        \item 1. 
        \item 2.
        \item 3.
    \end{enumerate}
    \item [{\bfseries 3.}]{\bfseries Ý thức làm việc của sinh viên:}
\end{enumerate}
\hspace{0.5\textwidth}
\begin{minipage}{0.5\textwidth}
    \noindent\begin{center}
        \textit{Hà Nội, ngày ....    tháng ... năm 2026} \\
        Giảng viên hướng dẫn\\ \vspace{2cm}
        \textbf{ThS. Nguyễn Danh Tú}
    \end{center}   
\end{minipage}

\newpage
\fontsize{22pt}{16pt}\selectfont
{\bfseries Lời cảm ơn}
\fontsize{13pt}{16pt}\selectfont
\bigskip

Lời đầu tiên, em xin gửi lời cảm ơn chân thành và sâu sắc nhất đến quý thầy cô Khoa \textbf{Toán - Tin}, Trường \textbf{Đại học Bách khoa Hà Nội}, những người đã tận tình truyền đạt kiến thức, định hướng tư duy và tạo điều kiện tốt nhất để em hoàn thành chương trình học trong suốt những năm qua.

Em xin bày tỏ lòng biết ơn sâu sắc đến \textbf{Thầy giáo Nguyễn Danh Tú}– người đã trực tiếp hướng dẫn, tận tình chỉ bảo và luôn đồng hành cùng em trong suốt quá trình thực hiện Đồ án 1, Đồ án 2 và đặc biệt là Đồ án Tốt nghiệp. Những góp ý chuyên môn sâu sắc, sự tận tâm trong công việc cùng những lời động viên, khích lệ đúng lúc của thầy đã trở thành nguồn động lực to lớn giúp em vượt qua những khó khăn, thử thách trong quá trình nghiên cứu và hoàn thiện đề tài. Sự tận tụy, trách nhiệm và tâm huyết của thầy là tấm gương để em học hỏi và noi theo.


Em cũng xin cảm ơn các bạn trong nhóm học tập, những người đã cùng thảo luận, chia sẻ kinh nghiệm và hỗ trợ nhau trong suốt thời gian học tập 4 năm tại trường.

Cuối cùng, con xin gửi lời tri ân sâu sắc đến gia đình, đặc biệt là bố mẹ, những người luôn ở bên, động viên và là chỗ dựa vững chắc để con có thể toàn tâm hoàn thành đồ án này.

Mặc dù đã cố gắng hết sức, đồ án chắc chắn còn nhiều thiếu sót. Em rất mong nhận được sự góp ý từ quý thầy cô và các bạn để đề tài được hoàn thiện hơn.

\vspace{1cm}

\begin{minipage}{0.5\textwidth}
\end{minipage}
\hspace{0.5\textwidth}
\begin{minipage}{0.5\textwidth}
    \begin{center}
        \textit{Hà Nội, tháng 6 năm 2026} \\
        Tác giả Đồ án \\
        \vspace{2cm}
        \textbf{Nguyễn Hồng Nhung}
    \end{center}
\end{minipage}

\newpage
\fontsize{22pt}{16pt}\selectfont
{\bfseries Tóm tắt nội dung Đồ án}
\fontsize{13pt}{16pt}\selectfont
\bigskip
\bigskip

Đồ án tập trung nghiên cứu và xây dựng hệ thống nền tảng dữ liệu (Data Platform) phục vụ cho bài toán trong lĩnh vực ngân hàng. Thông qua việc thu thập, xử lý và trực quan hóa, đồ án giải quyết trọn vẹn bài toán tích hợp cả dữ liệu có cấu trúc (thông tin khách hàng, thẻ, giao dịch) thông qua luồng ETL và dữ liệu phi cấu trúc (hình ảnh Căn cước công dân, Sổ tiết kiệm) thông qua công nghệ nhận dạng quang học. Nội dung đồ án được triển khai chi tiết qua bốn phần chính như sau:

\begin{enumerate}
    \item \textbf{Tổng quan lý thuyết:} 
    Phần này trình bày các cơ sở lý thuyết nền tảng làm kim chỉ nam cho toàn bộ quá trình nghiên cứu và xây dựng hệ thống. Nội dung bao quát các khái niệm cốt lõi về Kho dữ liệu (Data Warehouse) truyền thống và sự tiến hóa sang kiến trúc đa tầng Medallion (Bronze -- Silver -- Gold) hiện đại. Đồ án cũng đi sâu vào lý thuyết của mô hình Data Vault 2.0 phục vụ quản lý lịch sử thay đổi (SCD) một cách linh hoạt và mô hình đa chiều (Star Schema) phục vụ phân tích dữ liệu chuyên sâu. Đồng thời, nền tảng lý thuyết về quy trình tích hợp dữ liệu (ETL), kỹ thuật trích xuất đặc trưng hình ảnh bằng công nghệ nhận dạng ký tự quang học (OCR) và vai trò của các công cụ điều phối luồng công việc (Workflow Orchestration) cũng được làm rõ.

    \item \textbf{Khảo sát và phân tích yêu cầu:} 
    Quá trình khảo sát đánh giá chi tiết thực trạng quản lý dữ liệu tại các hệ thống ngân hàng hiện nay, nơi dữ liệu thường bị phân mảnh ở nhiều hệ thống tác nghiệp (OLTP) và các quy trình xử lý hồ sơ chứng từ vẫn phụ thuộc nhiều vào thao tác nhập liệu thủ công. Từ việc xác định những hạn chế về tính đồng bộ, khả năng truy xuất lịch sử và mức độ tự động hóa, đồ án đặt ra các yêu cầu chức năng và phi chức năng khắt khe. Dựa trên đó, giải pháp được đề xuất là xây dựng hai luồng dữ liệu (Data Pipeline) độc lập nhưng thống nhất: một luồng tinh gọn xử lý dữ liệu phi cấu trúc bằng OCR và một luồng ETL toàn diện xử lý dữ liệu giao dịch có cấu trúc.

    \item \textbf{Phân tích và thiết kế hệ thống:} 
    Ở giai đoạn này, kiến trúc hệ thống được thiết kế chi tiết với cơ sở dữ liệu tập trung trên nền tảng SQL Server. Đối với dữ liệu có cấu trúc, đồ án thiết kế chiến lược nạp dữ liệu gia tăng (Incremental Loading) sử dụng chiến lược MNS để bóc tách các thay đổi hàng ngày. Dữ liệu sau đó được lưu trữ tại tầng Silver theo mô hình Data Vault (gồm Hub, Link, Satellite) nhằm đảm bảo tính toàn vẹn, trước khi chuyển hóa sang Star Schema tại tầng Gold để tối ưu hóa truy vấn. Đối với dữ liệu phi cấu trúc, đồ án thiết kế các lược đồ bảng tại tầng Bronze để lưu trữ trực tiếp văn bản thô được trích xuất từ ảnh chứng từ, đi kèm với hệ thống bảng Audit để giám sát và đánh giá độ tin cậy của thuật toán.

    \item \textbf{Xây dựng và triển khai hệ thống:} 
    Hệ thống được lập trình và vận hành thực tế thông qua việc cô lập hai môi trường Python ảo nhằm giải quyết triệt để sự xung đột tài nguyên giữa engine OCR và các công cụ dữ liệu. Đồ án ứng dụng Apache Airflow để điều phối tự động kịch bản chạy định kỳ và công cụ dbt để thực hiện các phép biến đổi SQL, đồng thời tích hợp hàng loạt các bài kiểm tra chất lượng dữ liệu (Data Quality checks) như kiểm soát giá trị rỗng, trùng lặp khóa và tính cập nhật. Cuối cùng, toàn bộ nguồn dữ liệu đã được làm sạch tại tầng Gold được kết nối với hệ thống Power BI để xây dựng các báo cáo quản trị (Dashboard) trực quan, hỗ trợ ra quyết định kinh doanh.
\end{enumerate}

Phần kết luận tổng hợp lại những giá trị cốt lõi mà nền tảng dữ liệu mang lại trong việc nâng cao năng lực khai thác thông tin và số hóa quy trình tác nghiệp ngân hàng. Định hướng phát triển trong tương lai mở ra các tiềm năng nâng cấp kỹ thuật như xử lý song song nhiều tài liệu phức tạp, tinh chỉnh các bước tiền xử lý ảnh và thiết kế giao diện ứng dụng để người dùng dễ dàng theo dõi, can thiệp vào kết quả trích xuất.



% \vspace{1cm}
% \begin{minipage}{0.5\textwidth}
% \end{minipage}
% \hspace{0.5\textwidth}
% \begin{minipage}{0.5\textwidth}
%     \noindent\begin{center}
%         \textit{Hà Nội, tháng 1 năm 2026} \\
%         Tác giả đồ án\\ \vspace{2cm}
%         \textbf{Nguyễn Hồng Nhung}
%     \end{center}   
% \end{minipage}

\newpage
\tableofcontents % Mục lục

\newpage
\listoffigures % Danh sách hình vẽ

\newpage
\pdfbookmark[0]{Danh sach bang}{lot}
\hypersetup{bookmarksdepth=-2}
\listoftables % Danh sách bảng
\hypersetup{bookmarksdepth=3}

\chapter*{Bảng ký hiệu và chữ viết tắt}
\addcontentsline{toc}{chapter}{\bfseries Bảng ký hiệu và chữ viết tắt}

\begin{longtable}{@{\hspace{-0.1cm}} l @{\hspace{1.2cm}} p{11.5cm}}
    AI    & Artificial Intelligence - trí tuệ nhân tạo\\
    API   & Application Programming Interface - giao diện lập trình ứng dụng\\
    Ad-hoc & Truy vấn/phân tích tùy biến, phát sinh tức thời theo nhu cầu cụ thể\\

    BI    & Business Intelligence \\
    CCCD  & Căn cước công dân\\
    CDC   & Change Data Capture - thu thập dữ liệu thay đổi\\
    CLAHE & Contrast Limited Adaptive Histogram Equalization - cân bằng histogram thích ứng (tiền xử lý ảnh)\\
    CRM   & Customer Relationship Management - quản lý quan hệ khách hàng\\
    DAG   & Directed Acyclic Graph - đồ thị có hướng không chu trình (luồng điều phối Airflow)\\
    DATN  & Đồ án tốt nghiệp\\
    dbt   & data build tool - công cụ biến đổi dữ liệu\\
    Dim   & Dimension - chiều dữ liệu\\
    DQ    & Data Quality - kiểm tra chất lượng dữ liệu\\
    DWH   & Data Warehouse - kho dữ liệu\\
    ELT   & Extract, Load, Transform - trích xuất, nạp, biến đổi\\
    ETL   & Extract, Transform, Load - trích xuất, biến đổi, nạp\\
    HSV   & Hue-Saturation-Value - không gian màu HSV\\
    KPI   & Key Performance Indicator - chỉ số hiệu suất chính\\
    KYC   & Know Your Customer - định danh khách hàng\\
    MCC   & Merchant Category Code - mã phân loại đơn vị chấp nhận thẻ\\
    MD5   & Message Digest 5 - hàm băm tạo hash key\\
    MNS   & Minus ghi nhận thay đổi (I/U/D) giữa snapshot hôm nay và snapshot lũy kế\\
    NHNN  & Ngân hàng Nhà nước\\
    NĐ-CP & Nghị định do Chính phủ ban hành (ký hiệu văn bản pháp luật)\\
    ISO 18245 & Tiêu chuẩn quốc tế quy định bộ mã phân loại đơn vị chấp nhận thẻ \\
    OCR   & Optical Character Recognition - nhận dạng ký tự quang học\\
    OLAP  & Online Analytical Processing - xử lý phân tích trực tuyến\\
    OLTP  & Online Transactional Processing - xử lý giao dịch trực tuyến\\
    PDY   & Previous Day - snapshot lũy kế ngày trước\\
    PCI DSS & Payment Card Industry Data Security Standard - tiêu chuẩn bảo mật dữ liệu thẻ thanh toán\\
    ROI   & Region of Interest - vùng quan tâm trên ảnh\\
    SCD   & Slowly Changing Dimension - chiều dữ liệu thay đổi chậm\\
    SCD2  & Slowly Changing Dimension Type 2 - SCD loại 2 (lưu lịch sử)\\
    SQL   & Structured Query Language - ngôn ngữ truy vấn có cấu trúc\\
    SSIS  & SQL Server Integration Services - công cụ tích hợp dữ liệu của SQL Server\\
    TDY   & Today - snapshot dữ liệu hôm nay\\
\end{longtable}

%-----------------------------------------------------
% MAIN MATTER (Đánh số Ả Rập, Bắt đầu có Header)
%-----------------------------------------------------
\clearpage
\pagenumbering{arabic} % Đánh số 1, 2, 3...
\setcounter{page}{1} % Bắt đầu lại từ trang 1
% Bật lại header cho phần thân. Phải ĐỊNH NGHĨA LẠI đầy đủ ở đây vì \pagestyle{plain}
% ở front matter (qua \fancyhf{} trong \fancypagestyle{plain}) đã xóa sạch nội dung header.
\pagestyle{fancy}
\fancyhf{}
\fancyhead[L]{\small\textit{Đồ án tốt nghiệp}} % Trái: Đồ án tốt nghiệp
\fancyhead[R]{\small\nouppercase{\leftmark}}   % Phải: Chương hiện tại và tên chương
\fancyfoot[C]{\thepage}                        % Giữa dưới: số trang
\renewcommand{\headrulewidth}{0.4pt}

\chapter{Mở đầu và cơ sở lý thuyết}
\label{ch:intro}

\section{Đặt vấn đề và lý do chọn đề tài}

\textbf{Bối cảnh dữ liệu trong lĩnh vực ngân hàng:} 
Trong kỷ nguyên số hóa, ngân hàng là một trong những ngành công nghiệp sản sinh ra khối lượng dữ liệu khổng lồ và đa dạng bậc nhất. Nguồn dữ liệu này không chỉ dừng lại ở các thông tin có cấu trúc rõ ràng như lịch sử giao dịch, thông tin hồ sơ khách hàng hay thẻ tín dụng, mà còn bao gồm một lượng lớn dữ liệu phi cấu trúc đến từ các chứng từ vật lý như Căn cước công dân (CCCD), sổ tiết kiệm hay hợp đồng. Điểm đặc trưng là khối dữ liệu khổng lồ này thường nằm phân tán ở nhiều hệ thống tác nghiệp độc lập, tồn tại dưới nhiều định dạng khác nhau, gây khó khăn cho việc quản lý tập trung.

\textbf{Các vấn đề hạn chế hiện tại:} 
Sự phân mảnh dữ liệu nêu trên đã và đang tạo ra nhiều rào cản lớn đối với các tổ chức tài chính. Việc thiếu vắng một nền tảng dữ liệu hợp nhất khiến quá trình khai thác, tổng hợp và phân tích thông tin toàn diện trở nên rời rạc. Bên cạnh đó, nhiều quy trình xử lý thông tin, đặc biệt là bóc tách dữ liệu từ chứng từ giấy, vẫn đang phải thực hiện hoàn toàn thủ công. Điều này không chỉ gây lãng phí thời gian mà còn tiềm ẩn rủi ro sai sót do yếu tố con người. Hơn thế nữa, hệ thống dữ liệu hiện hành thường thiếu khả năng lưu trữ lịch sử thay đổi của thực thể khiến việc truy vết trở nên bất khả thi. Sự thiếu hụt các cơ chế kiểm soát chất lượng dữ liệu (Data Quality) tự động cũng làm suy giảm độ tin cậy của thông tin đầu ra, ảnh hưởng trực tiếp đến các quyết định kinh doanh chiến lược.

\textbf{Nhu cầu thực tiễn:} 
Từ những bất cập thực tế đó, một nhu cầu cấp thiết được đặt ra là phải xây dựng một nền tảng dữ liệu (Data Platform) tập trung, hoạt động tự động và có kiến trúc dễ dàng mở rộng. Nền tảng này phải đảm nhiệm vai trò nền móng vững chắc để phục vụ cho các bài toán phân tích kinh doanh thông minh và hỗ trợ công tác quản trị dữ liệu của ngân hàng.

\textbf{Lý do chọn đề tài:} 
Xuất phát từ bối cảnh và những nhu cầu cấp thiết nêu trên, em đã quyết định lựa chọn đề tài nghiên cứu này. Các lý do chính bao gồm:
\begin{itemize}
    \item \textbf{Tính thực tiễn cao:} Đề tài bám sát và giải quyết một bài toán tích hợp dữ liệu đang thực sự tồn tại và gây đau đầu trong ngành tài chính - ngân hàng.
    \item \textbf{Cơ hội ứng dụng công nghệ hiện đại:} Đây là cơ hội tốt để nghiên cứu và đưa vào thực tiễn các công cụ, kiến trúc dữ liệu tiên tiến nhất hiện nay như Apache Airflow (điều phối), dbt (biến đổi dữ liệu), OCR (nhận dạng quang học) và kiến trúc đa tầng Medallion.
    \item \textbf{Khả năng kiểm chứng kết quả rõ ràng:} Sản phẩm của đồ án mang tính định lượng và trực quan cao, có thể dễ dàng kiểm chứng thông qua việc hệ thống Data Pipeline vận hành trơn tru và các chỉ số được trực quan hóa thành công trên Dashboard.
\end{itemize}

\section{Mục tiêu và phạm vi đồ án}

\subsection{Mục tiêu nghiên cứu}
Mục tiêu cốt lõi của đồ án là thiết kế và triển khai thành công một nền tảng dữ liệu (Data Platform) toàn diện từ đầu đến cuối (end-to-end) phục vụ chuyên sâu cho một nghiệp vụ ngân hàng. Để hiện thực hóa mục tiêu tổng thể này, hệ thống được thiết kế với hai luồng xử lý (pipeline) độc lập, bao gồm:
\begin{itemize}
    \item \textbf{Luồng dữ liệu có cấu trúc (Structured Pipeline):} Tự động hóa toàn bộ quy trình trích xuất, biến đổi và nạp dữ liệu. Luồng này vận hành từ hệ thống nguồn tác nghiệp (OLTP), đi qua vùng chờ Bronze, chuẩn hóa tại tầng Silver bằng mô hình Data Vault 2.0 và tổng hợp tại tầng Gold bằng mô hình Star Schema để sẵn sàng phục vụ phân tích.
    \item \textbf{Luồng dữ liệu phi cấu trúc (Unstructured Pipeline):} Ứng dụng công nghệ nhận dạng ký tự quang học (OCR) để tự động đọc và trích xuất thông tin từ các hình ảnh chứng từ vật lý (như Căn cước công dân, Sổ tiết kiệm), sau đó lưu trữ kết quả có cấu trúc vào tầng Bronze.
\end{itemize}
Bên cạnh đó, đồ án hướng tới việc ứng dụng kiến trúc đa tầng Medallion kết hợp cùng các mô hình dữ liệu tiên tiến nhằm đảm bảo tính toàn vẹn của lịch sử biến động dữ liệu và tính sẵn sàng mở rộng. Nền tảng cũng đặt mục tiêu tự động hóa hoàn toàn công tác điều phối, kiểm soát chặt chẽ chất lượng dữ liệu và cuối cùng là trực quan hóa các kết quả phân tích thành các Dashboard sinh động trên công cụ Power BI.

\subsection{Phạm vi và giới hạn của đồ án}
Mặc dù hướng tới một kiến trúc nền tảng dữ liệu hiện đại, song do giới hạn về mặt thời gian và tài nguyên thực hiện của một đồ án tốt nghiệp, đề tài được giới hạn trong các phạm vi cụ thể như sau:
\begin{itemize}
    \item \textbf{Về nguồn dữ liệu:} Hệ thống sử dụng các tập dữ liệu giả lập và dữ liệu đã được ẩn danh hóa mô phỏng theo nghiệp vụ ngân hàng thực tế. Hệ thống không kết nối và không sử dụng dữ liệu thật từ môi trường vận hành (production) nhằm tuân thủ các nguyên tắc bảo mật thông tin.
    \item \textbf{Về phương thức xử lý:} Các luồng xử lý (pipeline) được thiết kế theo cơ chế xử lý theo lô (Batch Processing), hoạt động tự động theo lịch trình định sẵn. Hệ thống hiện tại chưa hỗ trợ phương thức xử lý dữ liệu theo thời gian thực (Real-time/Streaming processing).
    \item \textbf{Về luồng dữ liệu phi cấu trúc (OCR):} Trong khuôn khổ một đồ án về nền tảng dữ liệu, mục tiêu của luồng này không phải là nghiên cứu hay huấn luyện chuyên sâu các mô hình trí tuệ nhân tạo, mà là kế thừa, ứng dụng các mô hình OCR có sẵn vào bài toán xử lý và trích xuất thông tin từ chứng từ. Do đó, luồng nhận dạng hình ảnh chỉ dừng ở bước trích xuất thành công các trường thông tin cá nhân và lưu trữ kết quả tại tầng Bronze. Các dữ liệu này không được tích hợp tiếp vào luồng phân tích (Silver/Gold), bởi lẽ các thông tin định danh tương tự đã tồn tại sẵn ở nguồn cơ sở dữ liệu có cấu trúc; việc tích hợp ép buộc sẽ sinh ra hiện tượng dư thừa, trùng lặp và làm ảnh hưởng đến tính nhất quán của các insight phân tích đầu ra.

    \item \textbf{Về hạ tầng triển khai:} Toàn bộ nền tảng dữ liệu, từ cơ sở dữ liệu đến các công cụ điều phối, được đóng gói và vận hành cục bộ (On-premise/Local) thông qua Docker. Đồ án chưa triển khai mở rộng lên các nền tảng điện toán đám mây (Cloud Computing).
\end{itemize}

\section{Cơ sở lý thuyết}
\label{sec:background}
% Chỉ trình bày kiến thức nền tảng (mức khái niệm) dùng trong đồ án. Lý thuyết
% gắn với công cụ cụ thể (Airflow, dbt, PaddleOCR, Docker, Power BI, Data
% Quality) được trình bày ở Chương 4, mục "Công nghệ sử dụng".

\subsection{Tổng quan kiến thức về kho dữ liệu}

\subsubsection*{a) Vòng đời dữ liệu}

Data Engineer là kỹ sư phần mềm chuyên về dữ liệu, chịu trách nhiệm xây dựng và duy trì hạ tầng để dữ liệu có thể được thu thập, lưu trữ, xử lý và truy cập một cách đáng tin cậy. Nhiệm vụ cốt lõi của vị trí này bao gồm thiết kế và xây dựng các luồng dữ liệu tự động, thiết kế mô hình kho dữ liệu phù hợp với nhu cầu phân tích, đảm bảo chất lượng, tính nhất quán, khả năng truy vết dữ liệu, cũng như vận hành, giám sát và xử lý sự cố hệ thống liên tục.

Để thực hiện tốt các nhiệm vụ trên, việc nắm bắt vòng đời dữ liệu là vô cùng quan trọng. Vòng đời này diễn ra theo một trình tự tuyến tính bao gồm năm giai đoạn chính. Giai đoạn đầu tiên là thu thập, nơi dữ liệu được kéo về từ nhiều nguồn khác nhau như cơ sở dữ liệu, API, file hay thiết bị IoT. Tiếp đó, dữ liệu được đưa vào lưu trữ thô nhằm giữ nguyên dạng gốc trước khi bị can thiệp. Ở giai đoạn thứ ba, dữ liệu trải qua quá trình xử lý và làm sạch, bao gồm các bước chuẩn hóa, kiểm tra chất lượng và mô hình hóa. Sau khi đã hoàn thiện, dữ liệu bước vào giai đoạn phục vụ để cung cấp cho người dùng cuối thông qua các hệ thống BI, Machine Learning hoặc API. Cuối cùng, dữ liệu sẽ được đưa vào lưu trữ lâu dài hoặc xóa bỏ dựa trên các chính sách bảo mật và tuân thủ quy định.

\begin{figure}[H]
    \centering
    \includegraphics[width=0.8\textwidth]{img1/Vòng đời data.jpg}
    \caption{Vòng đời dữ liệu}
    \label{fig:data-lifecycle}
\end{figure}


Trong quá trình vận hành luồng dữ liệu, có hai phương pháp xử lý phổ biến nhất là xử lý theo lô (Batch processing) và xử lý luồng (Streaming processing)

\begin{table}[H]
    \centering
    \renewcommand{\arraystretch}{1.3}
    \begin{tabular}{|p{3.5cm}|p{5cm}|p{5.5cm}|}
        \hline
        \textbf{Tiêu chí} & \textbf{Batch} & \textbf{Streaming} \\ \hline
        \textbf{Độ trễ} & Phút $\rightarrow$ Giờ & Mili giây $\rightarrow$ Giây \\ \hline
        \textbf{Chi phí} & Thấp & Cao hơn \\ \hline
        \textbf{Độ phức tạp} & Thấp & Cao \\ \hline
        \textbf{Use case} & Báo cáo, phân tích lịch sử & Cảnh báo, fraud detection \\ \hline
        \textbf{Công cụ điển hình} & Airflow, Spark & Kafka, Flink, Spark Streaming \\ \hline
    \end{tabular}
    \caption{So sánh phương pháp xử lý Batch và Streaming}
\end{table}

Liên quan đến quy trình tích hợp và luân chuyển dữ liệu vào kho, hai phương pháp tiếp cận chính thường được cân nhắc là ETL (Extract - Transform - Load) và ELT (Extract - Load - Transform). Với quy trình ETL truyền thống, dữ liệu được trích xuất và biến đổi hoàn toàn ở bên ngoài (thông qua các công cụ riêng biệt như Informatica hay SSIS) trước khi nạp vào kho. Phương pháp này rất lý tưởng khi hệ thống kho đích có tài nguyên hạn chế hoặc khi doanh nghiệp cần kiểm soát chặt chẽ tính bảo mật của dữ liệu đầu vào. Tuy nhiên, hạn chế của ETL là không lưu lại dữ liệu thô gốc, dẫn đến khó khăn trong việc tái xử lý nếu logic nghiệp vụ thay đổi. Để khắc phục điều này, quy trình ELT đảo ngược thứ tự bằng cách nạp trực tiếp dữ liệu thô vào kho (thường là các nền tảng dữ liệu đám mây hiện đại), sau đó tận dụng sức mạnh tính toán của chính hệ thống này để thực hiện biến đổi thông qua các câu lệnh SQL hoặc công cụ như dbt. Nhờ giữ lại toàn bộ dữ liệu gốc, ELT mang lại sự linh hoạt tối đa khi cần tái xử lý và phục hồi dữ liệu.

\begin{table}[H]
    \centering
    \renewcommand{\arraystretch}{1.3}
    \begin{tabular}{|p{3cm}|p{5cm}|p{6cm}|}
        \hline
        \textbf{Tiêu chí} & \textbf{ETL} & \textbf{ELT} \\ \hline
        \textbf{Thứ tự bước} & Extract $\rightarrow$ Transform $\rightarrow$ Load & Extract $\rightarrow$ Load $\rightarrow$ Transform \\ \hline
        \textbf{Vị trí biến đổi} & Ngoài kho (công cụ riêng) & Trong kho (SQL/dbt) \\ \hline
        \textbf{Dữ liệu thô} & Không lưu & Lưu lại \\ \hline
        \textbf{Phù hợp} & Kho truyền thống, tài nguyên hạn chế & Kho hiện đại, cloud \\ \hline
    \end{tabular}
    \caption{So sánh quy trình tích hợp dữ liệu ETL và ELT}
\end{table}


\subsubsection*{b) Kho dữ liệu và hệ sinh thái}

Trước khi đi sâu vào sự tiến hóa của các kiến trúc nền tảng dữ liệu, cần phân biệt rõ hai hệ thống xử lý cốt lõi, vốn là tiền đề cho sự ra đời của các Kho dữ liệu: OLTP và OLAP.

\begin{itemize}
    \item \textbf{OLTP (Online Transaction Processing - Xử lý giao dịch trực tuyến):} Là hệ thống phục vụ các tác vụ vận hành và giao dịch hằng ngày của doanh nghiệp. Điểm đặc trưng của OLTP là được tối ưu hóa để ghi và đọc nhanh từng bản ghi đơn lẻ với tần suất liên tục (ví dụ: tạo một giao dịch chuyển tiền, cập nhật số dư). Dữ liệu trong hệ thống này được chuẩn hóa ở mức độ cao (thường là chuẩn 3NF) nhằm tránh dư thừa và đảm bảo tính toàn vẹn của các giao dịch.
    \item \textbf{OLAP (Online Analytical Processing - Xử lý phân tích trực tuyến):} Là hệ thống chuyên phục vụ công tác phân tích dữ liệu và hỗ trợ ra quyết định. Khác với OLTP, OLAP được tối ưu hóa để thực hiện các câu lệnh truy vấn đọc, quét và tính toán tổng hợp trên một khối lượng dữ liệu khổng lồ. Dữ liệu trong OLAP thường được phi chuẩn hóa và tổ chức theo các mô hình đa chiều (như Star Schema hoặc Snowflake Schema) để tăng tốc độ truy xuất cho các báo cáo kinh doanh.
\end{itemize}

Chính vì sự khác biệt hoàn toàn về mục đích và cấu trúc lưu trữ này, hệ thống phân tích (OLAP) cần phải được tách biệt khỏi hệ thống giao dịch (OLTP). Để đáp ứng nhu cầu phân tích ngày càng phức tạp, kiến trúc nền tảng dữ liệu đã trải qua ba giai đoạn tiến hóa mang tính cách mạng: từ Data Warehouse, Data Lake cho đến kiến trúc Lakehouse hiện đại.

\begin{figure}[H]
    \centering
    \includegraphics[width=0.8\textwidth]{img1/Các loại kiến trúc .png}
    \caption{Tiến trình tiến hóa từ kho dữ liệu truyền thống đến kiến trúc Lakehouse}
    \label{fig:data-architecture-evolution}
\end{figure}

\begin{table}[H]
    \centering
    \renewcommand{\arraystretch}{1.4}
    \resizebox{\textwidth}{!}{%
    \begin{tabular}{|p{3cm}|p{3.8cm}|p{3.8cm}|p{3.8cm}|}
        \hline
        \textbf{Tiêu chí} & \textbf{Data Warehouse} & \textbf{Data Lake} & \textbf{Lakehouse} \\ \hline
        \textbf{Loại dữ liệu} & Dữ liệu có cấu trúc & Mọi loại dữ liệu (Thô, phi cấu trúc) & Đa dạng (Cả cấu trúc và phi cấu trúc) \\ \hline
        \textbf{Cơ chế Schema} & Schema-on-write (Cố định trước khi nạp) & Schema-on-read (Linh hoạt khi truy vấn) & Kết hợp, hỗ trợ tính toàn vẹn (ACID) \\ \hline
        \textbf{Mục đích chính} & Phân tích BI, báo cáo & Lưu trữ gốc, Khoa học dữ liệu (ML/AI) & Hợp nhất BI và Machine Learning \\ \hline
        \textbf{Nhược điểm} & Chi phí cao, khó xử lý dữ liệu phi cấu trúc & Dễ trở thành "Data Swamp" do thiếu quản trị & Đòi hỏi công nghệ vận hành và thiết lập phức tạp \\ \hline
    \end{tabular}}
    \caption{So sánh đặc điểm của Data Warehouse, Data Lake và Lakehouse}
\end{table}

% TODO: OLTP vs OLAP; Data Warehouse / Data Lake / Lakehouse.

\subsection{Mô hình hóa dữ liệu }
Để quản lý và tổ chức luồng dữ liệu hiệu quả bên trong các hệ thống nền tảng dữ liệu hiện đại, đặc biệt là trên kiến trúc Lakehouse, \textbf{kiến trúc đa tầng Medallion (Medallion Architecture)} hiện đang được áp dụng rộng rãi như một tiêu chuẩn thiết kế cốt lõi. Về mặt bản chất, Medallion là mô hình tổ chức kho lưu trữ thành nhiều tầng kế tiếp nhau, trong đó mỗi tầng đại diện cho một mức độ gia tăng về chất lượng và độ phức tạp trong xử lý. 

Theo mô hình này, dòng chảy dữ liệu sẽ đi qua một hành trình tinh luyện xuyên suốt: bắt đầu từ trạng thái thô nguyên bản (tầng Bronze), tiếp tục được làm sạch, chuẩn hóa và hợp nhất (tầng Silver), cho đến khi trở thành nguồn dữ liệu chất lượng cao, sẵn sàng phục vụ trực tiếp cho các nghiệp vụ phân tích và ra quyết định (tầng Gold). Việc chia tách rõ ràng thành các tầng chất lượng không chỉ giúp hệ thống kiểm soát và truy vết dữ liệu dễ dàng hơn, mà còn tạo nền móng vững chắc để áp dụng các kỹ thuật mô hình hóa dữ liệu chuyên sâu tại từng tầng tương ứng.

\begin{figure}[H]
    \centering
    \includegraphics[width=0.8\textwidth]{img1/Medallion architecture.png}
    \caption{Kiến trúc Medallion Architecture}
    \label{fig:medallion-architecture}
\end{figure}

\subsubsection*{a) Mô hình: Data Vault 2.0}
% TODO: Hub / Link / Satellite; Business key & Hash key (MD5); Satellite SCD2.


Data Vault 2.0 là phương pháp mô hình hóa dữ liệu được thiết kế chuyên biệt cho tầng lưu trữ trung gian (tầng Silver) trong kiến trúc nền tảng dữ liệu. Mục tiêu cốt lõi của phương pháp này là cung cấp khả năng lưu trữ toàn bộ lịch sử thay đổi của dữ liệu, đồng thời dễ dàng mở rộng và tích hợp thông tin từ nhiều nguồn khác nhau. Nếu như mô hình hình sao (Star Schema) được tối ưu hóa cho tốc độ truy vấn phân tích, thì Data Vault lại đặt ưu tiên hàng đầu vào tính linh hoạt và khả năng kiểm toán (audit) dữ liệu.

Kiến trúc của một mô hình Data Vault 2.0 được cấu thành từ ba thành phần cốt lõi, hoạt động gắn kết chặt chẽ với nhau. \begin{itemize}
    \item \textbf{Hub (Trung tâm):}
    \begin{itemize}
        \item Lưu trữ danh sách các khóa nghiệp vụ (Business Key) duy nhất của một thực thể (VD: mã khách hàng, mã thẻ).
        \item Tuyệt đối không chứa thuộc tính mô tả.
        \item Chỉ chứa khóa nghiệp vụ và các siêu dữ liệu metadata (ngày nạp, nguồn gốc).
        \item Mỗi thực thể nghiệp vụ (business entity) đều có một Hub riêng biệt.
    \end{itemize}
    
    \item \textbf{Link (Liên kết):}
    \begin{itemize}
        \item Lưu trữ mối quan hệ giữa các Hub với nhau (VD: khách hàng — thẻ, thẻ — giao dịch).
        \item Xử lý các quan hệ nhiều-nhiều một cách tự nhiên, loại bỏ sự cần thiết của các bảng trung gian phức tạp.
        \item Chỉ chứa giá trị băm (hash key) của các Hub liên quan kèm theo metadata.
    \end{itemize}
    
    \item \textbf{Satellite (Vệ tinh):}
    \begin{itemize}
        \item Lưu trữ toàn bộ các thuộc tính mô tả chi tiết của một Hub hoặc một Link (VD: tên, địa chỉ, số dư...).
        \item Một Hub hoặc Link có thể liên kết với nhiều Satellite, được phân chia theo nguồn cấp hoặc nhóm thuộc tính.
        \item Là thành phần duy nhất chịu trách nhiệm lưu trữ lịch sử thay đổi của dữ liệu theo thời gian 
    \end{itemize}
\end{itemize}
\begin{figure}[H]
    \centering
    \includegraphics[width=0.8\textwidth]{img1/Data vault.png}
    \caption{Mô hình Data Vault 2.0}
    \label{fig:data-vault-model}
\end{figure}
\subsubsection*{b) Mô hình chiều (Star Schema)}

Star Schema (mô hình hình sao) là một phương pháp mô hình hóa dữ liệu được sử dụng chủ yếu cho tầng phân tích (tầng Gold). Phương pháp này được thiết kế để tối ưu hóa cho các truy vấn tổng hợp dữ liệu với tốc độ cao, đồng thời sở hữu cấu trúc trực quan, dễ hiểu đối với người dùng cuối. Cấu trúc cốt lõi của mô hình bao gồm một bảng sự kiện (Fact table) được đặt ở trung tâm, bao quanh bởi các bảng chiều (Dimension table), tạo nên hình dáng giống như một ngôi sao khi nhìn từ trên xuống.

Bảng sự kiện (Fact Table) đóng vai trò trung tâm trong mô hình, làm nhiệm vụ lưu trữ các sự kiện nghiệp vụ đã thực sự xảy ra trong hệ thống như các giao dịch, đơn đặt hàng hay lượt truy cập. Đặc điểm điển hình của bảng sự kiện là có dung lượng rất lớn và tăng trưởng liên tục theo thời gian. Về mặt cấu trúc, bảng này chứa các khóa ngoại trỏ đến các bảng chiều tương ứng, kèm theo các độ đo (measure) mang giá trị định lượng như số tiền hoặc số lượng. Dựa trên đặc thù lưu trữ, bảng sự kiện thường được phân thành các loại sau:
\begin{itemize}
    \item \textbf{Transaction fact (Sự kiện giao dịch):} Mỗi bản ghi đại diện cho một sự kiện đơn lẻ xảy ra tại một thời điểm cụ thể.
    \item \textbf{Periodic snapshot (Chụp nhanh định kỳ):} Ghi nhận và lưu trữ trạng thái của hệ thống theo một chu kỳ thời gian nhất định, ví dụ như số dư tài khoản vào cuối mỗi ngày.
    \item \textbf{Accumulating snapshot (Chụp nhanh lũy kế):} Được sử dụng để theo dõi toàn bộ tiến trình của một quy trình nghiệp vụ bao gồm nhiều bước nối tiếp nhau.
\end{itemize}

Bên cạnh bảng sự kiện và các bảng chiều nghiệp vụ, chiều thời gian (\texttt{dim\_date}) là một Dimension đặc biệt và gần như bắt buộc phải có trong mọi thiết kế Star Schema. Bảng này không chỉ lưu trữ thông tin ngày tháng đơn thuần mà còn chứa các thuộc tính phái sinh rất hữu ích cho việc phân tích như: thứ trong tuần, tuần trong năm, quý, cờ đánh dấu ngày lễ hoặc cờ đánh dấu ngày cuối tuần. Việc chuẩn hóa chiều thời gian cho phép hệ thống linh hoạt thực hiện các truy vấn phân tích theo nhiều mức độ chi tiết (granularity) khác nhau mà không cần phải thực hiện tính toán lại từ đầu.
\begin{figure}[H]
    \centering
    \includegraphics[width=0.8\textwidth]{img1/Star Schema.png}
    \caption{Mô hình Data Warehouse: Sơ đồ ngôi sao}
    \label{fig:star-schema}
\end{figure}

\subsection{Nạp dữ liệu gia tăng, CDC và SCD}

\subsubsection*{a) Nạp dữ liệu gia tăng (Incremental Loading)}

Trong quá trình luân chuyển và tích hợp dữ liệu vào kho, các hệ thống thường áp dụng hai chiến lược nạp dữ liệu cơ bản là nạp toàn bộ (Full Load) và nạp gia tăng (Incremental Load). Phương pháp Full Load thực hiện nạp lại toàn bộ dữ liệu từ hệ thống nguồn và ghi đè lên kho dữ liệu trong mỗi lần chạy. Mặc dù chiến lược này có ưu điểm là triển khai đơn giản, nó lại tiêu tốn rất nhiều tài nguyên tính toán và không thể lưu giữ được lịch sử thay đổi của hệ thống. 

Để khắc phục nhược điểm đó, chiến lược Incremental Load được áp dụng nhằm chỉ nạp những phần dữ liệu mới phát sinh hoặc có sự thay đổi kể từ lần nạp trước đó. Cách tiếp cận này giúp hệ thống tiết kiệm đáng kể tài nguyên xử lý, đồng thời duy trì và giữ được lịch sử biến động của dữ liệu theo thời gian.

\vspace{0.5em}
\subsubsection*{b) Change Data Capture (CDC)}

Để hiện thực hóa được quá trình nạp dữ liệu gia tăng, hệ thống cần ứng dụng kỹ thuật Change Data Capture (CDC). CDC là kỹ thuật nhằm phát hiện và ghi lại các thao tác thay đổi xảy ra trên dữ liệu nguồn, bao gồm các hành động Thêm mới (Insert), Cập nhật (Update) và Xóa (Delete). Mục tiêu cốt lõi của CDC là giúp hệ thống biết được chính xác bản ghi nào đã bị thay đổi, nội dung thay đổi là gì và diễn ra vào lúc nào, từ đó chỉ tập trung xử lý đúng phần dữ liệu có biến động.

Hiện nay, có bốn phương pháp triển khai CDC phổ biến với các cơ chế và ưu nhược điểm khác nhau:

\begin{table}[H]
    \centering
    \renewcommand{\arraystretch}{1.4}
    \resizebox{\textwidth}{!}{%
    \begin{tabular}{|p{3cm}|p{4cm}|p{3.5cm}|p{4cm}|}
        \hline
        \textbf{Phương pháp} & \textbf{Cơ chế} & \textbf{Ưu điểm} & \textbf{Nhược điểm} \\ \hline
        \textbf{Timestamp-based} & So sánh dựa trên cột \texttt{updated\_at} & Triển khai đơn giản & Dễ bỏ sót các bản ghi bị xóa (Delete) \\ \hline
        \textbf{Snapshot/Diff} & So sánh hai bản chụp (snapshot) liên tiếp & Không yêu cầu sửa đổi hệ thống nguồn & Tiêu tốn nhiều tài nguyên để so sánh \\ \hline
        \textbf{Log-based} & Đọc trực tiếp từ transaction log của Database & Đầy đủ thông tin nhất, tiêu tốn ít tài nguyên nguồn & Phụ thuộc nhiều vào kiến trúc của hệ thống Database \\ \hline
        \textbf{Trigger-based} & Gắn Trigger để ghi lại các thay đổi vào bảng audit & Độ chính xác cao & Làm tăng tải hoạt động của hệ thống nguồn \\ \hline
    \end{tabular}}
    \caption{So sánh các phương pháp Change Data Capture (CDC) phổ biến}
\end{table}

\vspace{0.5em}
\subsubsection*{c) Slowly Changing Dimension (SCD)}

Sau khi dữ liệu thay đổi được CDC bắt giữ và đẩy về kho, hệ thống phải đối mặt với bài toán Slowly Changing Dimension (SCD) — đây là bài toán chuyên xử lý sự thay đổi của dữ liệu chiều (Dimension) theo thời gian trong kho dữ liệu. Câu hỏi cốt lõi được đặt ra là: khi thông tin của một thực thể thay đổi (ví dụ: khách hàng đổi địa chỉ, nâng hạng thẻ), hệ thống nên ghi đè thông tin mới hay giữ lại lịch sử? Tùy thuộc vào nhu cầu phân tích, kho dữ liệu thường áp dụng ba phương pháp chính:

\begin{itemize}
    \item \textbf{SCD Type 1 (Ghi đè):} Phương pháp này cập nhật trực tiếp giá trị mới đè lên giá trị cũ. Do không lưu lại lịch sử, cách làm này rất đơn giản nhưng lại làm mất đi thông tin trong quá khứ. SCD Type 1 thường chỉ được dùng khi lịch sử thay đổi của thuộc tính không mang lại giá trị phân tích.
    
    \item \textbf{SCD Type 2 (Thêm bản ghi mới):} Khi có sự thay đổi, hệ thống sẽ thực hiện đóng bản ghi cũ và thêm một bản ghi mới chứa giá trị vừa được cập nhật. Cách tiếp cận này giúp lưu trữ được toàn bộ lịch sử thay đổi của dữ liệu. Để quản lý các phiên bản, bảng dữ liệu thường sử dụng thêm các cột như \texttt{effective\_from}, \texttt{effective\_to} nhằm xác định đâu là phiên bản hiện hành. SCD Type 2 rất phù hợp cho các bài toán cần phân tích theo trạng thái dữ liệu tại một thời điểm cụ thể trong quá khứ.
    
    \item \textbf{SCD Type 3 (Thêm cột):} Thay vì thêm dòng, phương pháp này giữ cả giá trị cũ và giá trị mới trong cùng một bản ghi bằng cách thêm cột mới (ví dụ: cột \texttt{previous\_value}). Phương pháp này có nhiều hạn chế vì nó chỉ có khả năng lưu lại được một lần thay đổi gần nhất của thực thể.
\end{itemize}
% TODO: Change Data Capture; mẫu snapshot/diff (tdy--pdy--mns), cờ I/U/D; idempotency.
\subsection{Các khái niệm cốt lõi trong bài toán xử lý ảnh và OCR}

Để trích xuất thông tin chính xác từ các chứng từ vật lý (như CCCD, Sổ tiết kiệm), hệ thống cần áp dụng chuỗi các kỹ thuật thị giác máy tính nhằm tiền xử lý ảnh và nhận dạng ký tự. Dưới đây là các khái niệm và kỹ thuật nền tảng được sử dụng:

\begin{itemize}
    \item \textbf{Khái niệm OCR:} 
    Là công nghệ chuyển đổi hình ảnh chứa văn bản (ảnh scan, ảnh chụp) thành định dạng văn bản có thể máy đọc và xử lý được. Quá trình OCR hiện đại thường được chia làm hai giai đoạn độc lập:
    \begin{itemize}
        \item \textit{Text Detection (Định vị văn bản):} Quét bức ảnh để xác định tọa độ của các vùng chứa chữ.
        \item \textit{Text Recognition (Nhận dạng văn bản):} Trích xuất và giải mã các hình ảnh cắt từ vùng chứa chữ thành các ký tự (string) thực tế.
    \end{itemize}

    \item \textbf{Không gian màu BGR, Grayscale và HSV:}
    \begin{itemize}
        \item \textit{BGR (Blue - Green - Red):} Là không gian màu mặc định được thư viện OpenCV sử dụng khi đọc ảnh (ngược với thứ tự RGB thông thường).
        \item \textit{Grayscale (Ảnh xám):} Biến đổi ảnh màu 3 kênh thành ảnh 1 kênh duy nhất mang giá trị cường độ sáng từ 0 (đen) đến 255 (trắng). Phép biến đổi này giúp giảm thiểu khối lượng tính toán và tập trung vào các đặc trưng hình học của nét chữ.
        \item \textit{HSV (Hue - Saturation - Value):} Phân tách hình ảnh thành Màu sắc, Độ bão hòa và Độ sáng. Không gian này đặc biệt hữu ích khi cần tách nền chứng từ dựa trên đặc trưng màu sắc hoặc xử lý các ảnh bị đổ bóng mờ.
    \end{itemize}

    \item \textbf{Tăng cường độ tương phản (CLAHE):}
    CLAHE là kỹ thuật cân bằng biểu đồ mức xám có giới hạn độ tương phản thích nghi. Thay vì áp dụng trên toàn bộ ảnh, CLAHE chia ảnh thành các khối nhỏ và cân bằng sáng cho từng khối. Điều này giúp làm nổi bật nét chữ trên các vùng nền tối mà không làm khuếch đại nhiễu ở các vùng nền đồng nhất.

    \item \textbf{Khử nhiễu qua thuật toán Non-local Means}: 
    Đây là thuật toán loại bỏ nhiễu hạt trên ảnh rất hiệu quả. Khác với các bộ lọc thông thường chỉ làm mờ ảnh xung quanh điểm xét, Non-local Means tính toán giá trị trung bình dựa trên tất cả các điểm ảnh có vùng lân cận tương đồng nhau trên toàn bộ bức ảnh. Kỹ thuật này giúp làm sạch nền chứng từ nhưng vẫn giữ được độ sắc nét tuyệt đối cho các cạnh và viền của chữ.

    \item \textbf{Ngưỡng hóa thích nghi (Adaptive Thresholding):}
    Là phương pháp chuyển ảnh xám thành ảnh nhị phân (chỉ có 2 màu đen/trắng). Do ảnh chụp thực tế thường có ánh sáng không đều (chỗ bị chói, chỗ đổ bóng), thuật toán này sẽ tự động tính toán một ngưỡng (threshold) riêng biệt cho từng vùng nhỏ cục bộ thay vì dùng một ngưỡng cố định cho toàn ảnh, từ đó giúp tách chữ ra khỏi nền một cách triệt để.

    \item \textbf{Phát hiện cạnh với thuật toán Canny:}
    Là kỹ thuật trích xuất các đường viền của đối tượng dựa trên sự thay đổi đột ngột về cường độ sáng của điểm ảnh. Trong bài toán chứng từ, thuật toán Canny đóng vai trò then chốt để phát hiện ra các khung viền của thẻ CCCD hoặc khung viền của từng dòng chữ.

    \item \textbf{Biến đổi phối cảnh:}
    Hình ảnh chụp từ điện thoại thường bị méo do góc chụp nghiêng. Phép biến đổi phối cảnh là một phép biến đổi hình học tính toán ma trận dịch chuyển dựa trên 4 điểm góc của chứng từ, qua đó "nắn" bức ảnh về dạng hình chữ nhật phẳng với góc nhìn thẳng từ trên xuống (bird's-eye view), khôi phục tỷ lệ thực của tài liệu.

    \item \textbf{Chỉnh nghiêng (Deskew):}
    Để OCR đọc chính xác, các dòng chữ phải nằm ngang. Hệ thống sử dụng phép biến đổi Hough (Hough Transform) để tìm kiếm các đường thẳng dài trong ảnh (như dòng kẻ, lề văn bản). Dựa trên góc nghiêng của các đường này, hệ thống sẽ tính toán góc bù và xoay lại toàn bộ bức ảnh về góc 0 độ chuẩn xác.

    \item \textbf{Điểm tin cậy OCR (Confidence Score):}
    Là chỉ số đánh giá (thường dao động từ 0 đến 1, hoặc 0-100\%) thể hiện mức độ "chắc chắn" của mô hình AI đối với kết quả dự đoán của một từ hoặc một ký tự. Chỉ số này rất quan trọng để thiết lập các ngưỡng tự động hóa: nếu điểm tin cậy thấp hơn ngưỡng cho phép, hệ thống sẽ gắn cờ cảnh báo (flag) yêu cầu con người tham gia đối chiếu lại.
\end{itemize}
\chapter{Khảo sát đề tài}
\label{ch:survey}

\section{Khảo sát bài toán nghiệp vụ và mô hình dữ liệu}

Môi trường vận hành của một ngân hàng hiện đại thường bao gồm rất nhiều hệ thống hoạt động song song nhằm đáp ứng các nghiệp vụ khác nhau, điển hình như hệ thống Core Banking (xử lý giao dịch), CRM (quản lý quan hệ khách hàng), Card Management (quản lý thẻ) và KYC/Document Management (quản lý chứng từ định danh). Đặc thù này dẫn đến việc dữ liệu bị phân tán ở nhiều nơi và thiếu tính đồng nhất, gây ra rào cản rất lớn cho việc tổng hợp và phân tích thông tin một cách toàn diện. Từ thực trạng đó, nhu cầu cấp thiết là phải hợp nhất các nguồn dữ liệu này vào một nền tảng trung tâm. 

Dựa trên nguồn dữ liệu được hợp nhất, hệ thống nền tảng dữ liệu trong đồ án được định hướng tập trung giải quyết chủ đề phân tích cốt lõi về \textbf{Phân khúc khách hàng và Marketing}. Mục tiêu nghiệp vụ của bài toán này là xác định các nhóm khách hàng tiềm năng nhằm phục vụ chiến lược up-sell các loại thẻ hoặc mời chào các gói vay. Quá trình phân tích sẽ dựa trên hành vi chi tiêu và hồ sơ tài chính của khách hàng. Kết quả đầu ra (insight) thu được từ hệ thống sẽ phục vụ trực tiếp cho các phòng ban chức năng, bao gồm Phòng Marketing, Quản lý sản phẩm (Thẻ \& Vay), Quản lý quan hệ khách hàng (Relationship Management) và Ban lãnh đạo.
\subsection{Dữ liệu có cấu trúc}

Trong bức tranh tổng thể, dữ liệu có cấu trúc đóng vai trò xương sống cho hầu hết các luồng phân tích nghiệp vụ ngân hàng. Đây là các dữ liệu được định nghĩa rõ ràng, lưu trữ dưới dạng bảng và có tính quan hệ chặt chẽ, bao gồm các thực thể chính sau:

\begin{itemize}
    \item \textbf{Khách hàng (\texttt{customers}):} Dữ liệu được trích xuất từ hệ thống CRM và Core Banking, đóng vai trò cung cấp thông tin nền tảng để phân khúc thu nhập, đánh giá rủi ro tín dụng và phân tích hành vi chi tiêu của từng cá nhân.
    \item \textbf{Thẻ (\texttt{cards}):} Nguồn dữ liệu từ hệ thống Card Management, cung cấp các thông số chi tiết như loại thẻ, hạn mức khả dụng, tuổi đời của thẻ, tình trạng bảo mật (PIN) và tình trạng sử dụng chip vật lý.
    \item \textbf{Giao dịch (\texttt{transactions}):} Thu thập từ Core Banking, đây là thực thể có dung lượng lớn nhất và phát sinh liên tục trong hệ thống. Nó phản ánh toàn bộ các hoạt động tài chính như doanh số, các lỗi giao dịch phát sinh và vị trí thực hiện.
    \item \textbf{Đơn vị chấp thuận thanh toán (\texttt{merchants}) và Danh mục MCC (\texttt{mcc\_codes}):} Dữ liệu đơn vị chấp thuận thanh toán giúp ngân hàng theo dõi sự phân bổ địa lý của các giao dịch và phát hiện các mẫu bất thường. Đồng thời, mã danh mục đơn vị chấp nhận thẻ (MCC) được phân loại theo tiêu chuẩn ISO 18245 giúp hệ thống dễ dàng nhóm các hành vi chi tiêu của khách hàng theo từng ngành nghề kinh doanh cụ thể.
\end{itemize}

\subsection{Dữ liệu phi cấu trúc}

Bên cạnh nguồn dữ liệu có cấu trúc, các ngân hàng còn sở hữu một khối lượng khổng lồ dữ liệu phi cấu trúc, chủ yếu tồn tại dưới dạng hình ảnh scan từ các quy trình nghiệp vụ vật lý. Điển hình trong phạm vi đồ án này là hai loại tài liệu:

\begin{itemize}
    \item \textbf{Căn cước công dân (CCCD):} Đây là tài liệu thu thập từ quy trình nhận biết khách hàng (KYC) khi người dùng yêu cầu mở tài khoản. Các thông tin trọng yếu cần bóc tách bao gồm: họ tên, ngày sinh, số CCCD, địa chỉ và ngày cấp.
    \item \textbf{Sổ tiết kiệm:} Tương ứng với các chứng từ giao dịch gửi tiết kiệm truyền thống, yêu cầu trích xuất các thông tin cốt lõi như số tài khoản, số tiền, mức lãi suất và kỳ hạn.
\end{itemize}

\begin{figure}[H]
    \centering
    \includegraphics[width=0.7\textwidth]{img2/cccd_chip_0406085313.png }
    \caption{Mẫu hình ảnh CCCD ngoài thực tế}
    \label{fig:real-cccd-sample}
\end{figure}

\begin{figure}[H]
    \centering
    \includegraphics[width=0.7\textwidth]{img2/Sổ tiết kiệm.png }
    \caption{Mẫu hình ảnh sổ tiết kiệm ngoài thực tế}
    \label{fig:real-savings-book-sample}
\end{figure}

Đặc điểm lớn nhất của loại hình dữ liệu này là chúng tồn tại dưới dạng ảnh pixels, dẫn đến việc không thể thực hiện các câu truy vấn trực tiếp. Để khai thác, hệ thống bắt buộc phải trải qua bước tiền xử lý bằng công nghệ nhận dạng ký tự quang học (OCR) nhằm chuyển đổi hình ảnh thành văn bản.


\section{Khảo sát hiện trạng và các vấn đề tồn đọng}

Qua khảo sát, hệ thống dữ liệu hiện tại của ngân hàng đang bộc lộ nhiều hạn chế cốt lõi, ảnh hưởng trực tiếp đến khả năng vận hành và phân tích:

\begin{itemize}
    \item \textbf{Dữ liệu phân tán, thiếu hợp nhất:} Các hệ thống nghiệp vụ (Core Banking, CRM...) hoạt động hoàn toàn rời rạc. Việc thiếu vắng "nguồn sự thật duy nhất" buộc các phòng ban phải làm báo cáo thủ công, dẫn đến mâu thuẫn số liệu và sai lệch định dạng.
    
    \item \textbf{Không lưu vết lịch sử:} Các hệ thống OLTP liên tục ghi đè trạng thái mới lên dữ liệu cũ. Điều này khiến hệ thống mất khả năng truy xuất và phân tích trạng thái của khách hàng tại các thời điểm trong quá khứ.
    
    \item \textbf{Lãng phí dữ liệu phi cấu trúc:} Các chứng từ định danh quan trọng (CCCD, sổ tiết kiệm) chỉ được lưu trữ dưới dạng ảnh scan. Việc không thể truy vấn tự động khiến ngân hàng tốn kém nhiều nhân lực cho việc xác minh thủ công.
    
    \item \textbf{Thiếu tự động hóa và kiểm soát chất lượng:} Việc chạy báo cáo bằng tay gây chậm trễ và dễ sinh lỗi. Đồng thời, hệ thống thiếu cơ chế bắt lỗi tự động (bản ghi trùng, giá trị rỗng), khiến dữ liệu "bẩn" lọt vào báo cáo làm sai lệch quyết định kinh doanh.
    
    \item \textbf{Thiếu dấu vết kiểm toán:} Các luồng nạp dữ liệu không hề ghi nhận lại lịch sử thực thi (thời gian chạy, trạng thái thành công/thất bại). Hậu quả là kỹ sư dữ liệu không có cơ sở để truy vết nguyên nhân khi xảy ra sự cố.
\end{itemize}

\section{Khảo sát các giải pháp và kiến trúc hiện có}

Qua khảo sát thực tế tại thị trường Việt Nam, hệ thống dữ liệu của các ngân hàng hiện được chia thành 3 nhóm kiến trúc chính, phản ánh các giai đoạn chuyển đổi số khác nhau.

\subsection{Nhóm 1: Kiến trúc truyền thống (Legacy DWH + ETL)}

\textbf{Đại diện điển hình:} Agribank, BIDV, VietinBank, Vietcombank (ở các giai đoạn trước khi tiến hành chuyển đổi).

\begin{figure}[H]
    \centering
    \begin{tikzpicture}[
        node distance=0.8cm,
        >=Stealth,
        box/.style={rectangle, draw=blue!80!black, fill=blue!5, thick, minimum width=8cm, minimum height=0.8cm, align=center, rounded corners=3pt},
        arrow/.style={->, thick, draw=gray!70, line width=1.2pt}
    ]
        % Nodes
        \node[box] (src) {\textbf{Core Banking} \\ \footnotesize (T24 / Oracle FLEXCUBE / Silverlake)};
        \node[box, below=of src] (etl) {\textbf{ETL} \\ \footnotesize (Informatica / IBM DataStage / SSIS)};
        \node[box, below=of etl] (dwh) {\textbf{Traditional DWH} \\ \footnotesize (Oracle / SQL Server)};
        \node[box, below=of dwh] (bi) {\textbf{BI} \\ \footnotesize (IBM Cognos / SAP BusinessObjects / MicroStrategy)};

        % Arrows
        \draw[arrow] (src) -- (etl);
        \draw[arrow] (etl) -- (dwh);
        \draw[arrow] (dwh) -- (bi);
    \end{tikzpicture}
    \caption{Sơ đồ kiến trúc dữ liệu truyền thống}
\end{figure}

\begin{itemize}
    \item \textbf{Đặc điểm:} Sử dụng schema cứng, tốn nhiều chi phí bản quyền phần mềm. Hệ thống đáp ứng tốt các yêu cầu báo cáo định kỳ cho Ngân hàng Nhà nước (NHNN).
    \item \textbf{Hạn chế:} Chỉ cung cấp báo cáo T+1 (độ trễ 1 ngày). Không lưu trữ được lịch sử thay đổi của khách hàng, khó tích hợp đa kênh và thời gian xây dựng các báo cáo ad-hoc rất chậm (từ vài ngày đến hàng tuần).
\end{itemize}

\subsection{Nhóm 2: Đang chuyển đổi (Hybrid On-prem + Cloud)}

\textbf{Đại diện điển hình:} MB Bank, VPBank, HDBank, SeABank.

\begin{figure}[H]
    \centering
    \begin{tikzpicture}[
        node distance=0.8cm,
        >=Stealth,
        box/.style={rectangle, draw=teal!80!black, fill=teal!5, thick, minimum width=8cm, minimum height=0.8cm, align=center, rounded corners=3pt},
        arrow/.style={->, thick, draw=gray!70, line width=1.2pt}
    ]
        % Nodes
        \node[box] (src) {\textbf{Core Banking (On-premise)} \\ \footnotesize (Giữ nguyên để tuân thủ quy định NHNN)};
        \node[box, below=of src] (ingest) {\textbf{Ingestion} \\ \footnotesize (CDC / Kafka Streaming)};
        \node[box, below=of ingest] (dl) {\textbf{Data Lake} \\ \footnotesize (Azure Data Lake / AWS S3)};
        \node[box, below=of dl] (process) {\textbf{Processing} \\ \footnotesize (Spark / dbt)};
        \node[box, below=of process] (dwh) {\textbf{Cloud DWH} \\ \footnotesize (Azure Synapse / Amazon Redshift)};
        \node[box, below=of dwh] (bi) {\textbf{BI \& Analytics} \\ \footnotesize (Power BI / Tableau)};

        % Arrows
        \draw[arrow] (src) -- (ingest);
        \draw[arrow] (ingest) -- (dl);
        \draw[arrow] (dl) -- (process);
        \draw[arrow] (process) -- (dwh);
        \draw[arrow] (dwh) -- (bi);
    \end{tikzpicture}
    \caption{Sơ đồ kiến trúc lai (Hybrid) đang chuyển đổi}
\end{figure}

\begin{itemize}
    \item \textbf{Đặc điểm:} Hệ thống Core giữ tại local (on-premise) để tuân thủ luật lưu trữ dữ liệu trong nước, trong khi luồng phân tích được đẩy lên Cloud để tận dụng khả năng mở rộng. Đã bắt đầu ứng dụng Data Lake để lưu dữ liệu phi cấu trúc và manh nha áp dụng kiến trúc Medallion.
    \item \textbf{Hạn chế:} Tồn tại song song hai môi trường (on-prem và cloud) chưa thông suốt hoàn toàn. Chi phí vận hành tăng cao trong giai đoạn chuyển tiếp và kỹ năng nhân sự nội bộ chưa theo kịp tốc độ đổi mới công nghệ.
\end{itemize}

\subsection{Nhóm 3: Kiến trúc hiện đại (Modern Data Stack)}

\textbf{Đại diện điển hình:} Techcombank, TPBank, VIB.

\begin{figure}[H]
    \centering
    \begin{tikzpicture}[
        node distance=0.8cm,
        >=Stealth,
        box/.style={rectangle, draw=purple!80!black, fill=purple!5, thick, minimum width=8cm, minimum height=0.8cm, align=center, rounded corners=3pt},
        arrow/.style={->, thick, draw=gray!70, line width=1.2pt}
    ]
        % Nodes
        \node[box] (src) {\textbf{Đa nguồn (Multi-source)} \\ \footnotesize (Core Banking + CRM + Mobile App + Web + ATM)};
        \node[box, below=of src] (stream) {\textbf{Streaming \& Ingestion} \\ \footnotesize (Kafka / Pub-Sub)};
        \node[box, below=of stream] (dl) {\textbf{Cloud Data Lake} \\ \footnotesize (Google Cloud Storage / AWS S3)};
        \node[box, below=of dl] (transform) {\textbf{Orchestration \& Transform} \\ \footnotesize (dbt + Airflow / Cloud Composer)};
        \node[box, below=of transform] (dwh) {\textbf{Cloud DWH} \\ \footnotesize (BigQuery / Snowflake)};
        \node[box, below=of dwh] (serve) {\textbf{Consumption} \\ \footnotesize (Looker / Power BI / Tableau + ML Platform)};

        % Arrows
        \draw[arrow] (src) -- (stream);
        \draw[arrow] (stream) -- (dl);
        \draw[arrow] (dl) -- (transform);
        \draw[arrow] (transform) -- (dwh);
        \draw[arrow] (dwh) -- (serve);
    \end{tikzpicture}
    \caption{Sơ đồ ngăn xếp dữ liệu hiện đại (Modern Data Stack)}
\end{figure}

\begin{itemize}
    \item \textbf{Đặc điểm:} Đây là xu hướng đích đến của ngành. Hệ thống chuyển hoàn toàn từ ETL sang ELT (nạp dữ liệu thô vào kho trước rồi mới biến đổi).
    \item \textbf{Ưu điểm vượt trội:} Khả năng xử lý dữ liệu gần thời gian thực (near real-time) phục vụ phát hiện gian lận. Kiến trúc Medallion/Lakehouse được áp dụng bài bản có hệ thống, đồng thời tích hợp trực tiếp các luồng Machine Learning (ML pipeline) ngay trên nền tảng để chấm điểm tín dụng và phân khúc khách hàng.
\end{itemize}
\subsection{Ràng buộc pháp lý đặc thù tại Việt Nam}

Khác với xu hướng công nghệ của các ngân hàng trên thế giới, việc lựa chọn và thiết kế kiến trúc hệ thống dữ liệu tại Việt Nam chịu sự chi phối mạnh mẽ bởi các khung pháp lý và quy định đặc thù của cơ quan quản lý nhà nước. Cụ thể, kiến trúc dữ liệu bị ảnh hưởng trực tiếp bởi các quy định sau:

\begin{itemize}
    \item \textbf{Thông tư 09/2020/TT-NHNN (An toàn hệ thống thông tin):} Yêu cầu kiến trúc hệ thống phải hỗ trợ việc phân loại và thực hiện đánh giá rủi ro định kỳ.
    \item \textbf{Nghị định 13/2023/NĐ-CP (Bảo vệ dữ liệu cá nhân):} Đây là rào cản lớn nhất đối với điện toán đám mây, quy định dữ liệu khách hàng bắt buộc phải được lưu trữ trong nước, hạn chế việc ngân hàng đẩy dữ liệu lên các nền tảng Public Cloud của nước ngoài.
    \item \textbf{Thông tư 35/2016/TT-NHNN (An toàn bảo mật):} Đòi hỏi hệ thống nền tảng phải thiết lập cơ chế kiểm soát truy cập phân quyền chặt chẽ và có khả năng mã hóa toàn bộ dữ liệu nhạy cảm.
    \item \textbf{Yêu cầu báo cáo NHNN:} Bất kể kiến trúc nào được áp dụng, ngân hàng vẫn phải duy trì một luồng hệ thống có khả năng xuất và nộp báo cáo định kỳ theo đúng quy chuẩn do Ngân hàng Nhà nước ban hành.
\end{itemize}

\textbf{Hệ quả thực tế:} Từ những ràng buộc pháp lý khắt khe nêu trên, đa số các ngân hàng tại Việt Nam hiện nay đều ưu tiên lựa chọn triển khai mô hình đám mây riêng (Private Cloud) hoặc kiến trúc lai (Hybrid Cloud) nhằm kết hợp sự linh hoạt của điện toán đám mây với tính bảo mật của On-premise, thay vì chuyển dịch hoàn toàn lên Public Cloud như xu hướng của các ngân hàng phương Tây.

\section{Khảo sát các phương pháp xử lý ảnh hiện có}
Để trích xuất thông tin từ ảnh chứng từ như Căn cước công dân và Sổ tiết kiệm, hiện nay có ba hướng tiếp cận chính.

\begin{itemize}
    \item \textbf{Công cụ nhận dạng ký tự mã nguồn mở.} Tiêu biểu là Tesseract, nhận dạng dựa trên đặc trưng và mạng nơ-ron hồi tiếp. Ưu điểm là miễn phí và chạy hoàn toàn cục bộ, song độ chính xác giảm mạnh khi ảnh bị nghiêng, nhiễu hoặc nền phức tạp, và còn nhiều hạn chế với văn bản tiếng Việt có dấu, đòi hỏi khâu tiền xử lý ảnh kỹ lưỡng.

    \item \textbf{Công cụ nhận dạng dựa trên học sâu.} Các bộ công cụ như PaddleOCR hay EasyOCR kết hợp bước phát hiện vùng chứa chữ và bước nhận dạng nội dung. Nhóm này cho độ chính xác cao hơn với ảnh thực tế, hỗ trợ tốt tiếng Việt và có thể chạy trên máy tính phổ thông. Đây là hướng cân bằng giữa độ chính xác, chi phí và khả năng triển khai tại chỗ.

    \item \textbf{Dịch vụ nhận dạng trên nền tảng đám mây.} Các nền tảng như Google Cloud Vision, Google Gemini hay Groq cho độ chính xác rất cao nhờ được huấn luyện trên lượng dữ liệu lớn. Đổi lại, chúng phụ thuộc vào kết nối mạng, phát sinh chi phí theo lượt gọi và tiềm ẩn rủi ro bảo mật khi phải gửi dữ liệu định danh nhạy cảm ra bên ngoài.
\end{itemize}

Ngoài việc chọn công cụ nhận dạng, chất lượng kết quả còn phụ thuộc nhiều vào các kỹ thuật tiền xử lý ảnh, bao gồm căn chỉnh độ nghiêng và nắn lại phối cảnh để đưa chứng từ về dạng phẳng, chuyển đổi không gian màu và cân bằng độ sáng nhằm tăng độ tương phản, cùng việc khoanh vùng quan tâm để giới hạn phạm vi nhận dạng cho từng trường thông tin.

Trên cơ sở khảo sát này, đồ án lựa chọn PaddleOCR làm công cụ nhận dạng chính nhờ cân bằng tốt giữa độ chính xác, khả năng chạy cục bộ trên máy tính phổ thông và không phát sinh chi phí. Bên cạnh đó, ứng dụng web minh họa vẫn tích hợp thêm tùy chọn sử dụng dịch vụ đám mây để thuận tiện so sánh, đối chiếu kết quả.






\section{Yêu cầu đặt ra đối với hệ thống}
\subsection{Yêu cầu chức năng}
\begin{itemize}
    \item Thu thập và trích xuất tự động dữ liệu từ cơ sở dữ liệu và hình ảnh chứng từ.
    \item Tích hợp nạp dữ liệu gia tăng hàng ngày.
    \item Lưu trữ đầy đủ lịch sử thay đổi của các thực thể.
    \item Kiểm soát chất lượng dữ liệu và giám sát quy trình chuyển đổi, đảm bảo dữ liệu trong quá trình chuyển đổi không có bất kì sai sót nào ảnh hưởng đến quyết định Kinh doanh
\end{itemize}

\subsection{Yêu cầu phi chức năng}
\begin{itemize}
    \item \textbf{Tính lũy đẳng (Idempotency):} Chạy lại một luồng công việc nhiều lần không làm sai lệch hay nhân bản dữ liệu.
    \item \textbf{Chất lượng và Giám sát (DQ \& Audit):} Hệ thống phải tự động kiểm tra lỗi (Null, trùng lặp) và ghi log.
    \item \textbf{Tính tự động hóa:} Chạy theo lịch trình đã được định sẵn mà không cần can thiệp thủ công.
\end{itemize}

\chapter{Phân tích và thiết kế hệ thống}
\label{ch:analysis-design}

\section{Đề xuất kiến trúc giải pháp}
\label{sec:solution-arch}

Dựa trên những phân tích về hiện trạng và các vấn đề tồn đọng ở Chương 2, đồ án đề xuất giải pháp xây dựng hệ thống dữ liệu theo hướng \textbf{Modern Data Stack triển khai nội bộ (On-premise)} kết hợp với mô hình lưu trữ đa tầng và chiến lược nạp \textbf{ELT}. Đề xuất này dựa trên các cơ sở thiết kế sau:

\subsection{Lựa chọn kiến trúc tổng thể và chiến lược nạp}

\begin{itemize}
    \item \textbf{Kiến trúc Modern Data Stack On-premise:}
    \begin{itemize}
        \item \textit{Đáp ứng toàn diện yêu cầu:} Khắc phục hoàn toàn những giới hạn của kiến trúc Legacy và Hybrid, hệ thống mới đáp ứng đầy đủ các yêu cầu chức năng phức tạp như lưu vết lịch sử, tự động hóa quy trình, kiểm soát chất lượng dữ liệu và xử lý dữ liệu phi cấu trúc.
        \item \textit{Tối ưu chi phí:} Toàn bộ ngăn xếp (stack) sử dụng các công nghệ mã nguồn mở (Open-source), giúp hệ thống không phát sinh các khoản chi phí bản quyền đắt đỏ.
        \item \textit{Tuân thủ pháp lý:} Việc vận hành hoàn toàn trên hạ tầng nội bộ (on-premise) đảm bảo tuân thủ tuyệt đối quy định về lưu trữ dữ liệu cá nhân trong nước theo Nghị định 13/2023/NĐ-CP.
    \end{itemize}

    \item \textbf{Chiến lược nạp ELT (Extract $\rightarrow$ Load $\rightarrow$ Transform):}
    \begin{itemize}
        \item Hệ thống thực hiện trích xuất và nạp (Load) toàn bộ dữ liệu thô vào tầng Bronze trước, sau đó mới tiến hành biến đổi (Transform) bằng công cụ dbt ngay bên trong kho dữ liệu.
        \item Chiến lược này giúp phân tách rạch ròi trách nhiệm của các công cụ: Python chỉ chuyên trách Extract \& Load, trong khi dbt chỉ tập trung xử lý Transform. Điều này giúp hệ thống cực kỳ dễ bảo trì và gỡ lỗi (debug).
    \end{itemize}
\end{itemize}

\subsection{Lựa chọn mô hình tổ chức dữ liệu}

Hệ thống ứng dụng \textbf{kiến trúc Medallion (Bronze $\rightarrow$ Silver $\rightarrow$ Gold)} để phân tầng dữ liệu theo mức độ tinh sạch, mỗi tầng phục vụ một mục đích và nhóm người dùng riêng biệt. Trong đó, tầng \textbf{Bronze} đóng vai trò giữ lại dữ liệu thô nguyên bản, cho phép hệ thống dễ dàng tái xử lý lại thông tin khi logic nghiệp vụ thay đổi mà không cần phải trích xuất lại từ đầu. Các tầng tiếp theo được mô hình hóa chi tiết như sau:

\begin{itemize}
    \item \textbf{Mô hình hóa tầng Silver (Data Vault 2.0):}
    \begin{itemize}
        \item \textit{Lưu trữ lịch sử toàn diện:} Sử dụng bảng Satellite kết hợp kỹ thuật SCD Type 2 để ghi nhận mọi thay đổi, giúp hệ thống có khả năng trả lời chính xác câu hỏi "Dữ liệu tại thời điểm X trong quá khứ là gì?".
        \item \textit{Tính linh hoạt và mở rộng:} Cấu trúc tách biệt giữa khóa (Hub/Link) và thuộc tính (Satellite) giúp việc bổ sung thêm nguồn dữ liệu hoặc các trường thông tin mới không làm phá vỡ mô hình hiện hành.
        \item \textit{Tối ưu hiệu suất:} Thiết kế các bảng độc lập cho phép điều phối viên nạp dữ liệu song song thông qua cơ chế Pool, giúp giảm thiểu đáng kể tổng thời gian chạy luồng.
    \end{itemize}

    \item \textbf{Mô hình hóa tầng Gold (Star Schema):}
    \begin{itemize}
        \item \textit{Tối ưu truy vấn và báo cáo:} Với bảng sự kiện (Fact table) ở trung tâm và các bảng chiều (Dimension) bao quanh, Star Schema tạo ra cơ chế lọc chéo (cross-filter) tự nhiên trong Power BI. Thiết kế này hạn chế các phép kết nối (join) phức tạp, mang lại tốc độ truy vấn vượt trội.
        \item \textit{Phục vụ trực tiếp nghiệp vụ:} Cấu trúc mô hình sao được thiết kế để giải quyết trực tiếp chủ đề phân tích đã xác định ở Chương 2.
    \end{itemize}
\end{itemize}




\section{Kiến trúc tổng thể hệ thống}
\label{sec:overall-arch}
Sơ đồ kiến trúc tổng quan — hai pipeline độc lập, các thành phần như sau:
\begin{figure}[H]
    \centering
    \begin{tikzpicture}[
        >=Stealth,
        % Định nghĩa style cho các khối
        layer_box/.style={rectangle, draw=black, thick, minimum width=9cm, minimum height=1.2cm, align=center, outer sep=0pt},
        % Đã sửa lỗi LR Mode: Thay minimum width thành text width
        bronze_box/.style={rectangle, draw=black, thick, text width=8.5cm, align=left, inner sep=10pt, outer sep=0pt, fill=gray!5},
        arrow/.style={->, thick, draw=black},
        dashed_arrow/.style={->, thick, dashed, draw=black}
    ]

    % --- KHỐI SOURCE ---
    \node (source) at (0,0) {\textbf{SOURCE} (SQL Server + Ảnh tự sinh)};

    % --- KHỐI BRONZE LAYER ---
    \node[bronze_box, below=1.2cm of source] (bronze) {
        {\centering \textbf{BRONZE LAYER} \par} % Đã sửa cách căn giữa an toàn hơn
        \vspace{0.2em}
        \rule{\linewidth}{0.5pt} \\[0.8em]
        \textbf{Structured:}\\
        \hspace*{0.5cm} $\bullet$ \texttt{users\_tdy/pdy/mns}\\
        \hspace*{0.5cm} $\bullet$ \texttt{cards\_tdy/pdy/mns}\\
        \hspace*{0.5cm} $\bullet$ \texttt{transactions\_tdy/pdy/mns}\\
        \hspace*{0.5cm} $\bullet$ \texttt{mcc\_codes\_tdy/pdy/mns}\\[1em]
        \textbf{Unstructured:}\\
        \hspace*{0.5cm} $\bullet$ \texttt{id\_card\_results}\\
        \hspace*{0.5cm} $\bullet$ \texttt{savings\_book\_results}
    };

    % --- KHỐI SILVER, GOLD VÀ CONSUMPTION (Nằm sát nhau) ---
    \node[layer_box, below=2.5cm of bronze] (silver) {\textbf{SILVER LAYER} (Structured Only)\\Data Vault Model};
    \node[layer_box, anchor=north] (gold) at (silver.south) {\textbf{GOLD LAYER}\\Star Schema};
    \node[layer_box, anchor=north] (consumption) at (gold.south) {\textbf{CONSUMPTION}\\Power BI / Analytics};

    % --- VẼ CÁC MŨI TÊN ---
    % Mũi tên chia 3 nhánh từ Source xuống Bronze
    \draw[thick] (source.south) -- ++(0,-0.4) coordinate (fork);
    \draw[thick] ([xshift=-2.5cm]fork) -- ([xshift=2.5cm]fork);
    \draw[arrow] ([xshift=-2.5cm]fork) -- ([xshift=-2.5cm]fork |- bronze.north);
    \draw[arrow] (fork) -- (fork |- bronze.north);
    \draw[arrow] ([xshift=2.5cm]fork) -- ([xshift=2.5cm]fork |- bronze.north);

    % Mũi tên đứt nét từ Bronze xuống Silver
    \draw[dashed_arrow] ([xshift=-2cm]bronze.south) -- ([xshift=-2cm]silver.north);

    % Mũi tên đứt nét từ Bronze ra Direct Queries
    \draw[dashed_arrow] ([xshift=2cm]bronze.south) -- ++(0,-1.2) node[below] {\texttt{Direct Queries}};

    \end{tikzpicture}
    \caption{Kiến trúc luồng xử lý phân tầng của hệ thống}
    \label{fig:overall_architecture}
\end{figure}


% TODO: Luồng tổng quan: nguồn OLTP → Extract → Bronze → MNS → Silver
% (Data Vault) → Gold (Star Schema) → Power BI. Điều phối bằng Airflow.



\subsection{Pipeline dữ liệu phi cấu trúc (OCR)}
Do đặc thù ngành ngân hàng, các chứng từ như Căn cước công dân (CCCD) hay Sổ tiết kiệm chứa đựng lượng lớn dữ liệu cá nhân nhạy cảm. Việc sử dụng hình ảnh thực tế của khách hàng cho quá trình phát triển, thử nghiệm và trình diễn là vi phạm nghiêm trọng các quy định về bảo mật, đặc biệt là Nghị định 13/2023/NĐ-CP về bảo vệ dữ liệu cá nhân.

Để giải quyết triệt để rào cản pháp lý này, toàn bộ hình ảnh đầu vào được sử dụng để chạy thực nghiệm trong pipeline OCR đều là \textbf{dữ liệu giả lập}. Cụ thể, một tập lệnh Python tự động hóa đã được xây dựng để sinh ra các hình ảnh chứng từ dựa trên phôi mẫu chuẩn của nhà nước và ngân hàng. 

Quá trình giả lập này được thực hiện qua hai bước tinh vi:
\begin{itemize}
    \item \textit{Sinh nội dung ngẫu nhiên:} Điền các thông tin định danh hoàn toàn giả lập (họ tên, số ID, địa chỉ, số tiền...) vào các trường tương ứng trên phôi mẫu.
    \item \textit{Mô phỏng điều kiện thực tế:} Chủ động thêm vào các hiệu ứng nhiễu hình học và quang học bằng thuật toán (như làm mờ, xoay nghiêng, nhiễu hạt, chói sáng cục bộ hay đổ bóng).
\end{itemize}

\begin{figure}[H]
    \centering
    \includegraphics[width=0.7\textwidth]{img3/id_card_scan.jpg }
    \caption{Mẫu hình ảnh CCCD tự sinh}
    \label{fig:synthetic-cccd-sample}
\end{figure}


\begin{figure}[H]
    \centering
    \includegraphics[width=\textwidth]{img3/savings_book_scan.jpg }
    \caption{Mẫu hình ảnh sổ tiết kiệm tự sinh}
    \label{fig:synthetic-savings-book-sample}
\end{figure}

\begin{figure}[H]
\centering

%% ════════════════════════════════════════════════
%%  LUỒNG 1 — CCCD
%% ════════════════════════════════════════════════
\begin{subfigure}[t]{0.48\textwidth}
\centering
\begin{tikzpicture}[
  >=Stealth, node distance=0.38cm,
  iobox/.style={rectangle, rounded corners=2pt, draw=gray!60, fill=gray!15, text width=5.2cm, minimum height=0.62cm, align=center, font=\footnotesize},
  geom/.style={rectangle, rounded corners=2pt, draw=blue!50, fill=blue!10, text width=5.2cm, minimum height=0.62cm, align=center, font=\footnotesize},
  geom2/.style={rectangle, rounded corners=2pt, draw=blue!70, fill=blue!22, text width=5.2cm, minimum height=0.62cm, align=center, font=\footnotesize\itshape},
  enh/.style={rectangle, rounded corners=2pt, draw=teal!60, fill=teal!10, text width=5.2cm, minimum height=0.62cm, align=center, font=\footnotesize},
  enh2/.style={rectangle, rounded corners=2pt, draw=teal!70, fill=teal!22, text width=5.2cm, minimum height=0.62cm, align=center, font=\footnotesize\itshape},
  ocrbox/.style={rectangle, rounded corners=2pt, draw=orange!60, fill=orange!10, text width=5.2cm, minimum height=0.62cm, align=center, font=\footnotesize},
  ocrbox2/.style={rectangle, rounded corners=2pt, draw=orange!70, fill=orange!22, text width=5.2cm, minimum height=0.62cm, align=center, font=\footnotesize\itshape},
  postbox/.style={rectangle, rounded corners=2pt, draw=green!60, fill=green!10, text width=5.2cm, minimum height=0.62cm, align=center, font=\footnotesize},
  arr/.style={->, thick, draw=gray!60},
  lbl/.style={font=\scriptsize\bfseries, rotate=90, anchor=center},
]

\node[iobox] (inp) {\textbf{Ảnh CCCD}};

\node[geom,  below=0.5cm of inp] (g1) {Chỉnh nghiêng};
\node[geom,  below=of g1]        (g2) {Phát hiện biên thẻ};
\node[geom,  below=of g2]        (g3) {Căn phẳng phối cảnh};
\node[geom2, below=of g3]        (g4) {Chuẩn hóa chiều †};

\node[enh,  below=0.5cm of g4] (e1) {Chuyển ảnh xám};
\node[enh,  below=of e1]       (e2) {Khử nhiễu};
\node[enh,  below=of e2]       (e3) {Tăng tương phản cục bộ};
\node[enh,  below=of e3]       (e4) {Làm sắc nét};
\node[enh2, below=of e4]       (e5) {Ngưỡng hóa thích nghi †};

\node[ocrbox,  below=0.5cm of e5] (o1) {Nhận dạng ký tự toàn ảnh};
\node[ocrbox2, below=of o1]       (o2) {Trích xuất theo nhãn (ưu tiên) †};
\node[ocrbox,  below=of o2]       (o3) {Trích xuất theo vùng (dự phòng)};

\node[postbox, below=0.5cm of o3] (p1) {Chuẩn hóa \& phân tích cú pháp};
\node[postbox, below=of p1]       (p2) {Tính điểm tin cậy};

\node[iobox, below=of p2] (out) {\textbf{9 trường} + điểm tin cậy};

\foreach \a/\b in {
  inp/g1, g1/g2, g2/g3, g3/g4, g4/e1,
  e1/e2, e2/e3, e3/e4, e4/e5, e5/o1,
  o1/o2, o2/o3, o3/p1, p1/p2, p2/out}
{ \draw[arr] (\a)--(\b); }

% Vẽ nhãn nhóm CCCD
\path (g1.north west) ++(-0.2cm, 0) coordinate (g_top1);
\path (g4.south west) ++(-0.2cm, 0) coordinate (g_bot1);
\draw[decorate,decoration={brace,mirror,amplitude=4pt},draw=blue!50,thick]
  (g_top1) -- (g_bot1) node[midway,left=6pt,lbl,text=blue!60]{Hình học};

\path (e1.north west) ++(-0.2cm, 0) coordinate (e_top1);
\path (e5.south west) ++(-0.2cm, 0) coordinate (e_bot1);
\draw[decorate,decoration={brace,mirror,amplitude=4pt},draw=teal!60,thick]
  (e_top1) -- (e_bot1) node[midway,left=6pt,lbl,text=teal!60]{Tăng cường};

\path (o1.north west) ++(-0.2cm, 0) coordinate (o_top1);
\path (o3.south west) ++(-0.2cm, 0) coordinate (o_bot1);
\draw[decorate,decoration={brace,mirror,amplitude=4pt},draw=orange!60,thick]
  (o_top1) -- (o_bot1) node[midway,left=6pt,lbl,text=orange!60]{OCR};

\path (p1.north west) ++(-0.2cm, 0) coordinate (p_top1);
\path (p2.south west) ++(-0.2cm, 0) coordinate (p_bot1);
\draw[decorate,decoration={brace,mirror,amplitude=4pt},draw=green!60,thick]
  (p_top1) -- (p_bot1) node[midway,left=6pt,lbl,text=green!60]{Hậu xử lý};

\end{tikzpicture}
\caption{Căn cước công dân (CCCD)}
\end{subfigure}
\hfill
%% ════════════════════════════════════════════════
%%  LUỒNG 2 — SỔ TIẾT KIỆM
%% ════════════════════════════════════════════════
\begin{subfigure}[t]{0.48\textwidth}
\centering
\begin{tikzpicture}[
  >=Stealth, node distance=0.38cm,
  iobox/.style={rectangle, rounded corners=2pt, draw=gray!60, fill=gray!15, text width=5.2cm, minimum height=0.62cm, align=center, font=\footnotesize},
  geom/.style={rectangle, rounded corners=2pt, draw=blue!50, fill=blue!10, text width=5.2cm, minimum height=0.62cm, align=center, font=\footnotesize},
  enh/.style={rectangle, rounded corners=2pt, draw=teal!60, fill=teal!10, text width=5.2cm, minimum height=0.62cm, align=center, font=\footnotesize},
  ocrbox/.style={rectangle, rounded corners=2pt, draw=orange!60, fill=orange!10, text width=5.2cm, minimum height=0.62cm, align=center, font=\footnotesize},
  postbox/.style={rectangle, rounded corners=2pt, draw=green!60, fill=green!10, text width=5.2cm, minimum height=0.62cm, align=center, font=\footnotesize},
  arr/.style={->, thick, draw=gray!60},
  lbl/.style={font=\scriptsize\bfseries, rotate=90, anchor=center},
]

\node[iobox] (inp) {\textbf{Ảnh sổ tiết kiệm}};

\node[geom, below=0.5cm of inp] (g1) {Chỉnh nghiêng};
\node[geom, below=of g1]        (g2) {Phát hiện biên tài liệu};
\node[geom, below=of g2]        (g3) {Căn phẳng phối cảnh};

\node[enh, below=0.5cm of g3] (e1) {Chuyển ảnh xám};
\node[enh, below=of e1]       (e2) {Khử nhiễu};
\node[enh, below=of e2]       (e3) {Tăng tương phản cục bộ};
\node[enh, below=of e3]       (e4) {Làm sắc nét};

\node[ocrbox, below=0.5cm of e4] (o1) {Trích xuất theo 7 vùng định sẵn};
\node[ocrbox, below=of o1]       (o2) {Thử nhiều biến thể, chọn tốt nhất};

\node[postbox, below=0.5cm of o2] (p1) {Chuẩn hóa \& phân tích cú pháp};
\node[postbox, below=of p1]       (p2) {Tính điểm tin cậy};

\node[iobox, below=of p2] (out) {\textbf{7 trường} + điểm tin cậy};

\foreach \a/\b in {
  inp/g1, g1/g2, g2/g3, g3/e1,
  e1/e2, e2/e3, e3/e4, e4/o1,
  o1/o2, o2/p1, p1/p2, p2/out}
{ \draw[arr] (\a)--(\b); }

% Vẽ nhãn nhóm Sổ tiết kiệm
\path (g1.north west) ++(-0.2cm, 0) coordinate (g_top2);
\path (g3.south west) ++(-0.2cm, 0) coordinate (g_bot2);
\draw[decorate,decoration={brace,mirror,amplitude=4pt},draw=blue!50,thick]
  (g_top2) -- (g_bot2) node[midway,left=6pt,lbl,text=blue!60]{Hình học};

\path (e1.north west) ++(-0.2cm, 0) coordinate (e_top2);
\path (e4.south west) ++(-0.2cm, 0) coordinate (e_bot2);
\draw[decorate,decoration={brace,mirror,amplitude=4pt},draw=teal!60,thick]
  (e_top2) -- (e_bot2) node[midway,left=6pt,lbl,text=teal!60]{Tăng cường};

\path (o1.north west) ++(-0.2cm, 0) coordinate (o_top2);
\path (o2.south west) ++(-0.2cm, 0) coordinate (o_bot2);
\draw[decorate,decoration={brace,mirror,amplitude=4pt},draw=orange!60,thick]
  (o_top2) -- (o_bot2) node[midway,left=6pt,lbl,text=orange!60]{OCR};

\path (p1.north west) ++(-0.2cm, 0) coordinate (p_top2);
\path (p2.south west) ++(-0.2cm, 0) coordinate (p_bot2);
\draw[decorate,decoration={brace,mirror,amplitude=4pt},draw=green!60,thick]
  (p_top2) -- (p_bot2) node[midway,left=6pt,lbl,text=green!60]{Hậu xử lý};

\end{tikzpicture}
\caption{Sổ tiết kiệm}
\end{subfigure}

\vspace{0.5cm}
\raggedright
{\footnotesize † Bước đặc thù của pipeline CCCD, không áp dụng cho sổ tiết kiệm do kích thước và định dạng tài liệu khác biệt.}

\caption{So sánh luồng xử lý ảnh trong pipeline OCR phi cấu trúc}
\label{fig:ocr-both-pipelines}
\end{figure}

Về chi tiết hoạt động, hai luồng xử lý này có những đặc tả nghiệp vụ khác nhau:

\begin{itemize}
    \item \textbf{Pipeline xử lý ảnh Căn cước công dân (CCCD):}
    \begin{itemize}
        \item \textit{Tiền xử lý ảnh:} Hình ảnh chụp CCCD thường bị nhiễu do ánh sáng chói hoặc bóng râm. Hệ thống áp dụng thuật toán CLAHE để cân bằng độ tương phản, kết hợp Non-local Means để khử nhiễu hạt. Tiếp theo, thuật toán Canny Edge Detection và Perspective Transform được dùng để nắn phẳng tấm thẻ từ góc chụp nghiêng về dạng hình chữ nhật vuông vức.
        \item \textit{Trích xuất thông tin:} Mô hình nhận dạng quang học quét toàn bộ ảnh, định vị văn bản và kết hợp với các biểu thức chính quy (Regex) để trích xuất 9 trường dữ liệu cốt lõi (họ tên, ngày sinh, số CCCD, v.v.). Các kết quả có điểm tin cậy (Confidence Score) đạt yêu cầu sẽ được nạp trực tiếp vào bảng \texttt{id\_card\_results} tại tầng Bronze.
    \end{itemize}

    \item \textbf{Pipeline xử lý ảnh Sổ tiết kiệm:}
    \begin{itemize}
        \item \textit{Tiền xử lý ảnh:} Tương tự như CCCD, hệ thống làm sạch và chuẩn hóa hình ảnh sổ tiết kiệm. Tuy nhiên, luồng này bỏ qua các bước không cần thiết đối với giấy tờ dạng khổ lớn như chuẩn hóa chiều hay ngưỡng hóa thích nghi phức tạp, giúp giảm bớt tài nguyên và tăng tốc độ tính toán.
        \item \textit{Trích xuất thông tin:} Do định dạng sổ tiết kiệm đa dạng hơn, hệ thống ưu tiên bóc tách dựa vào 7 vùng định sẵn tương ứng với các trường thông tin (như số tài khoản, số tiền gửi, lãi suất, kỳ hạn). Hệ thống cũng tự động thử nghiệm nhiều biến thể trích xuất và giữ lại kết quả có điểm tin cậy cao nhất trước khi nạp vào bảng \texttt{savings\_book\_results}.
    \end{itemize}
\end{itemize}
\subsection{Pipeline dữ liệu có cấu trúc}
\begin{figure}[H]
    \centering
    \begin{tikzpicture}[
        >=Stealth,
        % ==========================================
        % Định nghĩa Style cho các khối
        % ==========================================
        layerbox/.style={rectangle, rounded corners=3pt, draw=gray!80, fill=gray!5, thick, text width=10.5cm, align=center, inner sep=8pt, font=\small},
        airbox/.style={rectangle, rounded corners=3pt, draw=orange!60, fill=orange!5, thick, text width=5.5cm, align=center, inner sep=6pt, font=\small},
        pgbox/.style={rectangle, rounded corners=3pt, draw=blue!50, fill=blue!5, thick, text width=3cm, align=center, inner sep=6pt, font=\small},
        auditbox/.style={rectangle, rounded corners=3pt, draw=purple!60, fill=purple!5, dashed, thick, text width=10.5cm, align=center, inner sep=8pt, font=\small},
        arr/.style={->, thick, draw=gray!80},
        lbl/.style={fill=white, font=\footnotesize\ttfamily, align=center, inner sep=2pt}
    ]

    % ==========================================
    % 1. MÔI TRƯỜNG DOCKER COMPOSE (HOST)
    % ==========================================
    \node[airbox] (airflow) at (2.5, 0) {
        \textbf{Apache Airflow} \\
        \texttt{banking\_structured\_dag @daily 02:00}
    };
    
    \node[pgbox] (pg) at (-3, 0) {
        \textbf{Postgres} \\
        (metadata DAG)
    };

    \draw[arr] (pg) -- (airflow);

    % Vẽ khung đứt nét bao quanh Docker Compose
    \path (pg.north west) ++(-0.4, 0.6) coordinate (D_TL);
    \path (airflow.south east) ++(0.4, -0.4) coordinate (D_BR);
    \draw[thick, dashed, draw=blue!70, rounded corners=4pt] (D_TL) rectangle (D_BR);
    \node[anchor=north west, font=\bfseries\small, text=blue!70] at (D_TL) {DOCKER COMPOSE (host)};

    % ==========================================
    % 2. MÔI TRƯỜNG SQL SERVER (HOST MACHINE)
    % ==========================================
    % Đã nới rộng tọa độ trục Y của các layer
    \node[layerbox] (oltp) at (0, -4.0) {
        \textbf{OLTP}
    };

    \node[layerbox] (bronze) at (0, -7.5) {
        \textbf{BRONZE — DATN.bronze} \\
        \texttt{*\_tdy} (hôm nay) $\xrightarrow{\text{Python}}$ \texttt{*\_pdy} $\cdot$ \texttt{*\_mns}
    };

    \node[layerbox] (silver) at (0, -11.0) {
        \textbf{SILVER — DATN.silver — Data Vault 2.0} \\
        Hub $\cdot$ Link $\cdot$ Satellite SCD2 (MD5 hash key)
    };

    \node[layerbox] (gold) at (0, -14.5) {
        \textbf{GOLD — DATN.gold — Star Schema} \\
        \texttt{dim\_customer $\cdot$ dim\_card $\cdot$ dim\_mcc $\cdot$ dim\_merchant} \\
        \texttt{dim\_date} $\cdot$ \texttt{fact\_transaction}
    };

    \node[auditbox] (audit) at (0, -17.2) {
        \texttt{audit.pipeline\_run\_log} $\xleftarrow{\text{ghi log}}$ \textbf{AuditedOperator} (mọi task)
    };

    % Vẽ khung liền nét bao quanh SQL Server
    \path (oltp.north west) ++(-0.4, 0.7) coordinate (S_TL);
    \path (audit.south east) ++(0.4, -0.4) coordinate (S_BR);
    \draw[thick, draw=purple!70, rounded corners=4pt] (S_TL) rectangle (S_BR);
    \node[anchor=north west, font=\bfseries\small, text=purple!70] at (S_TL) {SQL Server (host machine)};

    % ==========================================
    % 3. CÁC MŨI TÊN KẾT NỐI (ARROWS) VÀ ĐIỀU PHỐI
    % ==========================================
    
    % Airflow trigger xuống SQL Server (Mũi tên đứt nét trỏ thẳng vào OLTP)
    \draw[arr, dashed] (airflow.south) -- node[lbl] {trigger tasks} (airflow.south |- oltp.north);

    % Các mũi tên luân chuyển và biến đổi dữ liệu giữa các tầng (ETL/ELT)
    \draw[arr] (oltp.south) -- node[lbl] {Python + SQLAlchemy \\ \texttt{load\_bronze\_<entity>.py}} (bronze.north);
    
    \draw[arr] (bronze.south) -- node[lbl] {dbt \\ \texttt{tag:hub $\rightarrow$ tag:link $\rightarrow$ tag:satellite}} (silver.north);
    
    \draw[arr] (silver.south) -- node[lbl] {dbt \\ \texttt{tag:dim $\rightarrow$ fact\_transaction}} (gold.north);

    \end{tikzpicture}
    \caption{Kiến trúc tổng quan của hệ thống Pipeline xử lý dữ liệu có cấu trúc}
    \label{fig:structured_pipeline_architecture}
\end{figure}


\section{Thiết kế mô hình dữ liệu cho pipeline có cấu trúc}


\subsection{Nguồn OLTP}

\begin{figure}[H]
    \centering
    \includegraphics[width=\textwidth]{img3/OLTP_removed.png}
    \caption{Kiến trúc OLTP của hệ thống}
    \label{fig:bronze-layer-architecture}
\end{figure}

\subsection{Tầng Bronze}
\begin{itemize}
    \item \textbf{Mô hình kiến trúc phân tầng (Medallion Architecture):} Hệ thống được thiết kế theo mô hình đa tầng từ nguồn đến đích nhằm chuẩn hóa và kiểm soát chất lượng dữ liệu:
    \begin{itemize}
        \item \textit{Nguồn (Source):} Hệ thống cơ sở dữ liệu gốc từ SQL Server.
        \item \textit{Tầng Bronze (TDY/PDY/MNS):} Lưu trữ dữ liệu thô với cơ chế phân tách bảng.
        \item \textit{Tầng Silver (Data Vault 2.0):} Mô hình hóa dữ liệu cốt lõi, bao gồm các \texttt{Hubs} (khóa nghiệp vụ), \texttt{Links} (mối quan hệ) và \texttt{Satellites} (lưu trữ thuộc tính và lịch sử thay đổi).
        \item \textit{Tầng Gold (Star Schema):} Phục vụ phân tích với bảng sự kiện  và các bảng chiều xung quanh .
    \end{itemize}

\item \textbf{Luồng dữ liệu có cấu trúc:}
    \begin{itemize}
        \item \textit{Nguồn dữ liệu:} Trích xuất từ SQL Server thuộc schema \texttt{banking}
        \item \textit{Chiến lược triển khai tại tầng Bronze:} Mỗi thực thể được tổ chức thành một bộ 3 bảng độc lập để phục vụ cơ chế nạp gia tăng (Incremental Load) và kiểm soát thay đổi (Change Data Capture):
        \begin{itemize}
            \item \textbf{Thực thể Users:} Gồm \texttt{bronze.users\_tdy} (Snapshot hôm nay), \texttt{bronze.users\_pdy} (Snapshot lũy kế), và \texttt{bronze.users\_mns} (Change-set I/U).
            \item \textbf{Thực thể Cards:} Gồm \texttt{bronze.cards\_tdy} (Snapshot hôm nay), \texttt{bronze.cards\_pdy} (Snapshot lũy kế), và \texttt{bronze.cards\_mns} (Change-set I/U).
            \item \textbf{Thực thể Transactions:} Gồm \texttt{bronze.transactions\_tdy} (Giao dịch hôm nay), \texttt{bronze.transactions\_pdy} (Giao dịch ngày trước), và \texttt{bronze.transactions\_mns} (Change-set chỉ dành cho Insert).
            \item \textbf{Thực thể MCC Codes:} Gồm \texttt{bronze.mcc\_codes\allowbreak\_tdy} (Snapshot hôm nay), \texttt{bronze.mcc\_codes\allowbreak\_pdy} (Snapshot lũy kế), và  \newline  \texttt{bronze.mcc\_codes\allowbreak\_mns} (Change-set I/U).
        \end{itemize}
    \end{itemize}

    \item \textbf{Luồng dữ liệu phi cấu trúc:}
    \begin{itemize}
        \item \textit{Nguồn dữ liệu:} Hình ảnh vật lý của Căn cước công dân (CCCD) và Sổ tiết kiệm. Dữ liệu này được đưa qua pipeline xử lý nhận dạng quang học sử dụng AI (PaddleOCR).
        \item \textit{Chiến lược triển khai tại tầng Bronze:} Kết quả văn bản sau khi được OCR trích xuất sẽ được lưu trữ trực tiếp vào các bảng đích tương ứng:
        \begin{itemize}
            \item \texttt{bronze.id\_card\_results}: Lưu trữ toàn bộ kết quả OCR bóc tách từ ảnh Căn cước công dân.
            \item \texttt{bronze.savings\_book\_results}: Lưu trữ toàn bộ kết quả OCR bóc tách từ ảnh Sổ tiết kiệm.
        \end{itemize}
    \end{itemize}
\end{itemize}






\begin{figure}[H]
    \centering
    \includegraphics[width=\textwidth]{img3/bronze.png}
    \caption{Kiến trúc dữ liệu tầng Bronze}
    \label{fig:silver-layer-architecture}
\end{figure}




\subsection{Tầng Silver: Data Vault 2.0}
Tầng Silver được thiết kế theo mô hình kiến trúc Data Vault để lưu lịch sử thay đổi dữ liệu và dễ dàng mở rộng như hình
\begin{figure}[H]
    \centering
    \includegraphics[width=\textwidth]{img3/Silver_Layer.png}
    \caption{Kiến trúc dữ liệu tầng Silver}
    \label{fig:gold-layer-architecture}
\end{figure}

\subsection{Tầng Gold: Star Schema}
Tầng Gold được thiết kế theo Start Schema với các bảng Dim và Fact 
\begin{figure}[H]
    \centering
    \includegraphics[width=\textwidth]{img3/Gold_Layer.png}
    \caption{Kiến trúc dữ liệu tầng Gold}
    \label{fig:oltp-source-model}
\end{figure}

\subsection{Schema Audit}

Bên cạnh các schema phục vụ lưu trữ dữ liệu nghiệp vụ, hệ thống còn thiết kế riêng một schema \texttt{audit} để phục vụ công tác quản trị và vận hành. 

\textbf{Mục đích của Schema Audit:}
\begin{itemize}
    \item Ghi lại toàn bộ lịch sử thực thi của mọi task trong pipeline (bao gồm thông tin luồng chạy, thời gian thực thi và kết quả đầu ra).
    \item Hỗ trợ đắc lực cho kỹ sư dữ liệu trong việc theo dõi, gỡ lỗi và đánh giá độ tin cậy cũng như hiệu năng hoạt động của hệ thống.
\end{itemize}

\textbf{Cấu trúc bảng chính: \texttt{audit.pipeline\_run\_\allowbreak log}}

Đây là bảng nhật ký trung tâm, được thiết kế để lưu vết mọi hoạt động của luồng dữ liệu. Cấu trúc bảng được chia thành các nhóm cột chi tiết như sau:

\begin{itemize}
    \item \textbf{Nhóm định danh task:} Gồm các cột \texttt{dag\_id}, \texttt{task\_id}, \texttt{run\_date}, \texttt{attempt}. Nhóm này giúp định danh chính xác tác vụ đang chạy, thuộc luồng nào và ở lần chạy thứ mấy.
    \item \textbf{Nhóm trạng thái:} Cột \texttt{status} ghi nhận tình trạng hiện tại của task với các giá trị chuẩn như: \texttt{started}, \texttt{success}, \texttt{failed}, \texttt{skipped}.
    \item \textbf{Nhóm thời gian:} Gồm \texttt{started\_at}, \texttt{ended\_at} và \texttt{duration\_sec}. Trong đó, \texttt{duration\_sec} là cột tính toán tự động giúp đo lường thời gian chạy thực tế của task.
    \item \textbf{Nhóm thống kê tác động dữ liệu:} Gồm \texttt{rows\_processed}, \texttt{rows\_inserted}, \texttt{rows\_updated}, \texttt{rows\_deleted}. Nhóm cột này cung cấp số liệu thực tế về khối lượng bản ghi được xử lý qua từng bước.
    \item \textbf{Nhóm thông tin lỗi \& debug:} Bao gồm \texttt{error\_message}, \texttt{log\_file\_path}, \texttt{airflow\_log\_url}, \texttt{host\_name} và đặc biệt là \texttt{extra\_metadata} (lưu dưới dạng định dạng JSON tự do cho phép mở rộng linh hoạt các thông tin log đặc thù mà không cần sửa bảng).
\end{itemize}

Để tối ưu hóa việc truy xuất và trực quan hóa tình trạng "sức khỏe" của toàn bộ hệ thống pipeline, schema \texttt{audit} được thiết kế thêm các View tổng hợp dữ liệu từ bảng log chính.

\textbf{Các View hỗ trợ phân tích:}
\begin{itemize}
    \item \texttt{audit.v\_latest\_run\_\allowbreak per\_task}: Cung cấp trạng thái thực thi mới nhất của từng task. View này đóng vai trò như một bảng điều khiển thời gian thực, giúp kỹ sư dữ liệu nắm bắt nhanh chóng hiện trạng của hệ thống (hiện có task nào đang treo hay vừa thất bại không).
    \item \texttt{audit.v\_task\_success\_\allowbreak rate\_30d}: Chuyên dùng để tính toán tỷ lệ thành công  và thống kê thời gian chạy (trung bình, tối đa) của các task trong vòng 30 ngày gần nhất. Tập dữ liệu từ view này đặc biệt quan trọng trong việc phát hiện các nút thắt cổ chai về hiệu năng và đánh giá độ ổn định dài hạn của các luồng xử lý.
\end{itemize}

\begin{figure}[H]
    \centering
    % Điều chỉnh width về tỷ lệ 0.9 hoặc \textwidth tùy thuộc độ phân giải ảnh gốc của bạn
    \includegraphics[width=0.9\textwidth]{img3/audit.png} 
    \caption{Sơ đồ cấu trúc dữ liệu của Schema Audit}
    \label{fig:email-alert-example}
\end{figure}

\section{Thiết kế cơ chế điều phối}
% TODO: Thiết kế luồng tác vụ điều phối trên Apache Airflow:
\subsection{Vận hành và điều phối}

Để vận hành toàn bộ nền tảng dữ liệu được thiết kế trong Đồ án tốt nghiệp, hệ thống sử dụng Apache Airflow làm công cụ điều phối trung tâm. Toàn bộ các tiến trình từ trích xuất, biến đổi dữ liệu, cho đến kiểm định chất lượng và ghi nhận nhật ký đều được tự động hóa và đóng gói hoàn chỉnh. Toàn bộ vòng đời của hệ thống dữ liệu trong đồ án được vận hành thông qua 3 luồng DAG cốt lõi:

\begin{itemize}
    \item \textbf{\texttt{banking\_structured\_dag}}: Là luồng xử lý xương sống, chịu trách nhiệm chuyển hóa dữ liệu có cấu trúc từ hệ thống nguồn OLTP đi qua lần lượt các tầng Bronze, Silver và Gold theo chuẩn kiến trúc Medallion.
    \begin{figure}[H]
    \centering
    % Điều chỉnh width về tỷ lệ 0.9 hoặc \textwidth tùy thuộc độ phân giải ảnh gốc của bạn
    \includegraphics[width=0.9\textwidth]{img3/banking_structured..png} 
    \caption{Sơ đồ workflow của \texttt{banking\_structured\_dag}}
    \label{fig:banking-structured-dag-workflow}
\end{figure}







    \item \textbf{\texttt{data\_quality\_dag}}: Đóng vai trò là chốt chặn kiểm duyệt, tự động kích hoạt để đánh giá các chỉ tiêu chất lượng dữ liệu ngay sau khi luồng dữ liệu có cấu trúc cập nhật xong, đảm bảo dữ liệu tại tầng Gold luôn đạt chuẩn trước khi đưa lên báo cáo.
    \begin{figure}[H]
    \centering
    % Điều chỉnh width về tỷ lệ 0.9 hoặc \textwidth tùy thuộc độ phân giải ảnh gốc của bạn
    \includegraphics[width=0.9\textwidth]{img3/Data_quality.png} 
    \caption{Sơ đồ workflow của \texttt{data\_quality\_dag}}
    \label{fig:data-quality-dag-workflow}
\end{figure}
    \item \textbf{\texttt{ocr\_unstructured\_dag}}: Vận hành luồng xử lý dữ liệu phi cấu trúc, tự động hóa quy trình nhận dạng quang học để bóc tách thông tin từ các hình ảnh chứng từ vật lý (Căn cước công dân, Sổ tiết kiệm) và nạp vào cơ sở dữ liệu.
\end{itemize}
    \begin{figure}[H]
    \centering
    % Điều chỉnh width về tỷ lệ 0.9 hoặc \textwidth tùy thuộc độ phân giải ảnh gốc của bạn
    \includegraphics[width=0.9\textwidth]{img3/ocr.png} 
    \caption{Sơ đồ workflow của \texttt{ocr\_unstructured\_dag}}
    \label{fig:ocr-unstructured-dag-workflow}
\end{figure}

Lịch trình thực thi, tần suất chạy cũng như mối liên hệ phụ thuộc giữa 3 luồng DAG này được thiết kế một cách logic nhằm tối ưu hóa tài nguyên hệ thống, được mô tả chi tiết trong Hình \ref{fig:airflow_dag_timeline}.

\begin{figure}[H]
    \centering
    \begin{tikzpicture}[
        >=Stealth,
        % ==========================================
        % Định nghĩa Style cho các đối tượng
        % ==========================================
        ax/.style={thick, draw=gray!80},
        tick/.style={thick, draw=gray!80},
        time/.style={font=\bfseries, text=black},
        dagbox/.style={rectangle, rounded corners=3pt, draw=blue!60, fill=blue!5, thick, text width=4.5cm, align=center, inner sep=6pt, font=\small},
        dqbox/.style={rectangle, rounded corners=3pt, draw=teal!60, fill=teal!5, thick, text width=4.5cm, align=center, inner sep=6pt, font=\small},
        indepbox/.style={rectangle, rounded corners=3pt, draw=orange!60, fill=orange!5, thick, text width=6.5cm, align=center, inner sep=6pt, font=\small},
        arr/.style={->, thick, draw=gray!80}
    ]

    % ==========================================
    % PHẦN 1: TIMELINE HẰNG NGÀY (CÓ PHỤ THUỘC)
    % ==========================================
    \node[anchor=west, font=\bfseries] at (-0.5, 1) {TIMELINE HẰNG NGÀY};
    
    % Vẽ trục thời gian ngang
    \draw[ax] (-0.5, 0) -- (11.5, 0);

    % --- Mốc 02:00 ---
    \node[time, above] at (1.5, 0.2) {02:00};
    \draw[tick] (1.5, 0.2) -- (1.5, 0);
    \draw[arr] (1.5, 0) -- (1.5, -0.8);
    
    \node[dagbox, anchor=north] (dag1) at (1.5, -0.8) {
        \textbf{\texttt{banking\_structured\_dag}} \\
        \footnotesize (tự động, \texttt{@daily})
    };

    % --- Mốc 04:00 ---
    \node[time, above] at (10, 0.2) {04:00};
    \draw[tick] (10, 0.2) -- (10, 0);
    \draw[arr] (10, 0) -- (10, -0.8);
    
    \node[dqbox, anchor=north] (dag2) at (10, -0.8) {
        \textbf{\texttt{data\_quality\_dag}} \\
        \footnotesize (tự động, \texttt{@daily}) \\
        \footnotesize chờ structured DAG \\
        \footnotesize thành công trong 24h
    };

    % --- Mốc ~03:30 (Tính toán vị trí chính giữa tự động) ---
    % Tìm điểm giữa của DAG 1 và DAG 2 để đặt mốc 3:30
    \path (dag1.east) -- (dag2.west |- dag1.east) coordinate[midway] (mid);
    
    \node[time, above] at (mid |- 0, 0.2) {$\sim$03:30};
    \draw[tick] (mid |- 0, 0.2) -- (mid |- 0, 0);
    
    % Mũi tên ngang nối từ DAG 1 sang DAG 2, ghi chữ "Kết thúc" ở giữa
    \draw[arr] (dag1.east) -- node[fill=white, font=\small\itshape, inner sep=4pt] (end1) {Kết thúc} (dag2.west |- dag1.east);
    
    % Mũi tên từ trục thời gian ~03:30 trỏ xuống chữ "Kết thúc"
    \draw[arr] (mid |- 0, 0) -- (end1.north);


    % ==========================================
    % PHẦN 2: TIMELINE THỦ CÔNG (ĐỘC LẬP)
    % ==========================================
    \node[anchor=west, font=\bfseries] at (-0.5, -4.5) {BẤT KỲ LÚC NÀO (Trigger thủ công)};
    
    % Vẽ một đoạn trục ngắn tượng trưng
    \draw[ax] (-0.5, -5.2) -- (4, -5.2);
    
    \draw[arr] (1.5, -5.2) -- (1.5, -6.0);
    
    \node[indepbox, anchor=north] (dag3) at (1.5, -6.0) {
        \textbf{\texttt{ocr\_unstructured\_dag}} \\
        \footnotesize (độc lập, không phụ thuộc 2 DAG trên)
    };

    \end{tikzpicture}
    \caption{Lịch trình điều phối các DAG trong Apache Airflow}
    \label{fig:airflow_dag_timeline}
\end{figure}
\subsection{Quản lý và giới hạn tài nguyên}

Trong quá trình điều phối, việc chạy đồng thời quá nhiều tác vụ (tasks) có thể dẫn đến cạn kiệt tài nguyên máy chủ (RAM, CPU) hoặc gây quá tải cho các hệ thống nguồn (như cơ sở dữ liệu OLTP đang phục vụ giao dịch cốt lõi). Để giải quyết vấn đề này, hệ thống sử dụng cơ chế \textbf{Pools} của Apache Airflow nhằm kiểm soát chặt chẽ mức độ đồng thời của các nhóm tác vụ đặc thù. 

Cấu hình chi tiết của 3 pool được triển khai trong hệ thống được trình bày trong Bảng \ref{tab:airflow_pools}:

\begin{table}[H]
    \centering
    \renewcommand{\arraystretch}{1.4}
    \begin{tabular}{|l|c|p{8.5cm}|}
        \hline
        
        \textbf{Tên Pool} & \textbf{Concurrency} & \textbf{Mục đích} \\
        \hline
        \texttt{source\_db\_pool} & 2 & Giới hạn số lượng kết nối trích xuất đồng thời vào cơ sở dữ liệu nguồn (OLTP DB). Tránh việc tạo ra quá nhiều session truy vấn cùng lúc làm ảnh hưởng đến hiệu năng giao dịch thực tế của ngân hàng. \\
        \hline
        \texttt{dbt\_pool} & 4 & Giới hạn số lượng các mô hình biến đổi dữ liệu (dbt models) được phép thực thi song song trong quá trình xây dựng tầng Silver và Gold. \\
        \hline
        \texttt{ocr\_pool} & 1 & Tác vụ nhận dạng quang học bằng AI (PaddleOCR) yêu cầu khối lượng tính toán cực kỳ lớn (CPU-bound nặng). Việc giới hạn mức 1 đảm bảo hệ thống chỉ xử lý tuần tự từng ảnh, tránh hiện tượng tràn bộ nhớ và treo toàn bộ pipeline. \\
        \hline
    \end{tabular}
    \caption{Cấu hình giới hạn tài nguyên trong Apache Airflow}
    \label{tab:airflow_pools}
\end{table}

\subsection{Cơ chế cảnh báo tự động}

Để đảm bảo tính chủ động trong công tác vận hành và rút ngắn thời gian xử lý sự cố, hệ thống điều phối Apache Airflow được tích hợp cơ chế thông báo tự động thông qua các hàm callback. Cơ chế này kết nối trực tiếp với các kênh giao tiếp của Email và hoạt động dựa trên 3 kịch bản chính:

\begin{itemize}
    \item \textbf{Cảnh báo thất bại} Bất kỳ tác vụ nào gặp lỗi và thất bại trong quá trình thực thi sẽ ngay lập tức kích hoạt callback để bắn cảnh báo khẩn cấp qua \textbf{Email}. Thông báo này thường đính kèm các ngữ cảnh quan trọng như tên luồng, tên tác vụ bị lỗi và đường dẫn (URL) trực tiếp tới file log để kỹ sư có thể truy cập và gỡ lỗi ngay lập tức.
    
    \item \textbf{Thông báo thử lại} Đối với các lỗi có thể do nguyên nhân khách quan tạm thời (như gián đoạn mạng, timeout), Airflow được cấu hình để thử lại. Khi quá trình này diễn ra, hệ thống sẽ báo cho đội ngũ biết rằng một tác vụ đang tự động khắc phục lỗi, tránh việc báo động giả làm trôi tin nhắn.
    
    \item \textbf{Xác nhận thành công:} Khi điểm neo cuối cùng của DAG hoàn thành, tức là toàn bộ pipeline đã chạy trọn vẹn và dữ liệu đã an toàn cập bến tầng Gold, một thông báo sẽ được gửi đi. Đây là tín hiệu xác nhận kho dữ liệu đã sẵn sàng để phục vụ cho các dashboard và báo cáo phân tích trong ngày.
\end{itemize}

\begin{figure}[H]
    \centering

    \includegraphics[width=0.9\textwidth]{img3/mail.png} 
    \caption{Hình ảnh cảnh báo được gửi về email cá nhân}
    \label{fig:audit-schema-diagram}
\end{figure}

\chapter{Xây dựng chương trình}
\label{ch:implementation}

\section{Công nghệ sử dụng}
\label{sec:tech}

\subsection{Hệ quản trị Cơ sở dữ liệu và Kho dữ liệu}


Trong kiến trúc của nền tảng, Microsoft SQL Server được lựa chọn làm "trái tim" lưu trữ và xử lý dữ liệu của toàn bộ hệ thống. Quyết định sử dụng công nghệ này dựa trên các yếu tố kỹ thuật và kiến trúc sau:

\begin{itemize}
    \item \textbf{Đóng vai trò kép:} SQL Server đảm nhận đồng thời hai nhiệm vụ quan trọng. Thứ nhất, nó đóng vai trò là nguồn dữ liệu tác nghiệp (OLTP) lưu trữ  Thứ hai, nó đóng vai trò là Kho dữ liệu đích (Data Warehouse) với tên database là \texttt{DATN}. Kho dữ liệu phân tích này được tổ chức quy củ thành 4 schema độc lập bao gồm: \texttt{bronze}, \texttt{silver}, \texttt{gold}, và \texttt{audit}.
    
    \item \textbf{Tương thích hệ sinh thái \& Kiến trúc Data Vault:} SQL Server tích hợp vô cùng mượt mà với công cụ \textbf{dbt} (Data Build Tool). Đây là nền tảng, bắt buộc phải có để xây dựng thành công kiến trúc Data Vault 2.0 tại tầng Silver. Ngoài ra, việc sử dụng nền tảng của Microsoft hoàn toàn bám sát với xu hướng công nghệ thực tế tại các Ngân hàng và tổ chức Tài chính lớn.
    
    \item \textbf{Mô hình triển khai mạng tối ưu:} Nhằm tối ưu hóa hiệu năng I/O cho cơ sở dữ liệu thay vì bọc qua một lớp ảo hóa, SQL Server được cài đặt chạy trực tiếp trên máy chủ vật lý (Host Windows). Các thành phần điều phối chạy trong Docker Container (như Airflow) sẽ kết nối ngược ra môi trường host sử dụng \texttt {driver ODBC 17}.
\end{itemize}

\begin{figure}[H]
    \centering

    \includegraphics[width=0.7\textwidth]{img4/sql server.png} 
    \caption{Hệ quản trị cơ sở dữ liệu SQL Server}
    \label{fig:sql-server-tool}
\end{figure}


\subsection{Công cụ điều phối và biến đổi dữ liệu}
% TODO: Trình bày các công nghệ cốt lõi để chạy Data Pipeline:
\begin{itemize}
    \item \textbf{Apache Airflow (Điều phối):} Lập lịch, quản lý phụ thuộc và giám sát trạng thái thực thi của các luồng công việc.
    \item \textbf{dbt - data build tool (Biến đổi):} Viết các model bằng SQL để thực hiện phép biến đổi dữ liệu từ tầng Bronze sang Silver và Gold, đồng thời quản lý tài liệu và kiểm thử (dbt tests).
\end{itemize}
\begin{figure}[H]
    \centering

    \includegraphics[width=0.7\textwidth]{img4/Gemini_Generated_Image_364pk0364pk0364p.png} 
    \caption{Công cụ điều phối và biến đổi dữ liệu}
    \label{fig:airflow-dbt-tools}
\end{figure}


\subsection{Xử lý ảnh và trích xuất tài liệu}
% TODO: Giới thiệu thư viện PaddleOCR:
\begin{itemize}
    \item Ứng dụng mô hình AI của PaddleOCR để nhận diện văn bản từ các ảnh chứng từ.
    \item Tiền xử lý ảnh (OpenCV) và trích xuất thông tin (CCCD, sổ tiết kiệm) chuyển đổi thành định dạng có cấu trúc.
\end{itemize}
\begin{figure}[H]
    \centering

    \includegraphics[width=0.7\textwidth]{img4/opencv.png} 
    \caption{Công cụ xử lý ảnh và trích xuất dữ liệu}
    \label{fig:opencv-ocr-tool}
\end{figure}

\subsection{Đóng gói và quản lý môi trường}
% TODO: Cách thức thiết lập môi trường chạy ổn định:
\begin{itemize}
    \item \textbf{Docker} Đóng gói các dịch vụ (Airflow, SQL Server) vào container để dễ dàng triển khai.
    \item \textbf{Quản lý môi trường Python:} Phân tách các thư viện xung đột bằng hai môi trường ảo độc lập: \texttt{.venv} (cho Airflow/dbt) và \texttt{.venv\_ocr} (cho PaddleOCR).
\begin{figure}[H]
    \centering

    \includegraphics[width=0.5\textwidth]{img4/dockerize-windows-app.png} 
    \caption{Công cụ đóng gói và quản lý môi trường}
    \label{fig:docker-environment-tool}
\end{figure}

\end{itemize}



\subsection{Trực quan hóa dữ liệu (BI)}
Xây dựng Báo cáo Phân tích Giao dịch. Kết nối nền tảng Power BI trực tiếp với tầng dữ liệu Gold. Hệ thống khai thác các bảng chiều và bảng fact để trực quan hóa lịch sử giao dịch và phân tích hành vi khách hàng.
\begin{figure}[H]
    \centering

    \includegraphics[width=0.5\textwidth]{img4/power-bi.png} 
    \caption{Công cụ trực quan hóa dữ liệu}
    \label{fig:power-bi-tool}
\end{figure}

\section{Quy trình xây dựng Pipeline dữ liệu có cấu trúc}


\subsection{Trích xuất và nạp dữ liệu vào tầng Bronze}

Tầng Bronze có nhiệm vụ tiếp nhận dữ liệu thô từ hệ thống OLTP và lưu
trữ nguyên trạng mà không thực hiện bất kỳ phép biến đổi nào. Toàn bộ
giá trị được giữ dưới dạng chuỗi như nguồn gốc, việc ép kiểu chỉ diễn
ra ở tầng Silver.

\subsubsection*{a) Cấu trúc ba bảng theo mẫu MNS}

Mỗi entity dữ liệu (users, cards, transactions, mcc\_codes) được tổ
chức thành ba bảng với vai trò tách biệt:

\begin{itemize}
    \item \textbf{*\_tdy (Today):} Lưu batch dữ liệu của ngày chạy hiện
          tại, được truncate và nạp lại từ đầu trong mỗi lần thực thi.
    \item \textbf{*\_pdy (Previous Day):} Snapshot lũy kế, lưu trạng thái
          gần nhất của mọi bản ghi từng xuất hiện qua các lần chạy trước.
    \item \textbf{*\_mns (Minus):} Bảng change-set chứa cặp
          \texttt{(id, operation\_flag)}, là đầu ra của bước so sánh và
          là đầu vào duy nhất mà tầng Silver đọc để xử lý.
\end{itemize}

\subsubsection*{b) Thuật toán tính chênh lệch gia tăng}

Sau khi batch mới được nạp vào \texttt{*\_tdy}, hệ thống thực hiện so
sánh với \texttt{*\_pdy} để xác định sự thay đổi và ghi kết quả vào
\texttt{*\_mns} với hai loại cờ:

\begin{itemize}
    \item \textbf{I (Insert):} Bản ghi xuất hiện trong \texttt{*\_tdy}
          nhưng chưa tồn tại trong \texttt{*\_pdy}, đây là dữ liệu mới
          lần đầu xuất hiện.
    \item \textbf{U (Update):} Bản ghi có trong cả hai bảng nhưng có ít
          nhất một cột thay đổi giá trị, phản ánh sự cập nhật từ phía
          nguồn.
\end{itemize}

Hệ thống không sinh cờ \textbf{D (Delete)} trong chế độ chạy theo ngày.
Nguyên nhân: khi lọc theo ngày, \texttt{*\_tdy} chỉ chứa các bản ghi
có hoạt động trong ngày đó, không phải toàn bộ source. Một khách hàng
vắng mặt hôm nay do không phát sinh giao dịch không đồng nghĩa với việc
bị xóa khỏi hệ thống nguồn, nếu đánh cờ D sẽ khiến tầng Silver
xóa nhầm hàng loạt bản ghi đang hoạt động.

Sau khi tính xong change-set, \texttt{*\_tdy} được upsert vào
\texttt{*\_pdy} để cập nhật snapshot lũy kế. Nhờ vậy, khi một khách
hàng quay lại giao dịch sau nhiều ngày vắng, so sánh đúng với
lần xuất hiện gần nhất và đánh cờ U chính xác thay vì nhầm thành I.

\subsubsection*{c) Xử lý riêng cho dữ liệu giao dịch}

Giao dịch là dữ liệu bất biến, một khi ghi nhận sẽ không bao giờ bị
sửa hoặc xóa. Do đó toàn bộ bản ghi trong \texttt{transactions\_tdy}
đều được đánh cờ \textbf{I} mà không cần so sánh với \texttt{*\_pdy}.
Tầng Silver tự khử trùng lặp theo business key nên việc chạy lại cùng
một ngày vẫn đảm bảo tính idempotent.

\subsubsection*{d) Hai chế độ vận hành}

Pipeline hỗ trợ hai chế độ chạy:

\begin{itemize}
    \item \textbf{Chế độ theo ngày:} Truyền tham số \texttt{--run-date
          YYYY-MM-DD} để lọc dữ liệu theo ngày giao dịch, dùng cho lịch
          chạy hằng ngày tự động qua Airflow.
    \item \textbf{Chế độ backfill:} Không truyền tham số, hệ thống đọc
          toàn bộ dữ liệu từ nguồn, dùng cho lần nạp dữ liệu đầu tiên
          hoặc khi cần tái tạo lại toàn bộ tầng Bronze.
\end{itemize}

Các entity không có cột ngày (\texttt{users}, \texttt{cards},
\texttt{mcc\_codes}) được lọc gián tiếp bằng cách join với bảng
\texttt{transactions} theo \texttt{run\_date}, đảm bảo chỉ nạp các bản
ghi thực sự có hoạt động trong ngày đó.


\subsection{Biến đổi dữ liệu tầng Silver (Data Vault)}

Tầng Silver tổ chức dữ liệu theo chuẩn Data Vault 2.0 gồm ba loại bảng:
Hub lưu business key của từng thực thể, Link ghi nhận mối quan hệ giữa
các Hub, Satellite lưu thuộc tính mô tả và lịch sử thay đổi. Mọi bản ghi
đều mang trường \texttt{load\_datetime} và \texttt{record\_source} để đảm
bảo khả năng truy vết nguồn gốc dữ liệu.

\subsubsection*{a) Quy tắc sinh Hash Key}

Hash Key là khóa kỹ thuật duy nhất đại diện cho mỗi business key trong
Data Vault. Hệ thống sử dụng thuật toán MD5 áp dụng thống nhất qua macro
dbt để đảm bảo kết quả nhất quán trên toàn bộ các model. Với Link cần
kết hợp nhiều business key, hàm băm được áp dụng trên chuỗi nối các
business key theo thứ tự cố định. Cách tiếp cận này cho phép tái tính
lại Hash Key ở bất kỳ tầng nào mà không cần join ngược lên Hub.

\subsubsection*{b) Mapping tầng Hub}

\begin{table}[H]
\centering
\footnotesize

\begin{tabular}{lllll}
\hline
\textbf{Bảng gốc} & \textbf{Trường nguồn} & \textbf{Quy tắc transform} & \textbf{Bảng đích} & \textbf{Trường đích} \\
\hline
users\_mns        & id          & MD5(id)            & hub\_customer    & hk\_customer    \\
users\_mns        & id          & Giữ nguyên         & hub\_customer    & customer\_id    \\
\hline
cards\_mns        & id          & MD5(id)            & hub\_card        & hk\_card        \\
cards\_mns        & id          & Giữ nguyên         & hub\_card        & card\_id        \\
\hline
transactions\_mns & id          & MD5(id)            & hub\_transaction & hk\_transaction \\
transactions\_mns & id          & Giữ nguyên         & hub\_transaction & transaction\_id \\
\hline
mcc\_codes\_mns   & mcc\_id     & MD5(mcc\_id)       & hub\_mcc         & hk\_mcc         \\
mcc\_codes\_mns   & mcc\_id     & Giữ nguyên         & hub\_mcc         & mcc\_id         \\
\hline
transactions\_tdy & merchant\_id & MD5(merchant\_id) & hub\_merchant    & hk\_merchant    \\
transactions\_tdy & merchant\_id & Giữ nguyên        & hub\_merchant    & merchant\_id    \\
\hline
\end{tabular}
\caption{Mapping Bronze $\rightarrow$ Hub}
\end{table}



\subsubsection*{c) Mapping tầng Link}

\begin{table}[H]
\centering
\footnotesize

\begin{tabular}{lllll}
\hline
\textbf{Bảng gốc} & \textbf{Trường nguồn} & \textbf{Quy tắc transform} & \textbf{Bảng đích} & \textbf{Trường đích} \\
\hline
transactions\_tdy & id, client\_id,        & \texttt{MD5(id || client\_id ||}     & link\_transaction & hk\_transaction\_link \\
                  & card\_id, merchant\_id, & \texttt{card\_id || merchant\_id ||} &                   &                       \\
                  & mcc                    & \texttt{mcc)}                         &                   &                       \\
transactions\_tdy & id                     & \texttt{MD5(id)}                      & link\_transaction & hk\_transaction       \\
transactions\_tdy & client\_id             & \texttt{MD5(client\_id)}              & link\_transaction & hk\_customer          \\
transactions\_tdy & card\_id               & \texttt{MD5(card\_id)}                & link\_transaction & hk\_card              \\
transactions\_tdy & merchant\_id           & \texttt{MD5(merchant\_id)}            & link\_transaction & hk\_merchant          \\
transactions\_tdy & mcc                    & \texttt{MD5(mcc)}                     & link\_transaction & hk\_mcc               \\
\hline
cards\_tdy        & client\_id, id         & \texttt{MD5(client\_id || id)}  & link\_customer\_card & hk\_customer\_card \\
cards\_tdy        & client\_id             & \texttt{MD5(client\_id)}        & link\_customer\_card & hk\_customer       \\
cards\_tdy        & id                     & \texttt{MD5(id)}                & link\_customer\_card & hk\_card           \\
\hline
\end{tabular}
\caption{Mapping Bronze $\rightarrow$ Link}
\end{table}


\subsubsection*{d) Mapping tầng Satellite}

\begin{table}[H]
\centering
\footnotesize

\begin{tabular}{lllll}
\hline
\textbf{Bảng gốc} & \textbf{Trường nguồn} & \textbf{Quy tắc transform} & \textbf{Bảng đích} & \textbf{Trường đích} \\
\hline
users\_mns   & id                 & MD5(id)                  & sat\_customer\_profile & hk\_customer        \\
users\_tdy   & current\_age       & TRY\_CAST $\rightarrow$ INT          & sat\_customer\_profile & current\_age        \\
users\_tdy   & retirement\_age    & TRY\_CAST $\rightarrow$ INT          & sat\_customer\_profile & retirement\_age     \\
users\_tdy   & birth\_year        & TRY\_CAST $\rightarrow$ INT          & sat\_customer\_profile & birth\_year         \\
users\_tdy   & birth\_month       & TRY\_CAST $\rightarrow$ INT          & sat\_customer\_profile & birth\_month        \\
users\_tdy   & gender             & CAST $\rightarrow$ NVARCHAR(20)      & sat\_customer\_profile & gender              \\
users\_tdy   & address            & CAST $\rightarrow$ NVARCHAR(255)     & sat\_customer\_profile & address             \\
users\_tdy   & latitude           & CAST $\rightarrow$ NVARCHAR(20)      & sat\_customer\_profile & latitude            \\
users\_tdy   & longitude          & CAST $\rightarrow$ NVARCHAR(20)      & sat\_customer\_profile & longitude           \\
users\_tdy   & per\_capita\_income & TRY\_CAST $\rightarrow$ DECIMAL(12,2) & sat\_customer\_profile & per\_capita\_income \\
users\_tdy   & yearly\_income     & TRY\_CAST $\rightarrow$ DECIMAL(12,2) & sat\_customer\_profile & yearly\_income      \\
users\_tdy   & total\_debt        & TRY\_CAST $\rightarrow$ DECIMAL(12,2) & sat\_customer\_profile & total\_debt         \\
users\_tdy   & credit\_score      & TRY\_CAST $\rightarrow$ INT          & sat\_customer\_profile & credit\_score       \\
users\_tdy   & num\_credit\_cards & TRY\_CAST $\rightarrow$ INT          & sat\_customer\_profile & num\_credit\_cards  \\
             & (tất cả cột trên) & MD5(nối chuỗi)           & sat\_customer\_profile & hashdiff            \\
             & —                 & SYSUTCDATETIME()         & sat\_customer\_profile & effective\_from     \\
             & —                 & Hằng số 9999-12-31       & sat\_customer\_profile & effective\_to       \\
\hline
\end{tabular}
\caption{Mapping Bronze $\rightarrow$ sat\_customer\_profile (SCD Type 2)}
\end{table}

\begin{table}[H]
\centering
\footnotesize

\begin{tabular}{lllll}
\hline
\textbf{Bảng gốc} & \textbf{Trường nguồn} & \textbf{Quy tắc transform} & \textbf{Bảng đích} & \textbf{Trường đích} \\
\hline
cards\_mns & id                    & MD5(id)                   & sat\_card\_detail & hk\_card                \\
cards\_tdy & card\_brand           & CAST $\rightarrow$ NVARCHAR(50)       & sat\_card\_detail & card\_brand             \\
cards\_tdy & card\_type            & CAST $\rightarrow$ NVARCHAR(50)       & sat\_card\_detail & card\_type              \\
cards\_tdy & card\_number          & CAST $\rightarrow$ NVARCHAR(20)       & sat\_card\_detail & card\_number            \\
cards\_tdy & expires               & TRY\_CAST $\rightarrow$ DATE          & sat\_card\_detail & expires                 \\
cards\_tdy & cvv                   & CAST $\rightarrow$ NVARCHAR(10)       & sat\_card\_detail & cvv                     \\
cards\_tdy & has\_chip             & CAST $\rightarrow$ NVARCHAR(10)       & sat\_card\_detail & has\_chip               \\
cards\_tdy & num\_cards\_issued    & TRY\_CAST $\rightarrow$ INT           & sat\_card\_detail & num\_cards\_issued      \\
cards\_tdy & credit\_limit         & TRY\_CAST $\rightarrow$ DECIMAL(12,2) & sat\_card\_detail & credit\_limit           \\
cards\_tdy & acct\_open\_date      & TRY\_CAST $\rightarrow$ DATE          & sat\_card\_detail & acct\_open\_date        \\
cards\_tdy & year\_pin\_last\_changed & TRY\_CAST $\rightarrow$ INT        & sat\_card\_detail & year\_pin\_last\_changed \\
           & (tất cả cột trên)    & MD5(nối chuỗi)            & sat\_card\_detail & hashdiff                \\
           & —                    & SYSUTCDATETIME()          & sat\_card\_detail & effective\_from         \\
           & —                    & Hằng số 9999-12-31        & sat\_card\_detail & effective\_to           \\
\hline
\end{tabular}
\caption{Mapping Bronze $\rightarrow$ sat\_card\_detail (SCD Type 2)}
\end{table}

\begin{table}[H]
\centering
\footnotesize

\begin{tabular}{lllll}
\hline
\textbf{Bảng gốc} & \textbf{Trường nguồn} & \textbf{Quy tắc transform} & \textbf{Bảng đích} & \textbf{Trường đích} \\
\hline
transactions\_tdy & id            & MD5(id)                    & sat\_transaction\_detail & hk\_transaction       \\
transactions\_tdy & date          & TRY\_CONVERT $\rightarrow$ DATETIME    & sat\_transaction\_detail & transaction\_datetime \\
transactions\_tdy & amount        & TRY\_CAST $\rightarrow$ DECIMAL(10,2)  & sat\_transaction\_detail & amount                \\
transactions\_tdy & use\_chip     & CAST $\rightarrow$ NVARCHAR(50)        & sat\_transaction\_detail & use\_chip             \\
transactions\_tdy & merchant\_city  & CAST $\rightarrow$ NVARCHAR(100)     & sat\_transaction\_detail & merchant\_city        \\
transactions\_tdy & merchant\_state & CAST $\rightarrow$ NVARCHAR(50)      & sat\_transaction\_detail & merchant\_state       \\
transactions\_tdy & zip           & CAST $\rightarrow$ NVARCHAR(10)        & sat\_transaction\_detail & zip                   \\
transactions\_tdy & errors        & CAST $\rightarrow$ NVARCHAR(100)       & sat\_transaction\_detail & errors                \\
\hline
\end{tabular}
\caption{Mapping Bronze $\rightarrow$ sat\_transaction\_detail (SCD Type 1)}
\end{table}

\begin{table}[H]
\centering
\footnotesize

\begin{tabular}{lllll}
\hline
\textbf{Bảng gốc} & \textbf{Trường nguồn} & \textbf{Quy tắc transform} & \textbf{Bảng đích} & \textbf{Trường đích} \\
\hline
mcc\_codes\_mns & mcc\_id     & MD5(mcc\_id)         & sat\_mcc\_detail & hk\_mcc      \\
mcc\_codes\_tdy & description & CAST $\rightarrow$ NVARCHAR(255) & sat\_mcc\_detail & description  \\
\hline
\end{tabular}
\caption{Mapping Bronze $\rightarrow$ sat\_mcc\_detail (SCD Type 1)}
\end{table}


\subsection{Biến đổi dữ liệu tầng Gold (Star Schema)}

Tầng Gold tổng hợp dữ liệu từ các Hub, Link và Satellite của tầng Silver
thành cấu trúc Star Schema phục vụ phân tích. Bên cạnh các trường ánh xạ
trực tiếp, mỗi bảng bổ sung thêm các trường tính toán sẵn nhằm giảm tải
logic xử lý cho công cụ báo cáo.

\subsubsection*{a) Xây dựng bảng Dimension}

\begin{table}[H]
\centering

\resizebox{\textwidth}{!}{%
\begin{tabular}{lllll}
\hline
\textbf{Bảng gốc} & \textbf{Trường nguồn} & \textbf{Quy tắc transform} & \textbf{Bảng đích} & \textbf{Trường đích} \\
\hline
Date spine & date\_day & FORMAT $\rightarrow$ YYYYMMDD $\rightarrow$ INT         & dim\_date & date\_key      \\
Date spine & date\_day & Giữ nguyên                      & dim\_date & full\_date     \\
Date spine & date\_day & YEAR()                          & dim\_date & year           \\
Date spine & date\_day & DATEPART(quarter)               & dim\_date & quarter        \\
Date spine & date\_day & MONTH()                         & dim\_date & month          \\
Date spine & date\_day & DATENAME(month)                 & dim\_date & month\_name    \\
Date spine & date\_day & DAY()                           & dim\_date & day\_of\_month \\
Date spine & date\_day & DATEPART(weekday)               & dim\_date & day\_of\_week  \\
Date spine & date\_day & DATENAME(weekday)               & dim\_date & day\_name      \\
Date spine & date\_day & CASE weekday IN (1,7) $\rightarrow$ 1/0     & dim\_date & is\_weekend    \\
\hline
\end{tabular}
}
\caption{Mapping Silver $\rightarrow$ dim\_date}
\end{table}

\texttt{dim\_date} được sinh tự động bằng kỹ thuật cross join date-spine,
tạo liên tục các ngày từ 2022-01-01 đến 3 năm sau ngày chạy mà không cần
file seed hay bảng tham chiếu ngoài.

\begin{table}[H]
\centering

\resizebox{\textwidth}{!}{%
\begin{tabular}{lllll}
\hline
\textbf{Bảng gốc} & \textbf{Trường nguồn} & \textbf{Quy tắc transform} & \textbf{Bảng đích} & \textbf{Trường đích} \\
\hline
hub\_customer          & customer\_id              & Giữ nguyên                                    & dim\_customer & customer\_id           \\
hub\_customer          & hk\_customer              & Giữ nguyên                                    & dim\_customer & hk\_customer           \\
sat\_customer\_profile & gender                    & Giữ nguyên                                    & dim\_customer & gender                 \\
sat\_customer\_profile & address                   & Giữ nguyên                                    & dim\_customer & address                \\
sat\_customer\_profile & yearly\_income            & Giữ nguyên                                    & dim\_customer & yearly\_income         \\
sat\_customer\_profile & credit\_score             & Giữ nguyên                                    & dim\_customer & credit\_score          \\
sat\_customer\_profile & current\_age              & Giữ nguyên                                    & dim\_customer & current\_age           \\
sat\_customer\_profile & retirement\_age           & Giữ nguyên                                    & dim\_customer & retirement\_age        \\
sat\_customer\_profile & birth\_year               & Giữ nguyên                                    & dim\_customer & birth\_year            \\
sat\_customer\_profile & birth\_month              & Giữ nguyên                                    & dim\_customer & birth\_month           \\
sat\_customer\_profile & latitude                  & TRY\_CAST $\rightarrow$ DECIMAL(10,6)                     & dim\_customer & latitude               \\
sat\_customer\_profile & longitude                 & TRY\_CAST $\rightarrow$ DECIMAL(10,6)                     & dim\_customer & longitude              \\
sat\_customer\_profile & per\_capita\_income       & Giữ nguyên                                    & dim\_customer & per\_capita\_income    \\
sat\_customer\_profile & total\_debt               & Giữ nguyên                                    & dim\_customer & total\_debt            \\
sat\_customer\_profile & num\_credit\_cards        & Giữ nguyên                                    & dim\_customer & num\_credit\_cards     \\
sat\_customer\_profile & total\_debt, yearly\_income & total\_debt / NULLIF(yearly\_income, 0)     & dim\_customer & debt\_to\_income\_ratio \\
sat\_customer\_profile & yearly\_income            & CASE: $<$30k/60k/100k $\rightarrow$ Low/Medium/High/Premium & dim\_customer & income\_segment        \\
sat\_customer\_profile & credit\_score             & CASE: $\geq$750/670/580 $\rightarrow$ Excellent/Good/Fair/Poor & dim\_customer & credit\_risk\_tier     \\
sat\_customer\_profile & retirement\_age, current\_age & retirement\_age - current\_age            & dim\_customer & years\_to\_retirement  \\
\hline
\end{tabular}
}
\caption{Mapping Silver $\rightarrow$ dim\_customer}
\end{table}

Chỉ lấy record có \texttt{effective\_to = 9999-12-31} từ Satellite (trạng
thái hiện tại), đảm bảo mỗi khách hàng chỉ có một dòng trong bảng chiều.

\begin{table}[H]
\centering

\resizebox{\textwidth}{!}{%
\begin{tabular}{lllll}
\hline
\textbf{Bảng gốc} & \textbf{Trường nguồn} & \textbf{Quy tắc transform} & \textbf{Bảng đích} & \textbf{Trường đích} \\
\hline
hub\_card          & card\_id                  & Giữ nguyên                                            & dim\_card & card\_id                 \\
hub\_card          & hk\_card                  & Giữ nguyên                                            & dim\_card & hk\_card                 \\
hub\_customer      & customer\_id              & Denormalize qua link\_customer\_card                  & dim\_card & customer\_id             \\
sat\_card\_detail  & card\_brand               & Giữ nguyên                                            & dim\_card & card\_brand              \\
sat\_card\_detail  & card\_type                & Giữ nguyên                                            & dim\_card & card\_type               \\
sat\_card\_detail  & credit\_limit             & Giữ nguyên                                            & dim\_card & credit\_limit            \\
sat\_card\_detail  & expires                   & Giữ nguyên                                            & dim\_card & expires                  \\
sat\_card\_detail  & has\_chip                 & Giữ nguyên                                            & dim\_card & has\_chip                \\
sat\_card\_detail  & num\_cards\_issued        & Giữ nguyên                                            & dim\_card & num\_cards\_issued       \\
sat\_card\_detail  & acct\_open\_date          & Giữ nguyên                                            & dim\_card & acct\_open\_date         \\
sat\_card\_detail  & year\_pin\_last\_changed  & Giữ nguyên                                            & dim\_card & year\_pin\_last\_changed \\
sat\_card\_detail  & acct\_open\_date          & DATEDIFF(YEAR, acct\_open\_date, GETDATE())           & dim\_card & card\_age\_years         \\
sat\_card\_detail  & year\_pin\_last\_changed  & YEAR(GETDATE()) - year\_pin\_last\_changed            & dim\_card & pin\_age\_years          \\
\hline
\end{tabular}
}
\caption{Mapping Silver $\rightarrow$ dim\_card}
\end{table}

Trường \texttt{card\_number} và \texttt{cvv} bị loại bỏ ở tầng Gold do
yêu cầu bảo mật dữ liệu thẻ (PCI DSS).

\begin{table}[H]
\centering

\resizebox{\textwidth}{!}{%
\begin{tabular}{lllll}
\hline
\textbf{Bảng gốc} & \textbf{Trường nguồn} & \textbf{Quy tắc transform} & \textbf{Bảng đích} & \textbf{Trường đích} \\
\hline
hub\_merchant            & merchant\_id    & Giữ nguyên                  & dim\_merchant & merchant\_id    \\
hub\_merchant            & hk\_merchant    & Giữ nguyên                  & dim\_merchant & hk\_merchant    \\
sat\_transaction\_detail & merchant\_city  & Lấy từ giao dịch gần nhất   & dim\_merchant & merchant\_city  \\
sat\_transaction\_detail & merchant\_state & Lấy từ giao dịch gần nhất   & dim\_merchant & merchant\_state \\
sat\_transaction\_detail & zip             & Lấy từ giao dịch gần nhất   & dim\_merchant & zip             \\
\hline
\end{tabular}
}
\caption{Mapping Silver $\rightarrow$ dim\_merchant}
\end{table}

\texttt{dim\_merchant} không có Satellite riêng ở tầng Silver. Thông tin
địa điểm được suy ra từ \texttt{sat\_transaction\_detail}, lấy bản ghi
giao dịch gần nhất theo thời gian thông qua hàm \texttt{ROW\_NUMBER()}.

\begin{table}[H]
\centering

\resizebox{\textwidth}{!}{%
\begin{tabular}{lllll}
\hline
\textbf{Bảng gốc} & \textbf{Trường nguồn} & \textbf{Quy tắc transform} & \textbf{Bảng đích} & \textbf{Trường đích} \\
\hline
hub\_mcc        & mcc\_id      & Giữ nguyên & dim\_mcc & mcc\_id          \\
hub\_mcc        & hk\_mcc      & Giữ nguyên & dim\_mcc & hk\_mcc          \\
sat\_mcc\_detail & description & Giữ nguyên & dim\_mcc & mcc\_description \\
\hline
\end{tabular}
}
\caption{Mapping Silver $\rightarrow$ dim\_mcc}
\end{table}

\subsubsection*{b) Xây dựng bảng Fact}

\begin{table}[H]
\centering

\resizebox{\textwidth}{!}{%
\begin{tabular}{lllll}
\hline
\textbf{Bảng gốc} & \textbf{Trường nguồn} & \textbf{Quy tắc transform} & \textbf{Bảng đích} & \textbf{Trường đích} \\
\hline
hub\_transaction         & transaction\_id       & Giữ nguyên                               & fact\_transaction & transaction\_id       \\
sat\_transaction\_detail & transaction\_datetime & FORMAT $\rightarrow$ YYYYMMDD $\rightarrow$ INT (FK dim\_date)    & fact\_transaction & date\_key             \\
hub\_customer            & customer\_id          & Giữ nguyên (FK dim\_customer)            & fact\_transaction & customer\_id          \\
hub\_card                & card\_id              & Giữ nguyên (FK dim\_card)                & fact\_transaction & card\_id              \\
hub\_merchant            & merchant\_id          & Giữ nguyên (FK dim\_merchant)            & fact\_transaction & merchant\_id          \\
hub\_mcc                 & mcc\_id               & Giữ nguyên (FK dim\_mcc)                 & fact\_transaction & mcc\_id               \\
sat\_transaction\_detail & transaction\_datetime & Giữ nguyên                               & fact\_transaction & transaction\_datetime \\
sat\_transaction\_detail & amount                & Giữ nguyên                               & fact\_transaction & amount                \\
sat\_transaction\_detail & use\_chip             & Giữ nguyên                               & fact\_transaction & use\_chip             \\
sat\_transaction\_detail & merchant\_city        & Giữ nguyên                               & fact\_transaction & merchant\_city        \\
sat\_transaction\_detail & merchant\_state       & Giữ nguyên                               & fact\_transaction & merchant\_state       \\
sat\_transaction\_detail & zip                   & Giữ nguyên                               & fact\_transaction & zip                   \\
sat\_transaction\_detail & errors                & Giữ nguyên                               & fact\_transaction & errors                \\
sat\_transaction\_detail & errors                & CASE: errors không rỗng $\rightarrow$ 1, còn lại $\rightarrow$ 0 & fact\_transaction & has\_error            \\
sat\_transaction\_detail & use\_chip             & CASE: = 'Chip Transaction' $\rightarrow$ 1, còn lại $\rightarrow$ 0 & fact\_transaction & is\_chip\_used     \\
\hline
\end{tabular}
}
\caption{Mapping Silver $\rightarrow$ fact\_transaction}
\end{table}

Bảng Fact sử dụng chiến lược nạp \texttt{delete+insert} theo ngày chạy,
đảm bảo tính idempotent khi pipeline chạy lại cùng một ngày. Khóa ngoại
\texttt{date\_key} được tính bằng cách định dạng \texttt{transaction\_datetime}
thành số nguyên dạng \texttt{YYYYMMDD} để join trực tiếp với \texttt{dim\_date}.




\section{Quy trình xây dựng Pipeline dữ liệu phi cấu trúc (OCR)}


% Bạn có thể viết một đoạn mở đầu ngắn (Intro) cho cả phần 4.x ở đây trước khi đi vào các tiểu mục.



\subsection{Tiền xử lý ảnh }
% Trình bày các kỹ thuật OpenCV đã dùng (cắt viền, chỉnh nghiêng, khử nhiễu, cân bằng sáng...).



Mục tiêu: Đưa ảnh chụp thực tế (xiên, lệch góc, sai hướng) về dạng tài liệu thẳng, chuẩn kích thước trước khi OCR.

\begin{figure}[H]
    \centering
    \begin{tikzpicture}[
        >=Stealth,
        node distance=0.7cm,
        % Style cho khối Đầu vào / Đầu ra
        iobox/.style={rectangle, rounded corners=3pt, draw=gray!70, fill=gray!10, thick, text width=10.5cm, align=center, inner sep=8pt, font=\bfseries},
        % Style cho khối Xử lý
        procbox/.style={rectangle, rounded corners=3pt, draw=blue!60, fill=blue!5, thick, text width=10.5cm, align=left, inner sep=8pt},
        % Style cho mũi tên
        arr/.style={->, thick, draw=gray!80}
    ]

    % --- Khối đầu vào ---
    \node[iobox] (input) {Ảnh gốc};

    % --- Các bước xử lý ---
    \node[procbox, below=of input] (step1) {
        \textbf{Chỉnh nghiêng (Deskew)} \\
        \footnotesize Phát hiện góc xiên qua Hough Lines, xoay ảnh về góc thẳng.
    };

    \node[procbox, below=of step1] (step2) {
        \textbf{Căn phẳng phối cảnh (Document Warp)} \\
        \footnotesize Tìm tọa độ 4 góc tài liệu, áp dụng biến đổi phối cảnh 
    };

    \node[procbox, below=of step2] (step3) {
        \textbf{Căn chỉnh chiều (Orientation)} \\
        \footnotesize Xoay ảnh về đúng chiều nhìn (Ví dụ CCCD: Dựa vào vị trí dải màu xanh đặc trưng).
    };

    \node[procbox, below=of step3] (step4) {
        \textbf{Chuẩn hóa kích thước (Width Norm)} \\
        \footnotesize Đưa ảnh về một độ rộng tiêu chuẩn cố định (CCCD: 1000px / Sổ tiết kiệm: 1200px).
    };

    % --- Khối đầu ra ---
    \node[iobox, below=of step4] (output) {Ảnh chuẩn hóa};

    % --- Vẽ mũi tên kết nối ---
    \foreach \a/\b in {input/step1, step1/step2, step2/step3, step3/step4, step4/output}
    { \draw[arr] (\a) -- (\b); }

    \end{tikzpicture}
    \caption{Quy trình chuẩn hóa hình học trong tiền xử lý ảnh}
    \label{fig:image-geometry-normalization}
\end{figure}

Sau bước này, mọi ảnh đầu vào đều là hình chữ nhật thẳng, đúng hướng, cùng tỷ lệ, tọa độ vùng trích xuất mới có ý nghĩa nhất quán.

Để minh họa trực quan cho quy trình trên, dưới đây là chi tiết kết quả hình ảnh đạt được sau từng bước biến đổi hình học:

\subsubsection*{a) Chỉnh nghiêng ảnh (Deskew)}
\label{sssec:deskew}

Ảnh CCCD chụp trong thực tế thường bị nghiêng một góc nhỏ so với phương ngang. Việc chỉnh nghiêng dựa trên một quan sát: thẻ CCCD chứa nhiều đường thẳng dài song song với cạnh thẻ (viền thẻ, mép dải họa tiết, các dòng chữ). Nếu đo được góc nghiêng của các đường thẳng này so với phương ngang thì cũng xác định được góc nghiêng của cả tấm thẻ, và chỉ cần xoay ngược lại để sửa. Quy trình gồm ba giai đoạn: (1)~ước lượng góc nghiêng, (2)~xoay ảnh, và (3)~kiểm chứng kết quả.

\paragraph{Giai đoạn 1: Ước lượng góc nghiêng:}
Thay vì dò đường thẳng trên toàn ảnh (dễ nhiễu bởi chữ, hoa văn và nền), hệ thống ưu tiên sử dụng dải màu xanh dương đặc trưng của CCCD. Ảnh được chuyển sang không gian màu HSV và lọc theo ngưỡng kênh Hue:
\begin{equation}
  85 \le H \le 135,
\end{equation}
tương ứng với dải màu xanh dương. Cạnh của vùng xanh luôn song song với cạnh thẻ nên là mốc đo góc ít nhiễu. Trường hợp ảnh không có đủ vùng xanh (bạc màu, chói sáng), phương án dự phòng là dò trên toàn ảnh xám đã làm mờ Gaussian, nhưng yêu cầu đoạn thẳng phải đủ dài (tối thiểu 25\% cạnh ngắn của ảnh) để chỉ bắt được viền thẻ thay vì nét chữ.

Trên vùng đã chọn, hệ thống áp dụng thuật toán dò biên Canny rồi biến đổi Hough xác suất (Probabilistic Hough Transform) để thu về tập các đoạn thẳng, mỗi đoạn cho bởi hai đầu mút $(x_1, y_1)$ và $(x_2, y_2)$. Góc của mỗi đoạn so với phương ngang được tính theo:
\begin{equation}
  \theta = \arctan\!\left(\frac{y_2 - y_1}{x_2 - x_1}\right) \cdot \frac{180}{\pi}.
\end{equation}

Vì trên thẻ có cả đường ngang lẫn đường dọc -- hai loại này lệch nhau $90^\circ$ nhưng đều phản ánh cùng một độ nghiêng -- nên góc được quy gọn về khoảng $[-45^\circ, 45^\circ]$:
\begin{equation}
  \theta' =
  \begin{cases}
    \theta - 90^\circ & \text{nếu } \theta > 45^\circ, \\
    \theta + 90^\circ & \text{nếu } \theta < -45^\circ, \\
    \theta            & \text{còn lại.}
  \end{cases}
\end{equation}
Sau khi quy gọn, một đường dọc nghiêng $88^\circ$ và một đường ngang nghiêng $-2^\circ$ đều cho cùng kết quả $-2^\circ$. Các góc có $|\theta'| > 30^\circ$ bị loại, vì độ nghiêng lớn cỡ đó không còn là ``nghiêng nhẹ'' mà là ảnh bị xoay $90^\circ$, được xử lý ở bước chuẩn hóa hướng riêng (dựa vào vị trí dải xanh để xoay $90/180/270^\circ$).

Góc nghiêng của ảnh được tổng hợp bằng trung vị có trọng số (weighted median) của các góc, với trọng số là chiều dài đoạn thẳng:
\begin{equation}
  w_i = \sqrt{(x_2 - x_1)^2 + (y_2 - y_1)^2}.
\end{equation}
Việc lấy trọng số theo chiều dài khiến các đoạn dài (viền thẻ) có ảnh hưởng lớn hơn các nét chữ, hoa văn ngắn; còn phép trung vị bền vững trước ngoại lệ hơn phép trung bình -- vài đoạn nhiễu lệch góc hẳn sẽ không kéo lệch kết quả.

\paragraph{Giai đoạn 2: Xoay ảnh không mất góc:}
Nếu góc ước lượng nhỏ hơn ngưỡng $0{,}75^\circ$ thì ảnh được giữ nguyên, vì xoay một góc quá nhỏ chỉ làm mờ chữ do nội suy. Khi xoay ảnh quanh tâm một góc $\alpha$, nếu giữ nguyên kích thước khung ảnh thì bốn góc của thẻ sẽ bị cắt mất; do đó khung ảnh đầu ra được mở rộng theo:
\begin{equation}
  W' = H \cdot |\sin\alpha| + W \cdot |\cos\alpha|, \qquad H' = H \cdot |\cos\alpha| + W \cdot |\sin\alpha|,
\end{equation}
với $(W, H)$ là kích thước ảnh gốc. Ma trận xoay affine đồng thời được tịnh tiến để tâm ảnh cũ trùng với tâm khung mới; phần nền lộ ra sau khi xoay được tô màu trắng, và phép nội suy bicubic được dùng để giữ độ sắc nét của chữ.

\paragraph{Giai đoạn 3: Kiểm chứng bằng góc dư:}
Góc ước lượng có thể sai dấu (không phân biệt được thẻ nghiêng trái hay phải trong một số trường hợp nhiễu). Thay vì tin tưởng kết quả đo, thuật toán tạo ra ba ứng viên: ảnh gốc (không xoay), ảnh xoay góc $+\alpha$ và ảnh xoay góc $-\alpha$. Với mỗi ứng viên, toàn bộ phép ước lượng góc ở Giai đoạn~1 được chạy lại để đo góc nghiêng còn dư (residual), rồi chọn ứng viên có góc dư nhỏ nhất:
\begin{equation}
  \text{kết quả} = \operatorname*{arg\,min}_{k \,\in\, \{0,\, +\alpha,\, -\alpha\}} \left| \theta_{\mathrm{residual}}(k) \right|.
\end{equation}
Cơ chế ``thử cả hai chiều rồi đo lại'' này đảm bảo: nếu phép đo ban đầu đúng thì ảnh được sửa thẳng; nếu phép đo sai (sai dấu hoặc sai hẳn) thì ảnh gốc vẫn được giữ nguyên --- thuật toán không bao giờ làm ảnh xấu đi

\begin{figure}[H]
    \centering

    \includegraphics[width=0.5\textwidth]{img4/id_card_scan.jpg} 
    \caption{Ảnh trước khi xử lý Deskew}
    \label{fig:before-deskew}
\end{figure}

\begin{figure}[H]
    \centering

    \includegraphics[width=0.5\textwidth]{img4/sau deskew.png} 
    \caption{Ảnh sau xử lý deskew}
    \label{fig:after-deskew}
\end{figure}

\subsubsection*{b) Căn phẳng phối cảnh (Document Warp)}
\label{sssec:perspective-warp}

Ảnh CCCD chụp trong thực tế thường xuyên bị méo phối cảnh do góc nghiêng của thiết bị chụp --- hình dáng thẻ bị biến dạng thành hình thang thay vì hình chữ nhật. Bước tiền xử lý này có nhiệm vụ xác định chính xác tọa độ bốn góc của thẻ trên ảnh, từ đó thực hiện phép chiếu hình học để ``kéo phẳng'' tài liệu về đúng chuẩn chữ nhật thẳng đứng. Khâu này đóng vai trò then chốt giúp các vùng quan tâm (ROI) luôn khớp vị trí khi đưa vào mô hình trích xuất OCR.

\paragraph{Giai đoạn 1: Xác định tọa độ bốn góc thẻ:}
Hệ thống áp dụng thuật toán dò biên Canny và trích xuất các dải contour khép kín. Mỗi contour được xấp xỉ thành một đa giác hình học. Một contour sẽ được hệ thống phân loại là khung bao thẻ nếu nó xấp xỉ thành đúng bốn đỉnh và thỏa mãn đồng thời hai điều kiện:
\begin{itemize}
    \item Diện tích bao phủ chiếm ít nhất $25\%$ tổng diện tích bức ảnh (nhằm loại trừ các tứ giác nhiễu nhỏ nằm lẫn trong nền).
    \item Tỷ lệ khung hình (cạnh dài chia cạnh ngắn) nằm trong khoảng $[1{,}25;\; 1{,}95]$. Biên độ này được thiết lập dựa trên chuẩn kích thước vật lý của thẻ ID-1 quốc tế (tỷ lệ ISO $\approx 1{,}586$).
\end{itemize}
Trong trường hợp không có contour nào đạt chuẩn bốn đỉnh, thuật toán sẽ kích hoạt phương án dự phòng bằng cách tính toán hình chữ nhật bao có diện tích nhỏ nhất của contour lớn nhất.

\begin{figure}[H]
    \centering

    \includegraphics[width=0.5\textwidth]{img4/viền.png} 
    \caption{Xác định khung viền thẻ}
    \label{fig:card-border-detection}
\end{figure}

\paragraph{Giai đoạn 2: Biến đổi phối cảnh: }
Sau khi thu thập được tọa độ bốn góc thực tế của thẻ, ký hiệu là $\{P_{\text{TL}}, P_{\text{TR}}, P_{\text{BR}}, P_{\text{BL}}\}$, hệ thống thiết lập một ma trận biến đổi phối cảnh $\mathbf{H}$ kích thước $3 \times 3$. Ma trận này có nhiệm vụ ánh xạ tuyến tính bốn góc thực tế về bốn góc vuông của một ảnh đích có kích thước $W_{\text{out}} \times H_{\text{out}}$ (với bề ngang chuẩn $W_{\text{out}} = 1\,000\,\text{px}$):
\begin{equation}
    \begin{pmatrix}P_{\text{TL}} \\ P_{\text{TR}} \\ P_{\text{BR}} \\ P_{\text{BL}}\end{pmatrix}
    \;\xrightarrow{\mathbf{H}}\;
    \begin{pmatrix}(0,\,0) \\ (W_{\text{out}}-1,\,0) \\ (W_{\text{out}}-1,\,H_{\text{out}}-1) \\ (0,\,H_{\text{out}}-1)\end{pmatrix}.
\end{equation}
Toàn bộ dữ liệu điểm ảnh sau đó được chiếu qua ma trận $\mathbf{H}$, trả về một bức ảnh thẻ phẳng phiu như được chụp trực diện từ trên xuống. Nếu thuật toán thất bại trong việc tìm bốn góc hợp lệ ở khâu trước, phép biến đổi này sẽ bị bỏ qua và ảnh chỉ được thu phóng bề ngang về mức tiêu chuẩn.
\begin{figure}[H]
    \centering

    \includegraphics[width=0.5\textwidth]{img4/sau phát hiện khung.png} 
    \caption{Ảnh sau biến đổi phối cảnh}
    \label{fig:after-perspective-transform}
\end{figure}
\subsubsection*{c) Căn chỉnh chiều ảnh (Orientation)}
\label{sssec:orientation}

Sau khi warp phẳng, thẻ có thể vẫn đang nằm ngang hoặc lộn ngược nếu ảnh gốc được chụp xoay. Bước này có nhiệm vụ xác định đúng chiều mặt trước của thẻ và tự động xoay ảnh về đúng hướng.

Ý tưởng thuật toán dựa trên một đặc điểm nhận diện hình ảnh cố định của CCCD: dải màu xanh dương (chứa quốc huy và họa tiết header) luôn nằm ở cạnh trên của mặt trước thẻ. Do đó, nếu định vị được cạnh nào của ảnh đang tập trung nhiều pixel màu xanh nhất, ta có thể suy ra hướng hiện tại của thẻ và góc cần bù để xoay cho đúng.

Ảnh được chuyển đổi sang không gian màu HSV, sau đó phân tách để lấy ra các điểm ảnh thuộc dải phổ màu xanh dương ($85 \le H \le 135$). Với mỗi cạnh trong bốn cạnh của ảnh, hệ thống tính toán tỷ lệ điểm xanh xuất hiện trong vùng không gian bao phủ 22\% chiều rộng (hoặc chiều cao) tính từ mép cạnh đó đổ vào trong. 

Cạnh nào có mật độ điểm xanh cao nhất và đồng thời vượt qua ngưỡng tối thiểu $3\%$ sẽ được xác định là cạnh chứa dải header; từ đó suy ra góc cần xoay để đưa dải xanh về phía trên:
\begin{equation}
  \theta_{\text{xoay}} =
  \begin{cases}
    0^\circ   & \text{nếu dải xanh ở cạnh trên (ảnh đã đúng chiều),} \\
    180^\circ & \text{nếu dải xanh ở cạnh dưới (ảnh bị lộn ngược),} \\
    90^\circ  & \text{nếu dải xanh ở cạnh trái (xoay cùng chiều kim đồng hồ),} \\
    -90^\circ & \text{nếu dải xanh ở cạnh phải (xoay ngược chiều kim đồng hồ).}
  \end{cases}
\end{equation}

Trong trường hợp không có vùng biên nào đạt ngưỡng mật độ điểm xanh $3\%$ (thường xảy ra đối với ảnh chụp bản sao đen trắng, thẻ bị chói sáng mất màu hoặc là mặt sau của thẻ), thuật toán sẽ bỏ qua thao tác này và giữ nguyên hướng ban đầu của ảnh.

Kết thúc bước định hướng, hình ảnh thẻ đã hoàn toàn ngay ngắn, đúng chiều trực diện và sẵn sàng chuyển sang khâu chuẩn hóa kích thước trước khi đi vào mô hình trích xuất văn bản OCR.
\begin{figure}[H]
    \centering
    \begin{tikzpicture}[
        >=Stealth, 
        node distance=2.5cm % Khoảng cách giữa 2 ảnh để lấy chỗ vẽ mũi tên
    ]
        
        % --- Ảnh 1: Trước khi xoay ---
        \node[inner sep=0pt] (img1) {
            % Lưu ý: Hãy đảm bảo tên file trong thư mục img4 là truoc_xoay.png hoặc truoc_xoay.jpg
            \includegraphics[height=5.5cm]{img4/trước xoay.png} 
        };
        \node[below=0.2cm of img1, font=\small] {a) Ảnh thẻ bị ngang};

        % --- Ảnh 2: Sau khi xoay ---
        \node[inner sep=0pt, right=of img1] (img2) {
            % Lưu ý: Hãy đảm bảo tên file trong thư mục img4 là sau_xoay.png hoặc sau_xoay.jpg
            \includegraphics[height=5.5cm]{img4/sau xoay.png} 
        };
        \node[below=0.2cm of img2, font=\small] {b) Ảnh thẻ sau khi chuẩn hóa};

        % --- Mũi tên chỉ chiều biến đổi ---
        \draw[->, line width=1.5pt, draw=blue!70] 
            (img1.east) -- 
            node[above, font=\footnotesize\bfseries, text=black] {Xoay $270^\circ$} 
            (img2.west);
            
    \end{tikzpicture}
    \caption{Kết quả tự động nhận diện dải màu xanh và căn chỉnh chiều ảnh}
    \label{fig:orientation_result}
\end{figure}

\subsubsection*{d) Chuẩn hóa kích thước ảnh (Width Normalization)}
\label{sssec:width-norm}

Ảnh CCCD đầu vào trong thực tế có thể được thu thập từ nhiều nguồn thiết bị khác nhau (điện thoại di động, máy quét chuyên dụng, webcam văn phòng\ldots) dẫn đến sự phân mảnh lớn về độ phân giải. Sự chênh lệch kích thước này gây ra trở ngại trực tiếp cho bước trích xuất: các vùng quan tâm (ROI) chứa thông tin nghiệp vụ thường được neo dựa trên các tọa độ khuôn mẫu cố định. Nếu không có sự đồng nhất về mặt kích thước, việc cắt ảnh theo tọa độ pixel sẽ bị sai lệch nghiêm trọng, dẫn đến hiện tượng mất chữ hoặc cắt lẹm sang các vùng nhiễu.

Để giải quyết triệt để vấn đề này, hệ thống áp dụng kỹ thuật chuẩn hóa kích thước, đưa toàn bộ ảnh thẻ về chung một bề rộng tiêu chuẩn $W_{\text{out}} = 1000\,\text{px}$. Chiều cao $H_{\text{out}}$ được tính toán tự động dựa trên tỷ lệ khung hình gốc nhằm đảm bảo thẻ không bị méo lệch vật lý:
\begin{equation}
    H_{\text{out}} = \left\lfloor H_{\text{gốc}} \times \frac{W_{\text{out}}}{W_{\text{gốc}}} \right\rfloor
\end{equation}


\subsubsection*{e) Xác định vùng ROI}
\label{sssec:roi}

Mọi thẻ CCCD có cùng bố cục in theo chuẩn nhà nước, và sau khi ảnh đã
được warp phẳng và chuẩn hóa về bề rộng 1000\,px, vị trí tương đối của
từng trường là hằng số trên mọi ảnh. Nhờ đó, tọa độ các vùng ROI chỉ cần
xác định một lần từ một ảnh mẫu đã chuẩn hóa: quan sát vị trí
từng trường, đo tọa độ pixel rồi quy về tỷ lệ so với kích thước ảnh:
\begin{equation}
  x_{\text{ratio}} = \frac{x_{\text{px}}}{W}, \qquad
  y_{\text{ratio}} = \frac{y_{\text{px}}}{H}.
\end{equation}
Dùng tỷ lệ thay vì pixel tuyệt đối giúp ROI tự co giãn theo chiều cao
thực tế của từng ảnh, vì chiều cao có thể dao động nhẹ dù chiều rộng đã
cố định ở 1000\,px.

\begin{figure}[H]
    \centering

    \includegraphics[width=0.5\textwidth]{img4/xác định vùng ROI.png} 
    \caption{Thẻ sau khi xử lý xác định các vùng ROI}
    \label{fig:roi-region-detection}
\end{figure}



\subsection{Tăng cường chất lượng hình ảnh}

Từ ảnh thẻ đã chuẩn hóa hình học, hệ thống tạo thêm năm biến thể tăng
cường theo chuỗi tuần tự, mỗi bước xây dựng trên kết quả của bước trước:

\begin{itemize}
  \item \textbf{Grayscale}: chuyển ảnh màu về ảnh xám, loại bỏ nhiễu màu
        không liên quan đến văn bản, đồng thời là đầu vào cho các bước tiếp theo.
  \item \textbf{CLAHE}: áp dụng cân bằng độ tương phản cục bộ theo lưới
        ô nhỏ lên ảnh xám, giúp chữ trên nền loang lổ hoặc thiếu sáng
        trở nên rõ ràng hơn.
  \item \textbf{Khử nhiễu}: lọc nhiễu hạt trên ảnh CLAHE bằng phương pháp
        Fast Non-Local Means, giữ nguyên nét chữ trong khi làm mịn các
        vùng nền.
  \item \textbf{Làm sắc nét}: tăng cường biên ký tự bằng bộ lọc Laplacian
        trên ảnh đã khử nhiễu, giúp OCR phân biệt các nét gần nhau rõ hơn.
  \item \textbf{Nhị phân hóa thích nghi}: chuyển ảnh đã làm sắc nét về
        ảnh đen trắng bằng ngưỡng Gaussian thích nghi cục bộ, tạo ra
        văn bản trắng trên nền đen hoàn toàn là biến thể phù hợp nhất
        khi nền ảnh không đồng đều.
\end{itemize}

Năm biến thể này không thay thế ảnh gốc trong luồng chính.
PaddleOCR vẫn chạy trên ảnh màu đầy đủ vì ảnh màu giữ nhiều thông tin
nhất. Các biến thể tăng cường chỉ được đưa vào OCR khi một trường cụ thể
thất bại ở bước trích xuất, lúc đó vùng ROI của trường đó được cắt ra từ
ảnh CLAHE để OCR lại. Riêng trường họ tên, vốn dễ bị watermark che và
nhạy cảm với chất lượng ảnh sẽ lần lượt thử trên tất cả các biến thể,
ưu tiên kết quả hợp lệ về mặt nghiệp vụ trước, điểm tin cậy OCR sau.


\begin{figure}[H]
    \centering

    \includegraphics[width=1\textwidth]{img4/tăng cường chất lượng.png} 
    \caption{Ảnh sau tăng cường chất lượng}
    \label{fig:image-quality-enhancement}
\end{figure}







\subsection{Trích xuất văn bản bằng PaddleOCR}
Hệ thống dùng PaddleOCR (PP-OCRv3) để trích xuất 9 trường thông tin theo
hai chế độ:

\textbf{Chế độ chính: Neo theo nhãn:} OCR toàn ảnh, tìm box khớp với
nhãn in sẵn trên thẻ (``FULL NAME'', ``ID NO''\ldots), rồi lấy text của
box nằm bên phải cùng dòng làm giá trị.

\textbf{Chế độ dự phòng: Cắt ROI:} Trường nào chưa hợp lệ thì cắt
thẳng vùng tọa độ cố định trên ảnh tăng cường rồi OCR lại; trường
họ tên thử thêm nhiều biến thể ảnh, ưu tiên kết quả hợp lệ trước, điểm
OCR sau.

\begin{figure}[H]
    \centering

    \includegraphics[width=0.8\textwidth]{img4/khung xanh dương.png} 
    \caption{Các text box được phát hiện bằng chế độ dự phòng cắt ROI}
    \label{fig:roi-text-box-detection}
\end{figure}

\subsection{Hậu xử lý và làm sạch dữ liệu}

Văn bản thô trả về từ OCR thường chứa nhiễu ký tự, khoảng trắng thừa
hoặc ký tự bị nhận dạng sai. Hệ thống áp dụng chuỗi bốn bước hậu xử lý
để chuyển kết quả OCR thô thành dữ liệu có cấu trúc và nhất quán.

\textbf{Bước 1: Làm sạch văn bản cơ bản}

Toàn bộ chuỗi văn bản được chuẩn hóa Unicode về dạng NFC để đảm bảo
so khớp nhất quán. Các ký tự gây nhiễu phổ biến được sửa tự động, điển
hình là thay thế ký tự \texttt{|} bằng chữ \texttt{I}, lỗi nhầm lẫn
thường gặp khi OCR gặp chữ hoa I trên nền sáng. Khoảng trắng thừa và
ký tự đặc biệt không liên quan cũng được loại bỏ.

\textbf{Bước 2: Loại bỏ watermark}

Ảnh CCCD mẫu thường in chồng các chuỗi định danh như ``NOT A REAL ID'',
``DEMO'', ``SAMPLE''\ lên vùng văn bản, khiến các chuỗi này lẫn vào
kết quả OCR của trường họ tên hoặc địa chỉ. Hệ thống dùng biểu thức
chính quy phát hiện và xóa các mẫu watermark đã biết, xử lý
cả ở cấp độ chuỗi lẫn từng token riêng lẻ nhằm tránh xóa nhầm từ hợp lệ.

\textbf{Bước 3: Chuẩn hóa theo từng trường nghiệp vụ}

Mỗi trường thông tin được áp dụng luật chuẩn hóa riêng phù hợp với ngữ
nghĩa của nó. Trường giới tính được chuẩn hóa về hai giá trị thống nhất
\texttt{Male} / \texttt{Female}; quốc tịch về \texttt{Vietnam}; họ tên
được chuyển về dạng Title Case. Các trường địa chỉ chỉ gộp khoảng trắng
thừa mà không thay đổi nội dung.

\textbf{Bước 4: Kiểm tra tính hợp lệ}

Sau chuẩn hóa, từng trường được kiểm tra xem kết quả có hợp lý về mặt
nghiệp vụ hay không thông qua các luật đặc thù: số CCCD phải chứa
ít nhất 6 chữ số liên tiếp; các trường ngày phải khớp định dạng
\texttt{DD/MM/YYYY}; giới tính chỉ nhận một trong bốn giá trị hợp lệ;
trường văn bản tự do phải đủ độ dài tối thiểu. Trường nào không vượt qua
kiểm tra sẽ được đưa ngược lại bước OCR dự phòng với vùng ảnh tăng cường
tương ứng.


\subsection{Đánh giá độ tin cậy}

Sau khi trích xuất và hậu xử lý, mỗi bản ghi được gắn một điểm tin cậy
tổng hợp nhằm phản ánh mức độ đáng tin của toàn bộ kết quả nhận dạng.
Điểm này được tính từ hai thành phần độc lập.

\textbf{Thành phần 1: Điểm tin cậy OCR}

Trong quá trình nhận dạng, PaddleOCR trả về một điểm tin cậy cho từng
text box, thể hiện mức độ chắc chắn của mô hình với chuỗi ký tự vừa nhận
dạng. Điểm tin cậy OCR của toàn bản ghi là trung bình cộng các điểm này
trên tất cả các trường, nhận giá trị trong khoảng $[0, 1]$.

\textbf{Thành phần 2: Điểm hợp lệ nghiệp vụ}

Thành phần này đánh giá kết quả sau hậu xử lý có ý nghĩa về mặt dữ liệu
hay không, hoàn toàn độc lập với mô hình OCR. Điểm được tính theo công thức:

\begin{equation}
    \text{parse\_confidence} = \frac{\text{số trường vượt qua kiểm tra hợp lệ}}{\text{tổng số trường cần trích xuất}}
\end{equation}

Một trường được coi là hợp lệ khi giá trị sau chuẩn hóa thỏa mãn các
luật nghiệp vụ tương ứng, chẳng hạn số CCCD phải đủ chữ số, ngày tháng
phải đúng định dạng, giới tính chỉ nhận giá trị cho phép.

\textbf{Điểm tin cậy tổng hợp}

Hai thành phần được kết hợp theo trọng số cố định:
\vspace{-0.5em}
\begin{equation}
    \text{final\_confidence} = 0{,}70 \times \text{ocr\_confidence}
                             + 0{,}30 \times \text{parse\_confidence}
\end{equation}



Trọng số 70/30 ưu tiên tín hiệu từ mô hình nhận dạng nhưng vẫn phạt rõ
ràng khi kết quả OCR tuy có điểm cao nhưng không có nghĩa về nghiệp vụ.
Điểm tổng hợp này được lưu cùng mỗi bản ghi vào tệp CSV đầu ra và nạp
vào tầng Bronze. Ở bước kiểm tra chất lượng dữ liệu, các lô ảnh có điểm
tin cậy trung bình thấp hơn ngưỡng cho phép sẽ được ghi nhận cảnh báo
để người vận hành xem xét và xử lý thủ công.


\section{Trực quan hóa và Kết quả triển khai}

\subsection{Giao diện điều phối Apache Airflow}

Hệ thống sử dụng Apache Airflow để lập lịch và giám sát toàn bộ luồng
xử lý dữ liệu. Giao diện web cho phép theo dõi trạng thái từng DAG,
lịch sử chạy và log chi tiết của từng task.

% -------------------------------------------------------
% CHỤP 1: Trang DAGs — danh sách 3 DAG
%   banking_structured_dag | data_quality_dag | ocr_unstructured_dag
%   Cho thấy lịch chạy, trạng thái lần chạy gần nhất (success/failed)
% -------------------------------------------------------
\begin{figure}[H]
    \centering
    \includegraphics[width=\textwidth]{img4/giao diện airflow.png}
    \caption{Danh sách DAG trên giao diện Apache Airflow}
    \label{fig:airflow-dag-list}
\end{figure}

% -------------------------------------------------------
% CHỤP 2: Graph view của banking_structured_dag
%   Cho thấy toàn bộ chuỗi task: precheck → extract → MNS → dbt silver → dbt gold → notify
%   Tô màu xanh các task đã chạy thành công
% -------------------------------------------------------
\begin{figure}[H]
    \centering
    \includegraphics[width=\textwidth]{img4/banking_structured_dag.png}
    \caption{Graph view của DAG banking\_structured\_dag}
    \label{fig:airflow-dag-graph}
\end{figure}

% -------------------------------------------------------
% CHỤP 3 (tuỳ chọn): Log của một task bất kỳ (ví dụ extract_bronze_transactions)
%   Cho thấy số dòng được nạp, thông tin audit
% -------------------------------------------------------
\begin{figure}[H]
    \centering
    \includegraphics[width=\textwidth]{img4/log.png}
    \caption{Log chi tiết task trích xuất dữ liệu giao dịch}
    \label{fig:airflow-task-log}
\end{figure}


\subsection{Giao diện web demo xử lý tài liệu}

Hệ thống cung cấp giao diện web trực quan cho phép tải ảnh tài liệu lên
và nhận kết quả trích xuất theo thời gian thực. Ứng dụng hỗ trợ hai loại
tài liệu (CCCD và Sổ tiết kiệm) với 2 engine OCR tùy chọn: PaddleOCR,và Groq. ĐỐi với engine Groq, sau khi trích xuất, người dùng có thể đặt câu hỏi về
dữ liệu thông qua chatbot tích hợp.

% -------------------------------------------------------
% CHỤP 4: Giao diện upload — trang chính
%   Cho thấy: chọn loại tài liệu (CCCD / Sổ tiết kiệm),
%   vùng drag & drop, nút chọn engine OCR
% -------------------------------------------------------
\begin{figure}[H]
    \centering
    \includegraphics[width=0.85\textwidth]{img4/dd web.png}
    \caption{Giao diện tải ảnh tài liệu lên hệ thống}
    \label{fig:demo-upload}
\end{figure}

% -------------------------------------------------------
% CHỤP 5: Kết quả trích xuất CCCD
%   Bên trái: ảnh thẻ có khung ROI annotation
%   Bên phải: bảng các trường (Họ tên, Số CCCD, Ngày sinh...)
%             kèm badge điểm tin cậy (final_confidence)
% -------------------------------------------------------
\begin{figure}[H]
    \centering
    \includegraphics[width=0.8\textwidth]{img4/kết quả demo CCCD.png}
    \caption{Kết quả trích xuất thông tin CCCD}
    \label{fig:demo-result-idcard}
\end{figure}

% -------------------------------------------------------
% CHỤP 6: Kết quả trích xuất Sổ tiết kiệm
%   Tương tự CHỤP 5 nhưng với doc_type = savings_book
% -------------------------------------------------------
\begin{figure}[H]
    \centering
    \includegraphics[width=0.8\textwidth]{img4/kết quả demo STK.png}
    \caption{Kết quả trích xuất thông tin Sổ tiết kiệm}
    \label{fig:demo-result-savings}
\end{figure}

% -------------------------------------------------------
% CHỤP 7: Chatbot hỏi đáp
%   Giao diện chat phía dưới trang kết quả,
%   nhập câu hỏi về dữ liệu vừa trích xuất, chatbot trả lời bằng tiếng Việt
% -------------------------------------------------------
\begin{figure}[H]
    \centering
    \includegraphics[width=\textwidth]{img4/demo chatbot.png}
    \caption{Chatbot hỏi đáp về dữ liệu tài liệu đã trích xuất}
    \label{fig:demo-chatbot}
\end{figure}


\subsection{Hệ thống báo cáo BI}

Dữ liệu sau khi được xử lý qua toàn bộ pipeline được trực quan hóa
trên Power BI, kết nối trực tiếp với tầng Gold của kho dữ liệu.
Hệ thống báo cáo cung cấp các góc nhìn phân tích về hành vi giao dịch,
phân khúc khách hàng và rủi ro tín dụng.

% -------------------------------------------------------
% CHỤP 8: Trang tổng quan (Overview)
%   KPI cards: tổng giao dịch, tổng doanh thu, số khách hàng, tỷ lệ lỗi
%   Biểu đồ xu hướng giao dịch theo thời gian
% -------------------------------------------------------
\begin{figure}[H]
    \centering
    \includegraphics[width=\textwidth]{img4/bi_overview.png}
    \caption{Trang tổng quan hệ thống báo cáo Power BI}
    \label{fig:bi-overview}
\end{figure}

% -------------------------------------------------------
% CHỤP 9: Trang phân tích khách hàng
%   Phân khúc theo income_segment, credit_risk_tier
%   Phân bố debt_to_income_ratio
% -------------------------------------------------------
\begin{figure}[H]
    \centering
    \includegraphics[width=\textwidth]{img4/bi_customer.png}
    \caption{Phân tích phân khúc khách hàng}
    \label{fig:bi-customer}
\end{figure}

% -------------------------------------------------------
% CHỤP 10: Trang phân tích giao dịch
%   Top merchant, phân bố theo MCC, tỷ lệ chip vs non-chip
%   Bản đồ hoặc biểu đồ theo địa điểm nếu có
% -------------------------------------------------------
\begin{figure}[H]
    \centering
    \includegraphics[width=\textwidth]{img4/bi_transaction.png}
    \caption{Phân tích hành vi giao dịch}
    \label{fig:bi-transaction}
\end{figure}





\pagestyle{plain} % Hết chương 4: tắt header cho phần Kết luận và các mục sau (chỉ còn số trang)

%-----------------------------------------------------
% KẾT LUẬN
%-----------------------------------------------------
\chapter*{Kết luận}
\addcontentsline{toc}{chapter}{\bfseries Kết luận}

Xuất phát từ mục tiêu xây dựng một nền tảng dữ liệu thống nhất cho lĩnh vực ngân hàng, có khả năng tích hợp đồng thời dữ liệu có cấu trúc và dữ liệu phi cấu trúc, đồ án đã hoàn thành đầy đủ các mục tiêu đặt ra ban đầu. Cụ thể:

\begin{itemize}
    \item \textbf{Tích hợp trọn vẹn hai loại dữ liệu.} Đồ án đã xây dựng thành công hai luồng dữ liệu độc lập: một luồng ETL toàn diện xử lý dữ liệu giao dịch có cấu trúc (khách hàng, thẻ, giao dịch) và một luồng tinh gọn trích xuất dữ liệu phi cấu trúc (ảnh Căn cước công dân, Sổ tiết kiệm) bằng công nghệ nhận dạng ký tự quang học (OCR). Hai nguồn dữ liệu vốn rời rạc đã được đưa về cùng một nền tảng quản lý.

    \item \textbf{Hiện thực hóa kiến trúc đa tầng Medallion.} Hệ thống được tổ chức theo ba tầng Bronze -- Silver -- Gold rõ ràng. Tầng Bronze lưu dữ liệu thô; tầng Silver áp dụng mô hình Data Vault 2.0 (Hub -- Link -- Satellite) bảo đảm tính toàn vẹn và khả năng truy vết lịch sử; tầng Gold chuyển hóa sang mô hình đa chiều Star Schema nhằm tối ưu hóa cho phân tích và truy vấn.

    \item \textbf{Giải quyết bài toán nạp dữ liệu gia tăng.} Đồ án đã thiết kế và triển khai chiến lược nạp gia tăng dựa trên cơ chế MNS để bóc tách các thay đổi (thêm mới/cập nhật) theo từng ngày, thay vì nạp lại toàn bộ. Nhờ đó hệ thống xử lý dữ liệu một cách hiệu quả, có tính idempotent và quản lý được lịch sử biến đổi theo mô hình SCD.

    \item \textbf{Tự động hóa và điều phối luồng công việc.} Toàn bộ pipeline có cấu trúc được điều phối tự động bằng Apache Airflow theo lịch định kỳ, kết hợp công cụ dbt để thực hiện các phép biến đổi SQL một cách có kiểm soát và dễ bảo trì.

    \item \textbf{Bảo đảm chất lượng và khả năng giám sát dữ liệu.} Hệ thống tích hợp các bài kiểm tra chất lượng dữ liệu (Data Quality) như kiểm soát giá trị rỗng, trùng lặp khóa, độ trễ cập nhật, cùng hệ thống bảng Audit ghi nhận lịch sử mỗi lần chạy và đánh giá độ tin cậy của kết quả OCR.

    \item \textbf{Trực quan hóa phục vụ ra quyết định.} Dữ liệu đã được làm sạch tại tầng Gold được kết nối với Power BI để xây dựng các báo cáo quản trị (Dashboard) trực quan, hỗ trợ khai thác thông tin và ra quyết định kinh doanh.
\end{itemize}

Như vậy, đồ án không chỉ dừng lại ở việc đề xuất giải pháp về mặt lý thuyết mà đã triển khai thành một hệ thống hoàn chỉnh, chạy thông suốt từ khâu thu thập dữ liệu nguồn đến khâu trực quan hóa, đáp ứng đầy đủ các yêu cầu chức năng và phi chức năng đã đặt ra trong giai đoạn khảo sát.
\newpage
\section*{Kết quả của đồ án}

Sau quá trình nghiên cứu và triển khai, đồ án đã tạo ra một hệ thống nền tảng dữ liệu hoàn chỉnh với các sản phẩm cụ thể sau:

\begin{itemize}
    \item \textbf{Kho dữ liệu đa tầng hoàn chỉnh.} Hệ thống đã dựng được kho dữ liệu trên SQL Server với bốn lược đồ \texttt{bronze}, \texttt{silver}, \texttt{gold} và \texttt{audit}. Tầng Silver theo mô hình Data Vault gồm 5 Hub, 2 Link và 4 Satellite; tầng Gold theo mô hình Star Schema 

    \item \textbf{Luồng trích xuất dữ liệu phi cấu trúc.} Đồ án đã xây dựng module OCR trích xuất thông tin từ hai loại chứng từ là Căn cước công dân và Sổ tiết kiệm, kèm cơ chế tiền xử lý ảnh và bộ điểm tin cậy ba mức để đánh giá chất lượng kết quả.

    \item \textbf{Hệ thống điều phối tự động.} Ba luồng công việc được đóng gói thành các DAG trên Apache Airflow: pipeline có cấu trúc (Bronze $\rightarrow$ Silver $\rightarrow$ Gold), pipeline kiểm tra chất lượng dữ liệu và pipeline OCR. Mỗi tác vụ đều được bao bọc bởi cơ chế ghi nhật ký Audit, cho phép giám sát số dòng xử lý, thời gian chạy và trạng thái thành công/thất bại của từng bước.

    \item \textbf{Cơ chế bảo đảm chất lượng dữ liệu.} Hệ thống tích hợp các bài kiểm tra như phát hiện trùng lặp khóa chính, kiểm soát giá trị rỗng ở khóa ngoại, kiểm tra độ trễ cập nhật và ngưỡng tin cậy OCR, giúp phát hiện sớm các bất thường trong dữ liệu.

    \item \textbf{Ứng dụng web minh họa.} Đồ án phát triển một ứng dụng web demo cho phép tải lên ảnh chứng từ và trích xuất thông tin trực tuyến với hai lựa chọn engine (PaddleOCR, Groq), kèm trợ lý hỏi đáp về dữ liệu vừa trích xuất.

    \item \textbf{Hệ thống báo cáo trực quan.} Dữ liệu tại tầng Gold được kết nối với Power BI để xây dựng các bảng điều khiển (Dashboard) phân tích giao dịch, hỗ trợ trực quan hóa và khai thác thông tin phục vụ ra quyết định.

    \item \textbf{Giải pháp triển khai linh hoạt.} Toàn bộ hệ thống được đóng gói bằng Docker Compose với hai môi trường Python được cô lập, giải quyết triệt để xung đột tài nguyên giữa engine OCR và bộ công cụ dữ liệu, đồng thời giúp việc cài đặt và vận hành trở nên thuận tiện.
\end{itemize}

\section*{Kỹ năng đạt được}

Quá trình thực hiện đồ án đã giúp em tích lũy được nhiều kiến thức và kỹ năng chuyên môn có giá trị:

\begin{itemize}
    \item \textbf{Kỹ năng thiết kế kho dữ liệu.} Hiểu và vận dụng được các mô hình kiến trúc dữ liệu hiện đại như Medallion, Data Vault 2.0 và Star Schema, cũng như chiến lược nạp dữ liệu gia tăng và quản lý lịch sử thay đổi (SCD).

    \item \textbf{Kỹ năng kỹ thuật dữ liệu (Data Engineering).} Thành thạo việc xây dựng pipeline ETL/ELT, sử dụng công cụ dbt để biến đổi dữ liệu theo hướng khai báo và Apache Airflow để điều phối luồng công việc tự động.

    \item \textbf{Kỹ năng xử lý dữ liệu phi cấu trúc.} Nắm được quy trình nhận dạng ký tự quang học (OCR), các kỹ thuật tiền xử lý ảnh và phương pháp đánh giá độ tin cậy của kết quả trích xuất.

    \item \textbf{Kỹ năng lập trình và làm chủ công nghệ.} Củng cố năng lực lập trình Python, làm việc với SQL Server, đóng gói ứng dụng bằng Docker, quản lý mã nguồn với Git

    \item \textbf{Kỹ năng phân tích và trực quan hóa.} Biết cách khảo sát, phân tích yêu cầu nghiệp vụ và sử dụng Power BI để xây dựng báo cáo phục vụ ra quyết định.

    \item \textbf{Kỹ năng làm việc và nghiên cứu độc lập.} Rèn luyện tư duy giải quyết vấn đề một cách hệ thống, khả năng tự tìm hiểu công nghệ mới, quản lý tiến độ và trình bày kết quả nghiên cứu một cách khoa học.
\end{itemize}

\section*{Hướng phát triển của đồ án trong tương lai}

Mặc dù đã hoàn thành các mục tiêu đề ra, hệ thống vẫn còn nhiều tiềm năng để mở rộng và hoàn thiện trong tương lai:

\begin{itemize}
    \item \textbf{Nâng cao hiệu năng xử lý OCR.} Triển khai xử lý song song để tăng tốc độ trích xuất khi khối lượng tài liệu lớn, đồng thời bổ sung cơ chế thử lại cho các tài liệu xử lý thất bại.

    \item \textbf{Cải thiện độ chính xác trích xuất.} Tinh chỉnh sâu hơn các bước tiền xử lý ảnh và bổ sung khả năng nhận diện nhiều mẫu (template) chứng từ khác nhau, giúp hệ thống thích ứng với đa dạng định dạng tài liệu trong thực tế.

    \item \textbf{Mở rộng nguồn và loại dữ liệu.} Tích hợp thêm các loại chứng từ ngân hàng khác và kết nối với nhiều hệ thống nguồn, hướng tới mô hình thu thập dữ liệu thay đổi (CDC) theo thời gian thực thay vì nạp theo lô hằng ngày.

    \item \textbf{Hoàn thiện giao diện người dùng.} Phát triển giao diện ứng dụng cho phép người dùng nghiệp vụ dễ dàng theo dõi, kiểm tra và can thiệp, hiệu chỉnh kết quả trích xuất trước khi đưa vào kho dữ liệu.

    \item \textbf{Tăng cường khả năng phân tích nâng cao.} Ứng dụng các mô hình học máy (Machine Learning) trên nền dữ liệu đã được làm sạch để phục vụ các bài toán như phát hiện gian lận giao dịch, phân khúc khách hàng hay dự báo xu hướng.

    \item \textbf{Củng cố vận hành và bảo mật.} Bổ sung cảnh báo tự động khi pipeline gặp sự cố, tăng cường cơ chế phân quyền và bảo vệ dữ liệu nhạy cảm phù hợp với các quy định về an toàn thông tin trong lĩnh vực ngân hàng.
\end{itemize}

\renewcommand{\bibname}{Tài liệu tham khảo}
\begin{thebibliography}{99}

% --- Mô hình hóa kho dữ liệu ---
\bibitem{linstedt2015}
D. Linstedt and M. Olschimke, \textit{Building a Scalable Data Warehouse with Data Vault 2.0}. Morgan Kaufmann, 2015.

\bibitem{kimball2013}
R. Kimball and M. Ross, \textit{The Data Warehouse Toolkit: The Definitive Guide to Dimensional Modeling}, 3rd ed. Wiley, 2013.

\bibitem{inmon2005}
W. H. Inmon, \textit{Building the Data Warehouse}, 4th ed. Wiley, 2005.

\bibitem{databricks_medallion}
Databricks, ``Medallion Architecture,'' Databricks Documentation. [Trực tuyến]. Địa chỉ: \url{https://docs.databricks.com/aws/en/lakehouse/medallion}. Truy cập: tháng 2/2026.

% --- Công cụ điều phối và biến đổi dữ liệu ---
\bibitem{airflow_docs}
Apache Software Foundation, ``Apache Airflow Documentation.'' [Trực tuyến]. Địa chỉ: \url{https://airflow.apache.org/docs/}. Truy cập: tháng 6/2026.

\bibitem{dbt_docs}
dbt Labs, ``dbt Documentation.'' [Trực tuyến]. Địa chỉ: \url{https://docs.getdbt.com/}. Truy cập: tháng 4/2026.

\bibitem{mssql_docs}
Microsoft, ``SQL Server Documentation.'' [Trực tuyến]. Địa chỉ: \url{https://learn.microsoft.com/en-us/sql/sql-server/}. Truy cập: tháng 3/2026.

% --- Nhận dạng ký tự và xử lý ảnh ---
\bibitem{ppocr2020}
Y. Du \textit{et al.}, ``PP-OCR: A Practical Ultra Lightweight OCR System,'' \textit{arXiv preprint arXiv:2009.09941}, 2020.

\bibitem{paddleocr_repo}
PaddlePaddle, ``PaddleOCR,'' GitHub repository. [Trực tuyến]. Địa chỉ: \url{https://github.com/PaddlePaddle/PaddleOCR}. Truy cập: tháng 3/2026.

\bibitem{opencv}
G. Bradski, ``The OpenCV Library,'' \textit{Dr. Dobb's Journal of Software Tools}, 2000. [Trực tuyến]. Địa chỉ: \url{https://docs.opencv.org/}.

% --- Tiêu chuẩn ngành ---
\bibitem{iso18245}
International Organization for Standardization, \textit{ISO 18245:2023 -- Retail financial services -- Merchant category codes}. ISO, 2023.

\bibitem{pcidss}
PCI Security Standards Council, \textit{Payment Card Industry Data Security Standard (PCI DSS), v4.0}, 2022. [Trực tuyến]. Địa chỉ: \url{https://www.pcisecuritystandards.org/}.

\end{thebibliography}







\end{document}